\PassOptionsToPackage{english}{babel}
\documentclass{notes}
\usepackage[utf8]{inputenc}
\usepackage{mathtools}
\usepackage{mathrsfs}
\usepackage{enumitem}
\usepackage{adjustbox}
\usepackage{pdflscape}
\usepackage{longtable}
\usepackage{float}
\usepackage{needspace}
\usepackage{microtype}
\usepackage{booktabs}
\usepackage{array}
\usepackage{tikz}
\usetikzlibrary{positioning,arrows.meta,shapes.geometric,fit,calc,matrix}
\usepackage{mathmongo-macros}
\hypersetup{colorlinks=true,linkcolor=blue!55!black,citecolor=teal!55!black,urlcolor=blue!55!black}
\hypersetup{pdftitle={Study Notes: Spectral Theory in Banach and C* Algebras},pdfauthor={Enrique Díaz Ocampo},pdflang={en-US}}

% The shared style defines its theorem boxes in Spanish; the English
% edition changes only their displayed names.
\renewcommand{\cthdefinicionname}{Definition}
\renewcommand{\cthteoremaname}{Theorem}
\renewcommand{\cthcorolarioname}{Corollary}
\renewcommand{\cthlemaname}{Lemma}
\renewcommand{\cthproposicionname}{Proposition}
\renewcommand{\cthejemploname}{Example}
\renewcommand{\cthnotaname}{Note}
\renewcommand{\cthobservacionname}{Observation}
\renewtcolorbox{context}{notesbase,title=Context,colback=gray!6,colframe=gray!60}
\renewtcolorbox{reflection}{notesbase,title=Reflection,colback=yellow!10,colframe=yellow!65}

\newcommand{\A}{\mathcal A}
\newcommand{\B}{\mathcal B}
\newcommand{\K}{\mathcal K}
\newcommand{\Calkin}{\mathcal Q}
\newcommand{\GL}{\operatorname{GL}}
\newcommand{\spn}{\operatorname{span}}
\newcommand{\alg}{\operatorname{alg}}
\newcommand{\Id}{I}
\newcommand{\spec}{\sigma}
\newcommand{\rad}{r}
\newcommand{\bbD}{\mathbb D}
\newcommand{\bbT}{\mathbb T}
\newcommand{\SO}{\operatorname{SO}}
\newcommand{\PSO}{\operatorname{PSO}}
\newcommand{\ran}{\operatorname{Ran}}
\newcommand{\Ker}{\operatorname{Ker}}
\newcommand{\Ind}{\operatorname{Ind}}
\newcommand{\clos}{\overline}

\newcommand{\origen}[1]{%
  \par\smallskip{\footnotesize\color{DIgray}\textbf{Correspondence with the source notes:} #1.}\par\smallskip
}

\newtcolorbox{dependencias}[1][]{
  notesbase,
  title={Logical dependencies},
  colback=DIlightblue,
  colframe=DIblue,
  #1
}
\newtcolorbox{investigacion}[1][]{
  notesbase,
  title={Example from the Haseman/SIO research line},
  colback=DIlightgreen,
  colframe=DIgreen,
  #1
}
\newtcolorbox{puente}[1][]{
  notesbase,
  title={Bridge to Section 2},
  colback=DIlightpurple,
  colframe=DIpurple,
  #1
}
\newtcolorbox{criterio}[1][]{
  notesbase,
  title={Key point to remember},
  colback=DIlightorange,
  colframe=DIorange,
  #1
}

\tikzset{
  flowbox/.style={draw=DIblue, rounded corners=3pt, very thick, fill=DIlightblue,
    align=center, text width=3.0cm, minimum height=1.15cm, inner sep=6pt},
  flowbox2/.style={draw=DIteal, rounded corners=3pt, very thick, fill=BootTealLight,
    align=center, text width=3.0cm, minimum height=1.15cm, inner sep=6pt},
  flowbox3/.style={draw=DIpurple, rounded corners=3pt, very thick, fill=DIlightpurple,
    align=center, text width=3.0cm, minimum height=1.15cm, inner sep=6pt},
  flowend/.style={draw=DIgreen, rounded corners=3pt, very thick, fill=DIlightgreen,
    align=center, text width=3.0cm, minimum height=1.15cm, inner sep=6pt},
  taskbox/.style={draw=DIorange, rounded corners=3pt, thick, fill=DIlightorange,
    align=center, text width=3.1cm, minimum height=0.95cm, inner sep=4pt},
  theorybox/.style={draw=DIblue, rounded corners=3pt, thick, fill=DIlightblue,
    align=center, text width=3.15cm, minimum height=0.95cm, inner sep=4pt},
  researchbox/.style={draw=DIgreen, rounded corners=3pt, thick, fill=DIlightgreen,
    align=center, text width=3.2cm, minimum height=0.95cm, inner sep=4pt},
  arr/.style={-{Latex[length=2.6mm]},thick,draw=DIgray},
  mergearr/.style={-{Latex[length=2.6mm]},thick,draw=DIpurple}
}

\newcommand{\proofdiagram}[4]{%
\begin{center}
\begin{tikzpicture}[node distance=7mm]
  \node[flowbox] (h) {\textbf{Assumptions / data}\\#1};
  \node[flowbox2,right=of h] (r) {\textbf{Earlier results}\\#2};
  \node[flowbox3,right=of r] (m) {\textbf{Main mechanism}\\#3};
  \node[flowend,right=of m] (c) {\textbf{Conclusion}\\#4};
  \draw[arr] (h) -- (r);
  \draw[arr] (r) -- (m);
  \draw[arr] (m) -- (c);
\end{tikzpicture}
\end{center}
}

\newcommand{\proofdiagrambranch}[5]{%
\begin{center}
\begin{tikzpicture}[node distance=7mm]
  \node[flowbox] (h1) {\textbf{Datum 1}\\#1};
  \node[flowbox,below=of h1] (h2) {\textbf{Datum 2}\\#2};
  \node[flowbox2,anchor=west] (r) at ($(h1.east)!0.5!(h2.east)+(9mm,0)$)
    {\textbf{Tool}\\#3};
  \node[flowbox3,right=of r] (m) {\textbf{Decisive step}\\#4};
  \node[flowend,right=of m] (c) {\textbf{Conclusion}\\#5};
  \draw[arr] (h1.east) -- (r.west);
  \draw[arr] (h2.east) -- (r.west);
  \draw[arr] (r) -- (m);
  \draw[arr] (m) -- (c);
\end{tikzpicture}
\end{center}
}

\newcommand{\examplediagram}[5]{%
\begin{center}
\begin{tikzpicture}[node distance=6mm and 7mm]
  \node[theorybox] (g) {\textbf{General concept / result}\\#1};
  \node[taskbox,right=of g] (k) {\textbf{Specific instance}\\#2};
  \node[researchbox,right=of k] (t) {\textbf{Theoretical expression}\\#3};
  \node[theorybox,below=of t] (s) {\textbf{Symbolic expression}\\#4};
  \node[taskbox,left=of s] (i) {\textbf{Interpretation}\\#5};
  \draw[arr] (g) -- (k);
  \draw[arr] (k) -- (t);
  \draw[arr] (t) -- (s);
  \draw[arr] (s) -- (i);
\end{tikzpicture}
\end{center}
}

\newcommand{\tasksource}[1]{\par\smallskip\textbf{Results used:} #1\par\smallskip}
\newcommand{\researchcite}[1]{\textcolor{DIgreen}{\textbf{Research connection:} #1}}

\title{Study Notes: Spectral Theory in Banach and $C^*$-Algebras}
\author{Enrique Díaz Ocampo}
\date{September 2026}

\begin{document}
\selectlanguage{english}
\maketitle
\tableofcontents

\begin{context}
These notes bring together \emph{Part 1: Spectral Theory in Banach Algebras}
\cite{CastillejosParte1}, \emph{Section 2: $C^*$-Algebras}
\cite{CastillejosCap2}, \emph{Assignment 1}, and \emph{Assignment 2}. Each result is studied
from its assumptions through its proof and a diagram that summarizes the argument.
Elementary examples come before connections to the literature: quotients by compact
operators, slowly oscillating functions, and Cauchy components of the Haseman problem.
Each connection identifies its underlying space and the additional assumptions it
requires, so the reader can trace an application back to the course results.
\end{context}

\chapter{Logical Map of the Two Topics}

\section{The central idea}

Let $\A$ be a unital complex Banach algebra and let $a\in\A$. Part 1 begins with a question:
when does $(\lambda1-a)^{-1}$ exist? The Neumann series answers this near an invertible
element; the resolvent organizes the answer for all scalars; and the spectrum records the values
at which inversion fails. Section 2 adds the $C^*$ identity and obtains consequences that do not
hold in an arbitrary Banach algebra.

\section{Map of Part 1}

\begin{center}
\begin{tikzpicture}[node distance=9mm and 8mm]
  \node[theorybox,text width=3.25cm] (ban) {Unital Banach algebra\\completeness and product};
  \node[theorybox,text width=3.25cm,right=of ban] (neu) {Neumann series\\$\|x\|<1\Rightarrow(1-x)^{-1}$};
  \node[theorybox,text width=3.25cm,right=of neu] (open) {$\GL(\A)$ is open\\continuous inversion};
  \node[theorybox,text width=3.25cm,below=of neu] (bound) {Bound for $|\lambda|>\|a\|$\\compact spectrum};
  \node[theorybox,text width=3.25cm,right=of bound] (res) {Resolvent\\identity and holomorphy};
  \node[theorybox,text width=3.25cm,below=of bound] (non) {Nonempty spectrum\\Liouville + Hahn--Banach};
  \draw[arr] (ban)--(neu);
  \draw[arr] (neu)--(open);
  \draw[arr] (open)--(res);
  \draw[arr] (neu)--(bound);
  \draw[arr] (open.south west)--(bound.north east);
  \draw[arr] (res.south) |- (non.east);
  \draw[arr] (bound)--(non);
\end{tikzpicture}
\end{center}

The arrow leading to nonemptiness requires two separate ingredients: holomorphy of the
resolvent and its decay at infinity, obtained from the Neumann series. Compactness uses
the openness of $\GL(\A)$ to establish closedness and the same series for boundedness.

\begin{center}
\begin{tikzpicture}[node distance=5mm and 8mm]
  \node[theorybox,text width=3.25cm] (non) {Nonempty spectrum};
  \node[theorybox,text width=3.25cm,right=of non] (gm) {Gelfand--Mazur\\division algebras};
  \node[theorybox,text width=3.25cm,right=of gm] (max) {Maximal ideals\\$\A/J\cong\mathbb C$ if $\A$ is commutative};
  \node[theorybox,text width=3.25cm,below=of max] (char) {Characters\\$\Omega(\A)$};
  \node[theorybox,text width=3.25cm,left=of char] (gel) {Gelfand transform\\$a\mapsto\widehat a$};
  \node[theorybox,text width=3.25cm,left=of gel] (rad) {Spectral radius\\resolvent + powers};
  \draw[arr] (non)--(gm);
  \draw[arr] (gm)--(max);
  \draw[arr] (max)--(char);
  \draw[arr] (char)--(gel);
  \draw[arr] (non)--(rad);
\end{tikzpicture}
\end{center}

The spectral radius lies on a separate branch: its formula uses holomorphy of the resolvent,
the spectral relation for powers, and the uniform boundedness principle; it does not follow
from the Gelfand transform. Identifying characters with maximal ideals requires
commutativity, whereas the initial spectral theory does not.

\section{Map of Section 2}

\begin{center}
\begin{tikzpicture}[node distance=5mm and 8mm]
  \node[theorybox,text width=3.25cm] (id) {$C^*$ identity\\$\|a^*a\|=\|a\|^2$};
  \node[theorybox,text width=3.25cm,right=of id] (nr) {For normal $a$\\$\|a\|=r(a)$};
  \node[theorybox,text width=3.25cm,right=of nr] (uq) {Uniqueness of the $C^*$ norm\\via the radius of $a^*a$};
  \draw[arr] (id)--(nr);
  \draw[arr] (nr)--(uq);
\end{tikzpicture}
\end{center}

\begin{center}
\begin{tikzpicture}[node distance=5mm and 8mm]
  \node[theorybox,text width=3.25cm] (a) {Normal $a$ in a unital $C^*$-algebra};
  \node[theorybox,text width=3.25cm,right=of a] (comm) {Commutative $C^*(a,1)$\\because $aa^*=a^*a$};
  \node[theorybox,text width=3.25cm,right=of comm] (gn) {Commutative\\Gelfand--Naimark};
  \node[theorybox,text width=3.25cm,below=of gn] (fc) {Continuous functional calculus\\$f\mapsto f(a)$};
  \node[theorybox,text width=3.25cm,left=of fc] (class) {Spectral classification\\and positive roots};
  \draw[arr] (a)--(comm);
  \draw[arr] (comm)--(gn);
  \draw[arr] (gn)--(fc);
  \draw[arr] (fc)--(class);
\end{tikzpicture}
\end{center}

The equality $\|a\|=r(a)$ requires normality and the $C^*$ identity.
Commutativity of $C^*(a,1)$ follows directly from $aa^*=a^*a$.
The Gelfand transform from Part 1 and the norm equality yield the
commutative isometric representation. Spectral permanence identifies the spectrum
of $a$ in $C^*(a,1)$ with its spectrum in the ambient algebra. The positive root and
polar decomposition then use functional calculus under the assumptions in their statements.

\section{Three reading paths}

\noindent\textbf{Invertibility and spectrum.} Neumann series $\Rightarrow$ stability of invertible
elements $\Rightarrow$ holomorphic resolvent; together with the estimate at infinity and
Liouville's theorem, this gives $\sigma_{\A}(a)\neq\varnothing$. For operators
$A\in\mathcal B(X)$, the same resolvent is $(\lambda I_X-A)^{-1}$.

\noindent\textbf{Commutativity and functions.} For commutative unital $\A$, maximal ideals
produce characters, which identify $\sigma_{\A}(a)$ with the range of $\widehat a$.
If $\A$ is also a $C^*$-algebra, the commutative representation yields functional calculus.

\noindent\textbf{Quotients and operators.} For a Banach space $X$, the closed ideal
$\mathcal K(X)\triangleleft\mathcal B(X)$ gives the Calkin quotient. Invertibility of
$A+\mathcal K(X)$ is equivalent to $A$ being Fredholm; this equivalence is a theorem of
operator theory applied after constructing the quotient, rather than an elementary theorem
about every Banach algebra. The research chapter discusses the Haseman operators.

\chapter{Spectral theory in Banach algebras}

In the unital results of this chapter, $1$ abbreviates the identity $1_\A$ of
the ambient algebra $\A$; we write $1_\A$ when the space to which an expression
belongs needs to be specified.

\section{Banach algebras, involution, and basic examples}

\begin{cthdefinicion}{(Banach algebra)}
A \emph{Banach algebra} is an associative algebra $\A$ over $\mathbb C$, equipped with a
norm such that $(\A,\|\cdot\|)$ is a Banach space and
\[
  \|ab\|\le \|a\|\,\|b\|,\qquad a,b\in\A.
\]
We say that $\A$ is \emph{unital} if there is an identity $1_\A$ with $\|1_\A\|=1$, and \emph{commutative} if
$ab=ba$ for every $a,b\in\A$.
\end{cthdefinicion}
\origen{Definition 1.1}

\begin{cthejemplo}{(Bounded functions as a Banach algebra)}
We seek a commutative model of the definition. Let $X$ be nonempty,
locally compact, and Hausdorff, and consider
\[
 C_b(X)=\{f:X\to\mathbb C: f\text{ is continuous and bounded}\},\qquad
 \|f\|_\infty=\sup_{x\in X}|f(x)|,
\]
with pointwise operations. For $f,g\in C_b(X)$,
$|f(x)g(x)|\le\|f\|_\infty\|g\|_\infty$ at each $x$, and therefore
$\|fg\|_\infty\le\|f\|_\infty\|g\|_\infty$.
The uniform limit of continuous bounded functions remains in $C_b(X)$,
which gives completeness; the constant function $1$ is the identity.
Thus $C_b(X)$ meets the definition of a Banach algebra and is commutative.
If $X$ is compact, every continuous function is bounded and $C_b(X)=C(X)$.
\end{cthejemplo}

\begin{cthejemplo}{(Composition of bounded operators)}
To contrast with pointwise multiplication, let $X$ be a nonzero Banach space.
We study $\mathcal B(X)$ (also denoted $\mathcal L(X)$), the algebra
of bounded linear operators $X\to X$, with composition and the operator
norm. The inequality $\|STx\|\le\|S\|\|T\|\|x\|$ implies
$\|ST\|\le\|S\|\|T\|$; moreover, the identity $I_X$ is a unit and
the norm limit of bounded operators remains in $\mathcal B(X)$.
The product can change when the order is reversed: in
$\mathcal B(\mathbb C^2)=M_2(\mathbb C)$, if $E_{ij}$ denotes the matrix
with a $1$ in position $(i,j)$ and zeros elsewhere,
$E_{12}E_{21}=E_{11}$, whereas $E_{21}E_{12}=E_{22}$.
This example shows why subsequent proofs cannot assume
commutativity of the entire algebra.
\end{cthejemplo}

\begin{investigacion}
\textbf{Connection with the literature: a direct example.}
Under Karlovich's hypotheses for a starlike curve $\Gamma$ and $1<p<\infty$, the boundary
problem leads to an element of $\mathcal B(L^p(\Gamma))$:
\[
 T=V_\alpha P_+ + M_G P_-\in \mathcal B(L^p(\Gamma)),
\]
where $V_\alpha f=f\circ\alpha$, $M_Gf=Gf$, and $P_\pm=(I\pm S_\Gamma)/2$.
The formula shows why sums and compositions are studied in a single algebra. In this
geometry one does not assume $S_\Gamma^2=I$ or that $P_\pm$ are algebraic projections;
the reduction by Sokhotski--Plemelj is explained at the end of the notes
\cite[pp.~463--464, eqs.~(1.1)--(1.2)]{Karlovich2017}.
\textbf{What to recognize:} $\mathcal B(L^p(\Gamma))$ realizes the abstract algebra of
bounded operators; the formula for $T$ combines elements of that algebra.
\end{investigacion}

\begin{cthdefinicion}{(Involution)}
An \emph{involution} on a complex algebra $\A$ is a map $a\mapsto a^*$ such that
\[
(\lambda a+b)^*=\overline\lambda a^*+b^*,\qquad
(a^*)^*=a,\qquad
(ab)^*=b^*a^*.
\]
These identities are required for all $a,b\in\A$ and $\lambda\in\mathbb C$.
If, in addition, $\A$ is a Banach algebra and $\|a^*\|=\|a\|$ for
every $a\in\A$, we call $\A$ an \emph{involutive Banach algebra}
in the sense of these notes.
\end{cthdefinicion}
\origen{Definition 1.4}

\begin{cthejemplo}{(Pointwise conjugation as an involution)}
In the algebra $C_b(X)$ from the preceding example, we want to identify a
concrete involution. For $f\in C_b(X)$ define
$f^*(x)=\overline{f(x)}$. If $f,g\in C_b(X)$, then
$(fg)^*(x)=\overline{f(x)g(x)}=g^*(x)f^*(x)$;
also $(f^*)^*=f$ and $\|f^*\|_\infty=\|f\|_\infty$.
Pointwise conjugation thus satisfies the definition of an isometric
involution; the equality $g^*f^*=f^*g^*$ reflects the commutativity of this model.
\end{cthejemplo}

\section{Invertibility and the Neumann series}

\begin{cthdefinicion}{(Group of invertible elements)}
Let $\A$ be a unital Banach algebra, with identity $1_\A$. Its group of invertible elements is
\[
 \GL(\A)=\{a\in\A: \exists b\in\A,\ ab=ba=1_\A\}.
\]
For $a\in\GL(\A)$, that $b$ is unique and is written $a^{-1}$.
\end{cthdefinicion}
\origen{Definition 1.7}

\begin{cthproposicion}{(Neumann series)}
Let $\A$ be a unital Banach algebra, with identity $1_\A$, and let $a\in\A$.
If $\|1_\A-a\|<1$, then $a\in\GL(\A)$ and
\[
 a^{-1}=\sum_{n=0}^{\infty}(1_\A-a)^n.
\]
\end{cthproposicion}
\origen{Proposition 1.8}

\begin{proof}[Proof]
Write $q=1_\A-a$. The hypothesis gives $\|q\|<1$ and, by submultiplicativity,
$\|q^n\|\le\|q\|^n$ for every $n\ge0$. Hence
\[
 \sum_{n=0}^{\infty}\|q^n\|
 \le \sum_{n=0}^{\infty}\|q\|^n<\infty.
\]
Completeness of $\A$ ensures that the partial sums
$s_N=\sum_{n=0}^{N}q^n$ converge in norm to an element $b\in\A$.
Since $a=1_\A-q$, multiplication on the left telescopes:
\[
 as_N=(1_\A-q)\sum_{n=0}^{N}q^n
 =\sum_{n=0}^{N}q^n-\sum_{n=1}^{N+1}q^n
 =1_\A-q^{N+1}.
\]
Multiplication on the right must also be checked:
\[
 s_Na=\sum_{n=0}^{N}q^n(1_\A-q)
 =\sum_{n=0}^{N}q^n-\sum_{n=1}^{N+1}q^n
 =1_\A-q^{N+1}.
\]
Here we have used only that $q$ commutes with its own powers, without requiring
$\A$ to be commutative. Finally,
$\|q^{N+1}\|\le\|q\|^{N+1}\to0$; continuity of
multiplication lets us pass to the limit in both identities to obtain
$ab=ba=1_\A$. Thus $a\in\GL(\A)$ and $b=a^{-1}$.
\end{proof}

\proofdiagram
{$\A$ unital Banach, $a\in\A$; $\|1_\A-a\|<1$}
{Completeness of $\A$; $\|(1-a)^n\|\le\|1-a\|^n$}
{The geometric series converges and the product telescopes}
{$a^{-1}=\sum_{n\ge0}(1_\A-a)^n$}

\begin{cthejemplo}{(Neumann series for a perturbed matrix)}
We want to invert a matrix close to the identity. In
$M_2(\mathbb C)$ with the operator norm, let $A=I_2-E$ with
$\|E\|<1$. The Neumann proposition applies because
$\|I_2-A\|=\|E\|<1$, and gives
\[
 A^{-1}=I_2+E+E^2+\cdots.
\]
For example, if
$E=\begin{pmatrix}0&t\\0&0\end{pmatrix}$ with $|t|<1$,
then $E^2=0$ and the series terminates at
$A^{-1}=I_2+E$; indeed, $(I_2-E)(I_2+E)=I_2$.
The norm condition ensures invertibility even when the series does not terminate.
\end{cthejemplo}

\begin{cthejemplo}{(Stable invertibility near an operator)}
Let $X$ be a nonzero Banach space and suppose an invertible operator
$A_0\in\mathcal B(X)$ is known. We want to invert $A_0+E$ when
$E\in\mathcal B(X)$ is small relative to $A_0$, in the sense that
$\|A_0^{-1}E\|<1$. The factorization
$A_0+E=A_0(I_X+A_0^{-1}E)$ allows us to apply Neumann to the second factor:
\[
 (A_0+E)^{-1}=(I_X+A_0^{-1}E)^{-1}A_0^{-1}
 =\sum_{n=0}^{\infty}(-A_0^{-1}E)^n A_0^{-1}.
\]
The formula computes an exact inverse and shows that the relevant scale of
the perturbation is $\|A_0^{-1}E\|$. The Fredholm question will be formulated
later through invertibility of a class in a quotient.
\end{cthejemplo}

\begin{cthnota}{(Where the hypotheses are used in the Neumann series)}
The inequality $\|1-a\|<1$ makes the series of norms summable; completeness of $\A$ produces
its sum inside $\A$; the identity lets us write $a=1-(1-a)$ and telescope.
Commutativity of the entire algebra is not used.
\end{cthnota}

\begin{cthproposicion}{(Local stability of invertibility)}
Let $\A$ be a unital Banach algebra. The set $\GL(\A)$ is open.
More precisely, if $a\in\GL(\A)$ and $b\in\A$ satisfy
\[
 \|b-a\|<\|a^{-1}\|^{-1},
\]
then $b\in\GL(\A)$.
\end{cthproposicion}
\origen{Remark 1.9}

\begin{proof}[Proof]
Put $c=a^{-1}(a-b)$. The hypothesis and submultiplicativity give
\[
 \|c\|\le \|a^{-1}\|\,\|a-b\|<1,
 \qquad b=a(1_\A-c).
\]
The preceding proposition gives a two-sided inverse
$(1_\A-c)^{-1}=\sum_{n\ge0}c^n$ inside $\A$. The product of
the two invertible factors is invertible, with
\[
 b^{-1}=(1_\A-c)^{-1}a^{-1}
 =\sum_{n=0}^{\infty}\bigl[a^{-1}(a-b)\bigr]^n a^{-1}.
\]
Therefore every ball $B(a,\|a^{-1}\|^{-1})$ is contained in
$\GL(\A)$, which is open.
\end{proof}

\proofdiagram
{$\A$ unital, $a\in\GL(\A)$ and $b\in\A$ close to $a$}
{Neumann series}
{Factor $b=a\,[1-a^{-1}(a-b)]$}
{$\|a^{-1}(a-b)\|<1\Rightarrow b$ invertible}

\begin{cthejemplo}{(Stability radius around $2I_n$)}
We want to make the invertibility ball of the preceding proposition concrete.
In $M_n(\mathbb C)$ with the induced operator norm, take
$a=2I_n$; then $a^{-1}=\frac12I_n$ and
$\|a^{-1}\|^{-1}=2$. If a matrix $b$ satisfies
$\|b-2I_n\|<2$, the local stability condition holds and
$b$ is invertible. Here the number $2$ measures how much the
matrix can vary before this criterion ceases to guarantee an inverse.
\end{cthejemplo}

\begin{cthcorolario}{(Invertibility outside the disk of radius $\|a\|$)}
Let $\A$ be a complex unital Banach algebra, with identity $1_\A$,
$a\in\A$, and $\lambda\in\mathbb C$. If $\|a\|<|\lambda|$, then
$\lambda1_\A-a\in\GL(\A)$.
\end{cthcorolario}
\origen{Corollary 1.10}

\begin{proof}[Proof]
The inequality $\|a\|<|\lambda|$ implies $\lambda\ne0$. For
$c=\lambda^{-1}a$ we have $\|c\|=\|a\|/|\lambda|<1$;
by the Neumann proposition, $1_\A-c$ has the two-sided inverse
$\sum_{n\ge0}c^n$. Since
$\lambda1_\A-a=\lambda(1_\A-c)$ and the scalar
$\lambda1_\A$ is invertible, so is $\lambda1_\A-a$.
More precisely, its two-sided inverse is
\[
 (\lambda1_\A-a)^{-1}
 =\lambda^{-1}\sum_{n=0}^{\infty}(\lambda^{-1}a)^n.
\]
The series converges in norm because $\|\lambda^{-1}a\|<1$.
\end{proof}

\proofdiagram
{$a\in\A$, $\lambda\in\mathbb C$; $|\lambda|>\|a\|$}
{$\lambda\ne0$ and $\|\lambda^{-1}a\|<1$}
{Neumann: $1_\A-\lambda^{-1}a\in\GL(\A)$}
{$\lambda1_\A-a\in\GL(\A)$}

\begin{cthejemplo}{(Invertibility of $2I_2-A$ outside the norm disk)}
Let us check the corollary with a matrix whose inverse can be seen
directly. Let
$A=\operatorname{diag}(1,-1)\in M_2(\mathbb C)$ with the operator norm;
$\|A\|=1$. For $\lambda=2$ we have $|\lambda|>\|A\|$ and
\[
  (2I_2-A)^{-1}=\operatorname{diag}(1,1/3).
\]
Diagonal multiplication verifies that this is a two-sided inverse.
The computation anticipates that there are no spectral values outside
the closed disk of radius $\|a\|$, a reading formalized after defining the spectrum.
\end{cthejemplo}

\section{Spectrum and resolvent}

\begin{cthdefinicion}{(Spectrum)}
Let $\A$ be a complex unital Banach algebra, with identity $1_\A$,
and let $a\in\A$. The \emph{spectrum of $a$ in $\A$} is
\[
 \spec_\A(a)=\{\lambda\in\mathbb C:\lambda1_\A-a\notin\GL(\A)\}.
\]
When there is no ambiguity, we write $\spec(a)$.
\end{cthdefinicion}
\origen{Definition 1.11}

\begin{cthobservacion}{(Spectral reading of the preceding corollary)}
Let $\A$ be a complex unital Banach algebra, $a\in\A$, and
$\lambda\in\mathbb C$. With the notation just defined, if
$|\lambda|>\|a\|$, then $\lambda\notin\sigma_\A(a)$.
\end{cthobservacion}

\begin{cthejemplo}{(Spectrum as the image of a continuous function)}
To relate the definition of spectrum to a familiar object, let
$X$ be nonempty compact Hausdorff and $f\in C(X)$. We want to decide when
$\lambda1-f$ is invertible in $C(X)$. If
$f(x)\ne\lambda$ for every $x\in X$, the function
$x\mapsto(\lambda-f(x))^{-1}$ is continuous and bounded by compactness,
and is its pointwise inverse. If $f(x_0)=\lambda$, no product with a
function in $C(X)$ can equal $1$ at $x_0$. Thus
$\spec_{C(X)}(f)=f(X)$: the spectrum coincides with the values
actually attained by the function.
\end{cthejemplo}

\pagebreak
\begin{cthproposicion}{(Compactness of the spectrum)}
Let $\A$ be a complex unital Banach algebra and $a\in\A$.
Then $\spec_\A(a)$ is compact and
\[
 \spec_\A(a)\subset\{\lambda\in\mathbb C:|\lambda|\le\|a\|\}.
\]
\end{cthproposicion}
\origen{Proposition 1.12}

\begin{proof}[Proof]
If $|\lambda|>\|a\|$, the preceding corollary gives
$\lambda1_\A-a\in\GL(\A)$. By the definition of spectrum, no
such $\lambda$ belongs to $\spec_\A(a)$; the spectrum is contained in the
closed disk of radius $\|a\|$ and is therefore bounded.
Moreover, the map $F:\mathbb C\to\A$, $F(\lambda)=\lambda1_\A-a$,
is continuous, since
$\|F(\lambda)-F(\mu)\|=|\lambda-\mu|\|1_\A\|$.
The stability proposition shows that $\GL(\A)$ is open, so
$F^{-1}(\GL(\A))=\mathbb C\setminus\spec_\A(a)$ is open.
The spectrum is closed and bounded in $\mathbb C$; by Heine--Borel, it is compact.
\end{proof}

\proofdiagram
{$\A$ unital, $a\in\A$; definition of $\spec_\A(a)$}
{Large scalars are in the resolvent set; $\GL(\A)$ is open}
{The spectrum's complement is open, and everything outside the disk lies in the resolvent set}
{$\spec(a)$ closed and bounded, hence compact}

\begin{cthejemplo}{(Spectrum of a triangular matrix)}
Let us verify the spectral bound for a matrix that is not diagonal.
Take $A=\begin{pmatrix}2&1\\0&3\end{pmatrix}\in M_2(\mathbb C)$.
The matrix $\lambda I_2-A$ is triangular and its determinant is
$(\lambda-2)(\lambda-3)$; it is therefore invertible exactly when
$\lambda\notin\{2,3\}$, and $\spec(A)=\{2,3\}$.
The compactness proposition also ensures $|\lambda|\le\|A\|$
for both values: the off-diagonal entry can change the
norm, but not these spectral values.
\end{cthejemplo}

\begin{cthdefinicion}{(Resolvent set and resolvent function)}
Let $\A$ be a complex unital Banach algebra and $a\in\A$.
The \emph{resolvent set} is $\rho_\A(a)=\mathbb C\setminus\spec_\A(a)$.
The \emph{resolvent function} is
\[
 R_a:\rho_\A(a)\to\A,\qquad R_a(z)=(z1_\A-a)^{-1}.
\]
When the ambient algebra is clear, we write $\rho(a)$.
\end{cthdefinicion}
\origen{Definition 1.13}

\noindent The conclusion of the corollary on invertibility outside the disk
can now also be written as
$|\lambda|>\|a\|\Rightarrow\lambda\in\rho_\A(a)$, because
$\rho_\A(a)=\mathbb C\setminus\spec_\A(a)$.

\begin{cthejemplo}{(Resolvent of a diagonal matrix)}
We want to compute the resolvent function of
$A=\operatorname{diag}(1,2)\in M_2(\mathbb C)$.
For $z\notin\{1,2\}$, the matrix $zI_2-A$ has nonzero diagonal entries,
and its inverse is found by inverting them separately:
\[
 R_A(z)=\operatorname{diag}\bigl((z-1)^{-1},(z-2)^{-1}\bigr).
\]
The formula is undefined precisely at the spectral values
$1$ and $2$; the resolvent thus records where invertibility fails.
\end{cthejemplo}

\begin{cthejemplo}{(The resolvent in the operator algebra)}
Let us place the same question in the algebra $\mathcal B(X)$ of a nonzero
Banach space $X$. Its identity is $I_X$, so for
$A\in\mathcal B(X)$ the preceding definitions read
\[
 \spec_{\mathcal B(X)}(A)
 =\{\lambda\in\mathbb C:\lambda I_X-A\notin\GL(\mathcal B(X))\},
 \qquad R_A(\lambda)=(\lambda I_X-A)^{-1}.
\]
Invertibility here means the existence of a \emph{bounded} linear inverse
$X\to X$. If $X=\mathbb C^2$ and
$A=\operatorname{diag}(1,2)$, we recover the calculation from the
preceding example. The abstract notation allows the same definition
to be applied to operators on infinite-dimensional spaces.
\end{cthejemplo}

\begin{cthlema}{(Resolvent identity and holomorphy)}
Let $\A$ be a complex unital Banach algebra and $a\in\A$.
The function $R_a:\rho_\A(a)\to\A$ is holomorphic and
\[
 R_a'(z)=-(z1-a)^{-2}.
\]
Moreover, for $z,w\in\rho(a)$,
\[
 R_a(z)-R_a(w)=(w-z)R_a(w)R_a(z).
\]
\end{cthlema}
\origen{Lemma 1.14}

\begin{proof}[Proof]
Write $b_z=z1_\A-a$ and $b_w=w1_\A-a$. Both are invertible
for $z,w\in\rho_\A(a)$, and $b_z-b_w=(z-w)1_\A$.
The general identity for two invertible elements,
$b_z^{-1}-b_w^{-1}=b_z^{-1}(b_w-b_z)b_w^{-1}$,
is checked because the right-hand side is
$b_z^{-1}b_wb_w^{-1}-b_z^{-1}b_zb_w^{-1}
=b_z^{-1}-b_w^{-1}$. Moreover, $b_z$
and $b_w$ commute, so their inverses commute as well. Thus
\[
 R_a(z)-R_a(w)=(w-z)R_a(z)R_a(w)
                =(w-z)R_a(w)R_a(z).
\]
To justify holomorphy as well, fix $w\in\rho_\A(a)$.
If $|z-w|\|R_a(w)\|<1$, the factorization
$b_z=b_w[1_\A+(z-w)R_a(w)]$ and Neumann give
\[
 R_a(z)=\sum_{n=0}^{\infty}(-1)^n(z-w)^nR_a(w)^{n+1}.
\]
This series converges in norm locally around $w$, so $R_a$
is holomorphic there. Dividing the preceding identity by $z-w$ and letting
$z\to w$ in the series yields
$R_a'(w)=-R_a(w)^2=-(w1_\A-a)^{-2}$.
\end{proof}

\proofdiagram
{$\A$ unital Banach, $a\in\A$; $z,w\in\rho_\A(a)$}
{Identity for inverses and local Neumann series}
{Obtain the resolvent identity and differentiate the local series}
{$R_a'(w)=-R_a(w)^2$}

\begin{cthejemplo}{(Derivative of the scalar resolvent)}
To check the sign in the lemma's formula, take
$\A=\mathbb C$ and $a=\alpha\in\mathbb C$.
For $z\ne\alpha$, $R_a(z)=(z-\alpha)^{-1}$ and
\[
 \frac{R_a(z+h)-R_a(z)}{h}
 =-\frac{1}{(z+h-\alpha)(z-\alpha)}
 \longrightarrow -(z-\alpha)^{-2}.
\]
This is the ordinary derivative of a reciprocal function; the lemma extends
the calculation to inverses in any unital Banach algebra.
\end{cthejemplo}

\begin{cthteorema}{(Nonemptiness of the spectrum)}
Let $\A$ be a complex unital Banach algebra and $a\in\A$. Then
\[
 \spec(a)\neq\varnothing.
\]
\end{cthteorema}
\origen{Theorem 1.15}

\begin{proof}[Proof]
Suppose $\spec_\A(a)=\varnothing$. Then
$R_a:\mathbb C\to\A$ is entire by the preceding lemma. In particular,
$R_a(0)=(-a)^{-1}\ne0$, since it is the inverse of an element of the
nonzero algebra. Hahn--Banach gives
$\varphi\in\A^*$ with $\varphi(R_a(0))\ne0$.
For $|z|>2\|a\|$, the Neumann series gives
\[
 R_a(z)=z^{-1}(1_\A-z^{-1}a)^{-1},\qquad
 \|R_a(z)\|\le\frac{1}{|z|}
 \sum_{n=0}^{\infty}\left(\frac{\|a\|}{|z|}\right)^n
 \le\frac{2}{|z|}.
\]
The scalar function $f=\varphi\circ R_a$ is entire and tends to zero
as $|z|\to\infty$. It is bounded outside a disk by the estimate,
and inside it by continuity and compactness. Liouville's theorem forces
$f$ to be constant; its limit at infinity makes that constant zero.
This contradicts $f(0)=\varphi(R_a(0))\ne0$. Hence
$\spec_\A(a)\ne\varnothing$.
\end{proof}

\proofdiagrambranch
{$\A$ complex unital Banach, $a\in\A$; suppose $\spec(a)=\varnothing$}
{$R_a$ would be entire and $R_a(z)\to0$ as $|z|\to\infty$}
{Hahn--Banach + Liouville applied to $\varphi\circ R_a$}
{Choose $\varphi$ with $\varphi(R_a(0))\ne0$; Liouville gives $\varphi\circ R_a=0$}
{$\spec(a)$ must be nonempty}

\begin{cthejemplo}{(Spectrum without using a determinant)}
As a familiar case, consider
$A=\begin{pmatrix}1&1\\0&1\end{pmatrix}\in M_2(\mathbb C)$.
The determinant of $\lambda I_2-A$ is $(\lambda-1)^2$, so
$\spec(A)=\{1\}$ and the spectrum is nonempty.
The nonemptiness theorem gives the same general conclusion without
using determinants, which are unavailable for an arbitrary
infinite-dimensional operator.
\end{cthejemplo}

\begin{cthnota}{(Where the hypotheses of nonemptiness are used)}
The complex scalar field allows Liouville's theorem to be applied to $\varphi\circ R_a$; completeness
guarantees the Neumann series and $\A$-valued holomorphy; Hahn--Banach separates a
nonzero value of $R_a$. The identity appears in $\lambda1-a$. Commutativity is
not needed.
\end{cthnota}

\section{Characters, maximal ideals, and the Gelfand transform}

\begin{cthteorema}{(Gelfand--Mazur rigidity)}
If $\A$ is a complex unital Banach algebra and every nonzero element
is invertible, then
$\A=\mathbb C1_\A$.
\end{cthteorema}
\origen{Theorem 1.16}

\begin{proof}[Proof]
Fix $a\in\A$. Nonemptiness of the spectrum gives
$\lambda\in\spec_\A(a)$ for some $\lambda\in\mathbb C$.
By definition, $\lambda1_\A-a$ is not invertible. Since every
nonzero element of $\A$ is invertible by hypothesis, necessarily
$\lambda1_\A-a=0$, that is, $a=\lambda1_\A$.
Thus the unital homomorphism $\mathbb C\to\A$,
$\lambda\mapsto\lambda1_\A$, is surjective; it is injective because
$1_\A\ne0$. It identifies $\A$ with $\mathbb C1_\A$.
\end{proof}

\proofdiagram
{$\A$ complex unital; every $a\in\A\setminus\{0\}$ is invertible}
{Nonemptiness of $\spec(a)$}
{Choose $\lambda\in\spec(a)$; then $\lambda1-a$ is not invertible}
{$\lambda1-a=0$, so $a=\lambda1$}

\begin{cthejemplo}{(The scalars as a division algebra)}
Let us test the theorem's hypothesis in the simplest model:
$\A=\mathbb C$ with absolute value. Every $z\ne0$ has inverse
$z^{-1}=1/z$, and the spectrum of $\alpha\in\mathbb C$ is
$\{\alpha\}$ because $\lambda-\alpha$ ceases to be invertible only
when it is zero. Gelfand--Mazur says that any complex unital
Banach algebra in which every nonzero element is invertible
is, algebraically, this same scalar model.
\end{cthejemplo}

\begin{cthdefinicion}{(Character and Gelfand spectrum)}
A \emph{character} of a commutative unital Banach algebra $\A$ is a nonzero
multiplicative linear functional $\varphi:\A\to\mathbb C$. The set of characters is denoted
\[
 \Omega(\A)=\{\varphi:\A\to\mathbb C:\varphi\text{ is a character}\},
\]
as a set. After continuity of the characters has been proved, it will be given
the weak-$*$ topology inherited from $\A^*$.
\end{cthdefinicion}
\origen{Definition 1.17}

\begin{cthejemplo}{(Point evaluation is a character)}
Let $X$ be nonempty compact Hausdorff and fix $x\in X$.
We want to check that evaluation at $x$ meets the definition of a character
of $C(X)$. The functional $\operatorname{ev}_x(f)=f(x)$ is linear and
\[
 \operatorname{ev}_x(fg)=f(x)g(x)
 =\operatorname{ev}_x(f)\operatorname{ev}_x(g),\qquad
 \operatorname{ev}_x(1)=1.
\]
It is therefore multiplicative and nonzero. Each point thus gives a
character; Assignment 2 will prove that $C(X)$ has no others.
\end{cthejemplo}

\begin{cthproposicion}{(Characters have norm one)}
Let $\A$ be a complex commutative unital Banach algebra and
$\varphi\in\Omega(\A)$. Then $\varphi$ is continuous and
\[
 \|\varphi\|=\varphi(1_\A)=1.
\]
\end{cthproposicion}
\origen{Proposition 1.18}

\begin{proof}[Proof]
Since $\varphi\ne0$, there is $b\in\A$ with $\varphi(b)\ne0$.
The equality $\varphi(b)=\varphi(1_\A b)=\varphi(1_\A)\varphi(b)$
forces $\varphi(1_\A)=1$. Fix $a\in\A$.
If $|\varphi(a)|>\|a\|$, the Neumann corollary shows that
$u=\varphi(a)1_\A-a$ is invertible. But
\[
 \varphi(u)=\varphi(a)\varphi(1_\A)-\varphi(a)=0.
\]
If $v=u^{-1}$, multiplicativity would give
$1=\varphi(1_\A)=\varphi(uv)=\varphi(u)\varphi(v)=0$,
a contradiction. Hence $|\varphi(a)|\le\|a\|$ for every
$a\in\A$. This bound proves continuity of $\varphi$ and
$\|\varphi\|\le1$; since $\|1_\A\|=\varphi(1_\A)=1$,
we obtain $\|\varphi\|\ge1$ and finally $\|\varphi\|=1$.
\end{proof}

\proofdiagram
{$\A$ commutative unital, $\varphi\in\Omega(\A)$, $a\in\A$}
{$\varphi(1)=1$; large scalars are in the resolvent set}
{If $|\varphi(a)|>\|a\|$, then $\varphi(a)1-a$ would be invertible but its image under $\varphi$ is $0$}
{$|\varphi(a)|\le\|a\|$ and $\|\varphi\|=1$}

\begin{cthejemplo}{(Norm of the evaluation functional on $C(X)$)}
Let $X$ be nonempty compact Hausdorff, fix $x\in X$, and consider
$\operatorname{ev}_x:C(X)\to\mathbb C$, $\operatorname{ev}_x(f)=f(x)$.
We want to calculate its norm. For any $f\in C(X)$,
\[
 |\operatorname{ev}_x(f)|=|f(x)|\le\|f\|_\infty,
\]
so $\|\operatorname{ev}_x\|\le1$. The constant function $1$
has uniform norm one and $\operatorname{ev}_x(1)=1$, hence
$\|\operatorname{ev}_x\|\ge1$. We conclude that
$\|\operatorname{ev}_x\|=1$, in agreement with the general proposition
on characters: evaluation at a point is a normalized continuous
functional.
\end{cthejemplo}

\begin{cthdefinicion}{(Topology of the character space)}
Let $\A$ be a complex commutative unital Banach algebra.
Since every $\varphi\in\Omega(\A)$ is continuous, $\Omega(\A)\subset\A^*$.
Its \emph{Gelfand topology} is the weak-$*$ topology
restricted from $\A^*$: a net $\varphi_i$ converges to $\varphi$ if
$\varphi_i(a)\to\varphi(a)$ for every $a\in\A$.
With this topology, $\Omega(\A)$ is compact and Hausdorff.
\end{cthdefinicion}
\origen{Definition 1.19}

\begin{proof}[Justification of compactness]
By the preceding proposition, every character has norm one, so
$\Omega(\A)$ lies in the unit ball of $\A^*$. If a net of
characters $\varphi_i$ converges weak-$*$ to $\varphi$ in that ball,
then $\varphi(1_\A)=\lim_i\varphi_i(1_\A)=1$, so
$\varphi\ne0$. Pointwise limits preserve linearity and, for
$a,b\in\A$,
\[
 \varphi(ab)=\lim_i\varphi_i(ab)
 =\lim_i\bigl(\varphi_i(a)\varphi_i(b)\bigr)
 =\varphi(a)\varphi(b).
\]
Thus $\Omega(\A)$ is closed in the weak-$*$ compact ball by
Banach--Alaoglu. The weak-$*$ topology is Hausdorff because two
distinct functionals differ on some element of $\A$.
\end{proof}

\begin{cthdefinicion}{(Two-sided ideal and quotient algebra)}
Let $\A$ be a Banach algebra. A subspace $J\subset\A$ is a
\emph{two-sided ideal} if $aJ\subset J$ and
$Ja\subset J$ for every $a\in\A$. Its classes $a+J$ form the quotient algebra
$\A/J$, with $(a+J)(b+J)=ab+J$. If $J$ is closed, the norm
\[
 \|a+J\|_{\A/J}=\inf_{j\in J}\|a+j\|
\]
makes $\A/J$ a Banach algebra; the proof of completeness is developed in
Assignment 1, Problem 4. A maximal ideal is a proper ideal that is not contained in another
larger proper ideal.
\end{cthdefinicion}

\begin{cthejemplo}{(Classes and norm in a simple quotient)}[unbreakable]
In arithmetic, $\mathbb Z/n\mathbb Z$ teaches us to identify elements
that differ by an ideal. Consider the normed version with
$X=\{p,q\}$ discrete, $\A=C(X)\cong\mathbb C^2$, and
$J=\{f:f(p)=0\}$. Two functions represent the same class
exactly when they have the same value at $p$. Moreover,
\[
 \|f+J\|_{\A/J}
 =\inf_{h\in J}\max\{|f(p)|,|f(q)+h(q)|\}
 =|f(p)|,
\]
since we can choose $h(q)=-f(q)$. The quotient retains only the
$p$ coordinate; the quotient norm measures the size of the
information that remains.
\end{cthejemplo}

\begin{cthproposicion}{(Maximal ideals and scalar quotients)}
Let $J$ be a maximal ideal of a unital Banach algebra $\A$.
Then $J$ is closed. If $\A$ is commutative, also
\[
 \A/J\cong\mathbb C.
\]
\end{cthproposicion}
\origen{Proposition 1.20}

\begin{proof}[Proof]
The closure $\overline J$ is still a two-sided ideal: if
$j_n\in J$ converges to $j$ and $a\in\A$, continuity of
multiplication gives $aj_n\to aj$ and $j_na\to ja$, both in
$\overline J$. If $J$ were not closed, maximality would require
$\overline J=\A$. Then there would be $j\in J$ with
$\|1_\A-j\|<1$. Neumann would make $j$ invertible, and
$1_\A=j^{-1}j\in J$, contradicting that $J$ is proper.
Hence $J$ is closed and $\A/J$ is a nonzero unital Banach algebra.
Its identity is $1_\A+J$ and has norm one: the quotient norm is at most
$\|1_\A\|=1$, while a distance less than one from $1_\A$ to $J$
would, by Neumann, produce an invertible element of $J$ and force
$J=\A$.

Now suppose $\A$ is commutative. The ideals of $\A/J$
correspond, under the quotient map, to the ideals
of $\A$ that contain $J$; by maximality, $\A/J$ has no
ideals other than $\{0\}$ and itself. Let $a+J\ne0$.
The ideal $(\A/J)(a+J)$ is nonzero and thus equals all of $\A/J$.
There is $b+J$ with $(b+J)(a+J)=1_\A+J$. Commutativity also gives
$(a+J)(b+J)=1_\A+J$, so every nonzero element
of the quotient is invertible. Gelfand--Mazur yields
$\A/J\cong\mathbb C$ as complex unital algebras.
\end{proof}

\proofdiagrambranch
{$\A$ unital Banach, $J\subset\A$ maximal}
{In the second assertion, $\A$ commutative}
{The closure of $J$ is still an ideal; $\A/J$ has no nontrivial ideals}
{Every nonzero element of $\A/J$ is invertible; apply Gelfand--Mazur}
{$J$ closed and $\A/J\cong\mathbb C$}

\begin{cthejemplo}{(A maximal ideal defined by evaluation)}[unbreakable]
Let $X$ be nonempty compact Hausdorff and fix $x\in X$.
We want to describe the quotient that forgets functions vanishing at $x$.
The set $J_x=\{f\in C(X):f(x)=0\}$ is the kernel of the character
$\operatorname{ev}_x$, and the map
\[
 C(X)/J_x\longrightarrow\mathbb C,\qquad f+J_x\longmapsto f(x)
\]
is well defined, multiplicative, and surjective because it includes
constant functions. Its kernel is zero, so it is an
isomorphism. Since the quotient is a field, $J_x$ is maximal:
the value at a point is all the information that survives.
\end{cthejemplo}

\begin{investigacion}
\textbf{Direct construction example.}
Let $X$ be an infinite-dimensional complex Banach space.
An operator $K\in\mathcal B(X)$ is \emph{compact} if
$K$ maps the unit ball to a relatively compact set. The set
$\mathcal K(X)$ of such operators is a closed two-sided ideal of $\mathcal B(X)$:
composition with a bounded operator preserves compactness, and the norm limit
of compact operators is compact. The quotient
\[
 \mathcal Q(X):=\mathcal B(X)/\mathcal K(X)
\]
is called the \emph{Calkin algebra} \cite[p.~2]{Karlovich2025}.
If $X$ is a Hilbert space, this quotient is also a
$C^*$-algebra; for a general Banach space, only its Banach algebra structure is used here.
\textbf{What to recognize:} $\mathcal K(X)$ plays exactly the role of the
closed ideal $J$ in the construction $\A/J$; $\mathcal Q(X)$ is the concrete quotient.
\end{investigacion}

\begin{investigacion}
\textbf{Connection with the literature: use of a quotient.}
Let $X$ be an infinite-dimensional complex Banach space.
An operator $A\in\mathcal B(X)$ is \emph{Fredholm} if its kernel is finite-dimensional,
its range is closed, and the codimension of its range is finite. The equivalence from
operator theory
\[
 A\text{ Fredholm}\quad\Longleftrightarrow\quad
 A+\mathcal K(X)\in\GL(\mathcal Q(X))
\]
identifies an operator question with a question of invertibility in a
quotient \cite[p.~2]{Karlovich2025}. In $X=\ell^2$, let $P$ be the rank-one projection
onto $\mathbb Ce_1$ and $A=I_X-P$. Then $A$ is not invertible, but
$A+\mathcal K(X)=I_X+\mathcal K(X)$ is: the difference $P$ disappears
in the quotient. This is an instructional elaboration based on the cited criterion.
\textbf{What to recognize:} the quotient $\A/J$ from Problem 1.4 reappears with
$\A=\mathcal B(X)$ and $J=\mathcal K(X)$; the equivalence with the Fredholm property is
an additional result from operator theory.
\end{investigacion}

\pagebreak
\begin{cthproposicion}{(Characters and maximal ideals)}
Let $\A$ be a complex commutative unital Banach algebra.
The assignment
\[
 \Omega(\A)\longrightarrow\{\text{maximal ideals of }\A\},\qquad
 \varphi\longmapsto\ker\varphi,
\]
is a bijection.
\end{cthproposicion}
\origen{Proposition 1.21}

\begin{proof}[Proof]
If $\varphi\in\Omega(\A)$, then $\varphi(1_\A)=1$ and
$\varphi$ is surjective because $\varphi(\lambda1_\A)=\lambda$
for every $\lambda\in\mathbb C$. Its kernel is a proper ideal, and the
first isomorphism theorem gives
$\A/\ker\varphi\cong\mathbb C$; since $\mathbb C$ has no
nonzero proper ideals, $\ker\varphi$ is maximal.

Conversely, if $J$ is maximal, the preceding proposition allows us to
choose a complex unital isomorphism
$\theta:\A/J\to\mathbb C$. Let
$q:\A\to\A/J$ be the quotient map. The composition
$\varphi=\theta\circ q$ is linear, multiplicative, and unital;
it is a character with $\ker\varphi=J$.
To check uniqueness, if $\psi$ is another character with the same
kernel, then for each $a\in\A$,
$a-\varphi(a)1_\A\in J=\ker\psi$. Therefore
$\psi(a)=\varphi(a)\psi(1_\A)=\varphi(a)$. Thus the two assignments
are inverses of one another.
\end{proof}

\proofdiagram
{$\A$ commutative unital; $\varphi\in\Omega(\A)$ or $J\subset\A$ maximal}
{Quotients $\A/J\cong\mathbb C$}
{A character's kernel has codimension $1$; conversely, use the quotient map}
{$\varphi\leftrightarrow\ker\varphi$}

\begin{cthejemplo}{(Recovering two points from ideals)}
To see the bijection between characters and maximal ideals,
take $X=\{p,q\}$ and $C(X)\cong\mathbb C^2$.
Evaluations give the distinct ideals
$J_p=\{f:f(p)=0\}$ and $J_q=\{f:f(q)=0\}$.
If $\varphi$ is any character, the indicator functions
$e_p,e_q$ satisfy $e_p+e_q=1$ and $e_pe_q=0$;
their images are scalar idempotents and exactly one equals $1$.
Since $f=f(p)e_p+f(q)e_q$, it follows that $\varphi$ is
$\operatorname{ev}_p$ or $\operatorname{ev}_q$.
In this case, the two points of the space are recovered as the
two maximal ideals of the algebra.
\end{cthejemplo}

\begin{cthteorema}{(The spectrum is read through characters)}
Let $\A$ be a complex commutative unital Banach algebra.
Then, for each $a\in\A$,
\[
 \spec(a)=\{\varphi(a):\varphi\in\Omega(\A)\}.
\]
\end{cthteorema}
\origen{Theorem 1.22}

\begin{proof}[Proof]
If $\lambda=\varphi(a)$ for a character $\varphi$, then
$\varphi(\lambda1_\A-a)=\lambda-\varphi(a)=0$.
An invertible element cannot have zero image under a unital character;
therefore $\lambda1_\A-a$ is not invertible and
$\lambda\in\spec_\A(a)$.

Conversely, let $\lambda\in\spec_\A(a)$ and put
$b=\lambda1_\A-a$. Since $b$ is not invertible, the principal ideal
$\A b$ is proper: if $1_\A=cb$ for some $c\in\A$,
commutativity would also give $bc=1_\A$. By Zorn's lemma,
$\A b$ is contained in a maximal ideal $J$: indeed, the union
of a chain of proper ideals containing $\A b$ remains
proper, since if it contained $1_\A$, that element would already belong
to one of its members. That union is an upper bound for the chain.
The preceding proposition gives a character $\varphi$ with $J=\ker\varphi$.
Since $b\in J$, we obtain
$0=\varphi(b)=\lambda-\varphi(a)$, whence
$\lambda=\varphi(a)$. Both inclusions are proved.
\end{proof}

\proofdiagrambranch
{$a\in\A$, $\varphi\in\Omega(\A)$ and $\lambda=\varphi(a)$}
{$\lambda\in\spec_\A(a)$}
{Characters $\leftrightarrow$ maximal ideals}
{$\lambda1-a$ lies in a maximal ideal exactly when some character annihilates it}
{$\lambda\in\spec(a)\Longleftrightarrow\exists\varphi\in\Omega(\A):\lambda=\varphi(a)$}

\begin{cthejemplo}{(Spectrum of the square function on $[0,1]$)}
Apply the spectral theorem via characters to
$X=[0,1]$ and $f(x)=x^2\in C(X)$.
Each evaluation $\operatorname{ev}_x$ takes on $f$ the value $x^2$,
and the pointwise invertibility criterion already checked for $C(X)$
rules out other spectral values. Therefore
\[
 \spec_{C(X)}(f)
 =\{\operatorname{ev}_x(f):x\in[0,1]\}
 =\{x^2:x\in[0,1]\}=[0,1].
\]
The spectrum recovers exactly the geometric range of the function.
\end{cthejemplo}

\begin{cthproposicion}{(Jacobson's lemma for $1-ab$)}
Let $\A$ be a unital Banach algebra with identity $1_\A$ and
$a,b\in\A$. Then
\[
 1_\A-ab\in\GL(\A)\quad\Longleftrightarrow\quad 1_\A-ba\in\GL(\A).
\]
If $1_\A-ab$ is invertible,
\[
 (1_\A-ba)^{-1}=1_\A+b(1_\A-ab)^{-1}a.
\]
\end{cthproposicion}
\origen{Proposition 1.23}

\begin{proof}[Proof]
Suppose $1_\A-ab$ is invertible and write
$c=(1_\A-ab)^{-1}$. The two inverse identities give
$c-abc=1_\A$ and $c-cab=1_\A$.
Define $d=1_\A+bca$. On the left,
\[
 (1_\A-ba)d
 =1_\A-ba+bca-babca
 =1_\A+b(c-1_\A-abc)a=1_\A.
\]
On the right, use the other identity for $c$:
\[
 d(1_\A-ba)
 =1_\A-ba+bca-bcaba
 =1_\A+b(c-1_\A-cab)a=1_\A.
\]
Thus $d$ is the two-sided inverse of $1_\A-ba$, proving
the formula in the statement. Starting instead from invertibility of
$1_\A-ba$, the same reasoning with $a$ and $b$ interchanged
proves the reverse implication.
\end{proof}

\proofdiagram
{$\A$ unital, $a,b\in\A$ and $1_\A-ab$ invertible}
{Associativity}
{Explicitly construct $1+b(1-ab)^{-1}a$}
{It is a two-sided inverse of $1-ba$}

\begin{cthejemplo}{(Jacobson's identity for noncommuting matrices)}[unbreakable]
We want to see the proposition without assuming the factors commute.
In $M_2(\mathbb C)$ take $A=E_{12}$ and
$B=\frac12E_{21}$, where $E_{ij}$ are the matrices with a
single entry $1$ in position $(i,j)$. Then
$AB=\frac12E_{11}$ and $BA=\frac12E_{22}$ are distinct, but
\[
 (I_2-AB)^{-1}=\operatorname{diag}(2,1),\qquad
 (I_2-BA)^{-1}=\operatorname{diag}(1,2).
\]
Jacobson's formula gives the second inverse from the first:
$I_2+B(I_2-AB)^{-1}A=I_2+E_{22}$.
Simultaneous invertibility follows from the structure of the product,
without requiring $AB=BA$.
\end{cthejemplo}

\begin{cthcorolario}{(Nonzero spectra of $ab$ and $ba$)}
Let $\A$ be a complex unital Banach algebra and $a,b\in\A$.
Then
\[
 \spec_\A(ab)\cup\{0\}=\spec_\A(ba)\cup\{0\},
 \qquad\text{equivalently}\qquad
 \spec_\A(ab)\setminus\{0\}=\spec_\A(ba)\setminus\{0\}.
\]
\end{cthcorolario}
\origen{Corollary 1.24}

\begin{proof}[Proof]
Let $\lambda\in\mathbb C\setminus\{0\}$ and use
the factorization
$\lambda1_\A-ab=\lambda(1_\A-\lambda^{-1}ab)$.
The scalar factor $\lambda1_\A$ is invertible, so
\[
\lambda1-ab\in\GL(\A)
\Longleftrightarrow
1-\lambda^{-1}ab\in\GL(\A)
\Longleftrightarrow
1-\lambda^{-1}ba\in\GL(\A)
\Longleftrightarrow
\lambda1-ba\in\GL(\A).
\]
The middle equivalence is the preceding lemma applied to
$\lambda^{-1}a$ and $b$; the last again uses the scalar factor
$\lambda\ne0$. By the definition of spectrum, $\lambda$ belongs
to $\spec_\A(ab)$ exactly when it belongs to
$\spec_\A(ba)$. Equality of the spectra away from zero
follows for every $\lambda\ne0$; adding $\{0\}$ to both
sets gives the first equality in the statement.
\end{proof}

\proofdiagram
{$a,b\in\A$, $\lambda\in\mathbb C\setminus\{0\}$}
{Jacobson's lemma}
{$\lambda1_\A-ab=\lambda(1_\A-\lambda^{-1}ab)$}
{The nonzero spectra coincide; adding $\{0\}$ gives the equality}

\begin{cthejemplo}{(Reversing the order across different dimensions)}
Let $A:\mathbb C^m\to\mathbb C^n$ and
$B:\mathbb C^n\to\mathbb C^m$. We want to compare the products
$AB\in M_n(\mathbb C)$ and $BA\in M_m(\mathbb C)$.
Jacobson's identity adapts to the two identities and, for
$\lambda\ne0$, gives
\[
 (I_m-\lambda^{-1}BA)^{-1}
 =I_m+\lambda^{-1}B(I_n-\lambda^{-1}AB)^{-1}A
\]
 if $I_n-\lambda^{-1}AB$ is invertible. As a concrete case, take
$A(t)=(t,0)$ from $\mathbb C$ to $\mathbb C^2$ and
$B(x,y)=x$ from $\mathbb C^2$ to $\mathbb C$. Then
$AB=\operatorname{diag}(1,0)$ and $BA=(1)$:
their nonzero spectra are both $\{1\}$, although only the first
contains $0$. This explains why the corollary adds $\{0\}$.
\end{cthejemplo}

\begin{cthdefinicion}{(Spectral radius)}
Let $\A$ be a complex unital Banach algebra and $a\in\A$.
The \emph{spectral radius} of $a$ is
\[
 r(a)=\max\{|\lambda|:\lambda\in\spec_\A(a)\}.
\]
\end{cthdefinicion}
\origen{Definition 1.25}

\begin{cthejemplo}{(Spectral radius and matrix norm)}
For $A=\operatorname{diag}(2,-3,i)\in M_3(\mathbb C)$,
the spectral values are its diagonal entries, so
$r(A)=\max\{2,3,1\}=3$. With the Euclidean operator norm, also
$\|A\|=3$, since it is diagonal. Equality is not automatic for
nonnormal matrices: if
$B=\begin{pmatrix}2&10\\0&-3\end{pmatrix}$,
its spectrum is still $\{2,-3\}$ and $r(B)=3$, whereas
$\|B\|\ge\|Be_2\|=\sqrt{109}>3$.
The radius measures spectral values and can be less than the norm.
\end{cthejemplo}

\begin{cthteorema}{(Spectral radius formula)}
Let $\A$ be a complex unital Banach algebra and $a\in\A$.
Then
\[
 r(a)=\lim_{n\to\infty}\|a^n\|^{1/n}
     =\inf_{n\ge1}\|a^n\|^{1/n}.
\]
\end{cthteorema}
\origen{Theorem 1.26}

\begin{proof}[Proof]
First check that the limit of norms exists. Put
$\beta_n=\|a^n\|$. Submultiplicativity gives
$\beta_{m+n}\le\beta_m\beta_n$. If $\beta_m=0$ for some
$m\ge1$, then $a^n=0$ for every $n\ge m$ and
$\lim_n\beta_n^{1/n}=\inf_{n\ge1}\beta_n^{1/n}=0$.
Otherwise, fix $m\ge1$ and write $n=qm+s$,
$0\le s<m$. Since
$\beta_n\le\beta_m^q\beta_s$ and there are only finitely many
possible $s$, taking $n$th roots and letting $n\to\infty$
gives
\[
 \limsup_{n\to\infty}\beta_n^{1/n}\le\beta_m^{1/m}.
\]
This holds for every $m$. On the other hand, each term
$\beta_n^{1/n}$ is at least $\inf_{m\ge1}\beta_m^{1/m}$.
Consequently,
\[
 L:=\lim_{n\to\infty}\|a^n\|^{1/n}
   =\inf_{n\ge1}\|a^n\|^{1/n}.
\]

Let us show that $r(a)\le L$. If $\lambda\in\spec_\A(a)$ and
$n\ge1$, the polynomial factorization
\[
 \lambda^n1_\A-a^n
 =(\lambda1_\A-a)
   \sum_{k=0}^{n-1}\lambda^{n-1-k}a^k
 \]
has commuting factors. If the product were invertible, each
factor would be: indeed, if $x$ and $y$ commute and $xy$ has
inverse $c$, the equality $x(xy)=(xy)x$ implies $xc=cx$ after
multiplication by $c$. Then
$x(yc)=(xy)c=1_\A$ and
$(yc)x=y(cx)=yx c=(xy)c=1_\A$, so $yc$ is
a two-sided inverse of $x$.
This contradicts the noninvertibility of $\lambda1_\A-a$. Thus
$\lambda^n\in\spec_\A(a^n)$, and the spectral bound already proved gives
$|\lambda|^n\le\|a^n\|$. Taking the maximum over the spectrum
and then taking roots yields $r(a)\le\|a^n\|^{1/n}$ for every
$n$, and therefore $r(a)\le L$.

For the other inequality, let $R=1/r(a)$ if $r(a)>0$ and
$R=\infty$ if $r(a)=0$. The function
$g(z)=(1_\A-za)^{-1}$ is defined and holomorphic on
$|z|<R$: for $z\ne0$ in that disk,
$|z^{-1}|>r(a)$ and
$1_\A-za=z(z^{-1}1_\A-a)$ is invertible; at $z=0$
it is the identity. If $z\ne0$, moreover,
$g(z)=z^{-1}R_a(z^{-1})$, which is holomorphic by the resolvent
lemma and holomorphy of $z\mapsto z^{-1}$ away from zero.
Near zero, Neumann gives
$g(z)=\sum_{n\ge0}a^nz^n$.
This series converges in norm for $|z|\|a\|<1$ and proves
holomorphy at zero as well.

Fix a real $t$ with $0<t<R$. For each $\varphi\in\A^*$,
$\varphi\circ g$ is holomorphic on $|z|<R$, and its Taylor series
at zero has coefficients $\varphi(a^n)$. Choose $s$ with
$t<s<R$. Cauchy's bound on the circle $|z|=s$ gives
$|\varphi(a^n)|\le M_{\varphi,s}s^{-n}$, where
$M_{\varphi,s}=\max_{|z|=s}|\varphi(g(z))|<\infty$.
Thus $t^n|\varphi(a^n)|\le M_{\varphi,s}(t/s)^n$ is
bounded for each $\varphi$. Apply Banach--Steinhaus on the Banach space
$\A^*$ to the continuous operators
$T_n:\A^*\to\mathbb C$,
$T_n(\varphi)=t^n\varphi(a^n)$. Pointwise boundedness implies
$\sup_n\|T_n\|<\infty$. Hahn--Banach gives
$\|T_n\|=t^n\|a^n\|$, so there is $C_t<\infty$ with
$t^n\|a^n\|\le C_t$ for every $n$. Taking roots and limits,
$L\le t^{-1}$. If $r(a)>0$, let $t\uparrow1/r(a)$ and
obtain $L\le r(a)$; if $r(a)=0$, we may take $t$
arbitrarily large and obtain $L=0$. Together with the first
inequality and the existence of $L$, this proves both equalities.
\end{proof}

\proofdiagrambranch
{$\A$ unital, $a\in\A$; $\lambda\in\spec(a)\Rightarrow\lambda^n\in\spec(a^n)$}
{Study $g(z)=(1-za)^{-1}$ on $|z|<R$, $R=1/r(a)$ or $\infty$}
{Holomorphy + Taylor + Banach--Steinhaus}
{Bound $t^n\|a^n\|$ for each $t<R$ and let $t\uparrow R$}
{$r(a)=\lim\|a^n\|^{1/n}$}

\begin{cthejemplo}{(The Volterra operator is quasinilpotent)}[unbreakable]
Let us study a nonzero operator whose spectral radius is zero.
On $C[0,1]$ with the uniform norm, define
$(Vf)(x)=\int_0^x f(t)\,dt$. Iterating the integral gives, for
$n\ge1$,
\[
 (V^nf)(x)=\frac1{(n-1)!}\int_0^x(x-t)^{n-1}f(t)\,dt,
 \qquad \|V^n\|\le\frac1{n!}.
\]
The bound follows by replacing $|f(t)|$ with $\|f\|_\infty$
and integrating the power.
By the spectral radius formula,
$r(V)=\lim_n\|V^n\|^{1/n}=0$.
Since the spectrum is nonempty and all its points have modulus
at most zero, $\spec(V)=\{0\}$.
This case distinguishes a quasinilpotent operator from the zero operator:
$V\ne0$ although its spectrum is $\{0\}$.
\end{cthejemplo}

\begin{cthnota}{(Interpretation of the formula for iterates)}
Let $X$ be a Banach space and $A\in\mathcal B(X)$. Then $\|A^n\|$
measures the largest effect of $n$ successive
applications. The limit $\lim_n\|A^n\|^{1/n}$ measures its asymptotic exponential rate and
equals $r(A)$. The Volterra example shows that this rate may be zero
even when $A\ne0$.
\end{cthnota}

\begin{cthdefinicion}{(Gelfand transform)}
Let $\A$ be a complex commutative unital Banach algebra.
The \emph{Gelfand transform} is
\[
 \Gamma:\A\to C(\Omega(\A)),\qquad
 a\mapsto\widehat a,\qquad \widehat a(\varphi)=\varphi(a).
\]
\end{cthdefinicion}
\origen{Definition 1.27}

\begin{cthejemplo}{(The transform recovers a function in $C(X)$)}
Take $X=[0,1]$ and $f(t)=t\in C(X)$ to make the definition concrete.
Each point $x\in X$ gives the evaluation character
$\operatorname{ev}_x$, and
\[
 \widehat f(\operatorname{ev}_x)
 =\operatorname{ev}_x(f)=x=f(x).
\]
Under the identification $\Omega(C(X))\cong X$, proved in Assignment 2,
the Gelfand transform is the function $f$ itself.
The formula shows how a function is read through all the
characters of the algebra.
\end{cthejemplo}

\begin{cthproposicion}{(The Gelfand transform is contractive)}
Let $\A$ be a complex commutative unital Banach algebra and
$a\in\A$. The Gelfand transform is linear, multiplicative,
unital, and
\[
 \|\widehat a\|_\infty=r(a)\le\|a\|.
\]
\end{cthproposicion}
\origen{Remark 1.28}

\begin{proof}[Proof]
For each $\varphi\in\Omega(\A)$, linearity and
multiplicativity of the character give
\[
 \widehat{\alpha a+b}(\varphi)
   =\alpha\widehat a(\varphi)+\widehat b(\varphi),\qquad
 \widehat{ab}(\varphi)=\widehat a(\varphi)\widehat b(\varphi),
 \qquad \widehat{1_\A}(\varphi)=1.
\]
Each $\widehat a$ is continuous in the weak-$*$ topology by its
definition. The theorem describing the spectrum through
characters identifies the range of $\widehat a$
with $\spec_\A(a)$. Hence
\[
 \|\widehat a\|_\infty=\sup_{\varphi\in\Omega(\A)}|\varphi(a)|
 =\max_{\lambda\in\spec(a)}|\lambda|=r(a)\le\|a\|.
\]
The maximum exists because the spectrum is compact and nonempty;
the last inequality is the spectral bound proved earlier.
\end{proof}

\proofdiagram
{Definition $\widehat a(\varphi)=\varphi(a)$}
{Linear and multiplicative characters; spectrum via characters}
{$\|\widehat a\|_\infty=\sup_{\varphi}|\varphi(a)|$}
{$\|\widehat a\|_\infty=r(a)\le\|a\|$}

\begin{investigacion}[unbreakable]
\textbf{Connection with the literature: a direct example of a commutative algebra.}
Karlovich considers $\SO(\mathbb R_+)$, the algebra of bounded continuous functions whose
oscillation on $[r,2r]$ tends to zero as $r\to0$ and as $r\to\infty$
\cite[pp.~465--466, §2]{Karlovich2017}. It is a commutative $C^*$-subalgebra of
$C_b(\mathbb R_+)$: it contains constants, is closed under addition, multiplication, and conjugation, and
the oscillation condition is preserved under uniform limits. The
$C^*$ identity will be recognized formally in Section 2; at this point all that is
needed is that the algebra is commutative and unital. For $a$ in this algebra,
$\widehat a(\xi)=\xi(a)$ evaluates $a$ at a character $\xi\in\Omega(\SO(\mathbb R_+))$.
Besides evaluations at $r>0$, the maximal ideal space has fibers over $0$ and
$\infty$. These fibers organize limiting behaviors that a single value $a(0)$ or
$a(\infty)$ would not describe. The Gelfand transform of this coefficient algebra
is a conceptual step toward Fredholm symbols, not the full Fredholm symbol.
\textbf{What to recognize:} the characters of $\A$ studied in the course
become ordinary evaluations and characters in the fibers of
$\Omega(\SO(\mathbb R_+))$ here. For the precise construction, see
\cite[§2, Propositions~2.1--2.2]{Karlovich2017}.
\end{investigacion}

\begin{cthdefinicion}{(Modular ideal)}
An ideal $I$ of an algebra $\A$ is called \emph{modular} if there is $u\in\A$ such that
\[
 a-ua\in I,\qquad a-au\in I,\qquad a\in\A.
\]
\end{cthdefinicion}
\origen{Definition 1.29}

\begin{cthejemplo}{(Local identity for an evaluation ideal)}
Let $X$ be locally compact Hausdorff and fix $x\in X$.
We want to check the modularity of the ideal
\[
 I_x=\{f\in C_0(X):f(x)=0\}.
\]
Choose a continuous compactly supported function $u$ with
$u(x)=1$, available by local compactness and the Hausdorff separation property;
in particular, $u\in C_0(X)$.
For every $f\in C_0(X)$ we have
$(f-uf)(x)=(f-fu)(x)=0$; therefore both differences belong to
$I_x$. The class $u+I_x$ acts as the identity of the quotient, although
$C_0(X)$ may lack a global identity.
\end{cthejemplo}

\begin{cthproposicion}{(Modularity criterion via the quotient)}
Let $\A$ be an algebra and $I\subsetneq\A$ a proper two-sided ideal.
Then $I$ is modular if and only if $\A/I$ has an identity.
\end{cthproposicion}
\origen{Proposition 1.31}

\begin{proof}[Proof]
If $I$ is modular, there is $u\in\A$ such that for every
$a\in\A$, we have $a-au\in I$ and $a-ua\in I$. In the
quotient, this is equivalent to
\[
 (a+I)(u+I)=au+I=a+I=ua+I=(u+I)(a+I).
\]
Therefore $u+I$ is a two-sided identity of $\A/I$.
Conversely, if $\A/I$ has an identity, it is a class
$u+I$ for some $u\in\A$. For each $a\in\A$, the
identity laws give $au+I=a+I=ua+I$, that is,
$a-au\in I$ and $a-ua\in I$. This $u$ proves that $I$ is
modular.
\end{proof}

\proofdiagram
{$I\subsetneq\A$ a proper two-sided ideal}
{Definition of multiplication in $\A/I$}
{$u$ modulo $I$ acts as the identity exactly when $a-au,a-ua\in I$}
{$I$ modular $\Longleftrightarrow \A/I$ unital}

\begin{cthejemplo}{(A unital quotient of $C_0(\mathbb R)$)}
To see the proposition in a nonunital algebra, take
$\A=C_0(\mathbb R)$ and
$I_0=\{f\in\A:f(0)=0\}$. The constant function $1$ does not belong
to $\A$, but $u(t)=\max\{1-|t|,0\}$ does and satisfies $u(0)=1$.
Evaluation induces
\[
 \A/I_0\longrightarrow\mathbb C,\qquad f+I_0\longmapsto f(0),
\]
a surjective isomorphism because the functions $\lambda u$ attain
any value $\lambda$. Thus $\A/I_0$ has identity $u+I_0$:
modularity detects an identity that appears only after taking the quotient.
\end{cthejemplo}

\chapter{\texorpdfstring{$C^*$}{C*}-algebras: spectral rigidity and functional calculus}

\section{The \texorpdfstring{$C^*$}{C*}-identity and special classes of elements}

\begin{cthdefinicion}{($C^*$-algebra)}
A $C^*$-algebra is an involutive Banach algebra $\A$ satisfying
\[
 \|a^*a\|=\|a\|^2,
 \qquad a\in\A.
\]
\end{cthdefinicion}
\origen{Definition 2.1}

\begin{cthejemplo}{(Operators on a Hilbert space and the $C^*$-identity)}
We seek a concrete model of the definition. If $H$ is a
Hilbert space, $\mathcal B(H)$ is the algebra of bounded linear operators
$H\to H$, with involution given by the adjoint $T\mapsto T^*$.
For $T\in\mathcal B(H)$, the inequality
$\|T^*T\|\le\|T^*\|\|T\|=\|T\|^2$ and the identity
\[
 \|T\|^2=\sup_{\|x\|=1}\langle T^*Tx,x\rangle
 \le\|T^*T\|
\]
give $\|T^*T\|=\|T\|^2$. Thus the operator norm and the adjoint
satisfy precisely the $C^*$-identity.
\end{cthejemplo}

\begin{cthobservacion}{(The identity is compatible with the involution)}
If a $C^*$-algebra $\A$ has an identity $1_\A$, then
$1_\A^*=1_\A$ and $\|1_\A\|=1$.
\end{cthobservacion}
\origen{Observation 2.2}

\begin{proof}[Proof]
For every $a\in\A$, the identity $(a1_\A)^*=1_\A^*a^*=a^*$ shows that
$1_\A^*$ acts as a left identity. Likewise, it acts as a right identity;
by uniqueness of the identity, $1_\A^*=1_\A$. The $C^*$-identity gives
$\|1_\A\|^2=\|1_\A^*1_\A\|=\|1_\A\|$, and $1_\A\ne0$ implies
$\|1_\A\|=1$.
\end{proof}

\begin{cthejemplo}{(A self-adjoint identity of norm one in matrix algebras)}
Let us illustrate the preceding observation with the identity of
$M_n(\mathbb C)$, $n\ge1$, equipped with the operator norm and the
conjugate-transpose adjoint. The matrix $I_n$ satisfies
$I_n^*=I_n$; moreover, $\|I_nx\|=\|x\|$ for every vector, so
$\|I_n\|=1$. The abstract assertion about $1_\A$ becomes here
a visible property of the identity matrix.
\end{cthejemplo}

\begin{investigacion}
\textbf{Direct example.} For the open logarithmic star $\Gamma$ in the Haseman
problem, $C_b(\Gamma)$, with the uniform norm and pointwise conjugation, is a
unital $C^*$-algebra; the subalgebra $\SO(\Gamma)$ consists of functions that
oscillate slowly at $0$ and $\infty$ \cite[§1, p.~464]{Karlovich2017}.
\textbf{What to recognize:} both are function algebras of the type just
defined; this observation only illustrates the $C^*$-identity and plays no
part in the course proofs.
\end{investigacion}

\begin{cthdefinicion}{($*$-homomorphism)}
Let $\A,\B$ be $C^*$-algebras. A $*$-homomorphism is a linear, multiplicative
homomorphism $\Phi:\A\to\B$ satisfying
\[
 \Phi(a^*)=\Phi(a)^*\qquad(a\in\A).
\]
If both algebras are unital, we say that $\Phi$ is unital if $\Phi(1_\A)=1_\B$.
\end{cthdefinicion}
\origen{Definition 2.4}

\begin{cthejemplo}{(Functions as multiplication operators)}
Let $X$ be locally compact and Hausdorff. We want to represent the algebra
$C_b(X)$ by operators on $\ell^2(X)$, the space of square-summable complex
families indexed by $X$. For $f\in C_b(X)$ we define
$(M_f\xi)(x)=f(x)\xi(x)$. The estimate
$\|M_f\xi\|_2\le\|f\|_\infty\|\xi\|_2$ proves that $M_f$ is bounded.
Moreover, $M_fM_g=M_{fg}$ and, by comparing inner products,
$M_f^*=M_{\overline f}$; also $M_1=I$.
Hence $f\mapsto M_f$ is a unital $*$-homomorphism.
Problem 2.1 checks that it even preserves the norm: a function algebra
can be represented faithfully as an operator algebra.
\end{cthejemplo}

\begin{cthdefinicion}{(Normal, self-adjoint, unitary, and projection)}
Let $\A$ be a $C^*$-algebra and $a\in\A$.
\begin{enumerate}[label=(\roman*)]
\item $a$ is \emph{normal} if $a^*a=aa^*$;
\item $a$ is \emph{self-adjoint} if $a^*=a$; we write
$\A_{\mathrm{sa}}=\{x\in\A:x=x^*\}$;
\item if $\A$ has an identity $1_\A$, $a$ is \emph{unitary} if
$a^*a=aa^*=1_\A$; we write $\mathcal U(\A)$ for the set of unitaries;
\item $a$ is a \emph{projection} if $a=a^*=a^2$; we write
$\mathcal P(\A)$ for the set of projections.
\end{enumerate}
\end{cthdefinicion}
\origen{Definition 2.5}

\begin{cthejemplo}{(Normality, self-adjointness, unitarity, and projection in $M_2$)}
Let us find a representative of each class in the definition in
$M_2(\mathbb C)$:
\[
\begin{array}{c|c}
\text{type}&\text{example}\\ \hline
\text{normal}&\begin{pmatrix}1&0\\0&i\end{pmatrix}\\[2mm]
\text{self-adjoint}&\begin{pmatrix}1&0\\0&-2\end{pmatrix}\\[2mm]
\text{unitary}&\begin{pmatrix}0&1\\1&0\end{pmatrix}\\[2mm]
\text{projection}&\begin{pmatrix}1&0\\0&0\end{pmatrix}
\end{array}
\]
The two diagonal matrices commute with their adjoints; the second
equals its adjoint. For the third, $U^*=U$ and $U^2=I_2$,
so $U^*U=UU^*=I_2$. The fourth satisfies
$P=P^*=P^2$. These checks show why every
self-adjoint element, unitary, or projection is normal, although the first
normal matrix is not self-adjoint.
\end{cthejemplo}

\begin{investigacion}
\textbf{Direct example in $L^2$; an interpretive analogy for Haseman.}
For $\Gamma_0=\mathbb R$ or the circle
$\mathbb T=\{z\in\mathbb C:|z|=1\}$, with arc-length measure on $L^2(\Gamma_0)$,
the Cauchy operator satisfies $S_{\Gamma_0}^2=I$ and defines
\[
 P_{\Gamma_0}^\pm=\frac12(I\pm S_{\Gamma_0}).
\]
The identity $S_{\Gamma_0}^2=I$ gives $(P_{\Gamma_0}^\pm)^2=P_{\Gamma_0}^\pm$ by
algebraic calculation. In this model, moreover, $S_{\Gamma_0}$ is self-adjoint and
$(P_{\Gamma_0}^\pm)^*=P_{\Gamma_0}^\pm$: these are orthogonal projections
\cite[Sec.~2, p.~22]{BFK2006}. The formula
$P_\pm=(I\pm S_\Gamma)/2$ reappears in the Haseman operator on
$L^p(\Gamma)$ for an open logarithmic star $\Gamma$, with
$1<p<\infty$ and data satisfying the hypotheses of \cite{Karlovich2017},
\[
 T=V_\alpha P_+ + GP_-\in\mathcal B(L^p(\Gamma)),
\]
which applies the shift $V_\alpha f=f\circ\alpha$ to the $P_+f$ term and
multiplication by $G$ to the $P_-f$ term. On the open star,
$S_\Gamma^2=I$ and orthogonality do not follow from the $\Gamma_0$ case
\cite[(1.2)]{Karlovich2017}.
\textbf{What to recognize:} $S^2=I$ ensures idempotence in the Hilbert-space
model and $S^*=S$ adds orthogonality; the Haseman formula is read term by
term without transferring these properties to the open star.
\end{investigacion}

\pagebreak
\section{Norm and spectrum of normal elements}

\begin{cthteorema}{(The norm of a normal element equals its spectral radius)}
Let $\A$ be a unital $C^*$-algebra and $a\in\A$ normal. Then
\[
 \|a\|=r(a).
\]
\end{cthteorema}
\origen{Theorem 2.6}

\begin{proof}[Proof]
Normality allows us to rearrange the factors:
\[
\|a^2\|^2
=\|(a^2)^*a^2\|
=\|(a^*)^2a^2\|
=\|(a^*a)^2\|
=\|a^*a\|^2
=\|a\|^4.
\]
Thus $\|a^2\|=\|a\|^2$. Since $a^2$ is also normal, induction gives
$\|a^{2^n}\|=\|a\|^{2^n}$. Applying the spectral radius formula along the
subsequence $2^n$,
\[
 r(a)=\lim_{n\to\infty}\|a^{2^n}\|^{1/2^n}=\|a\|.
\]
\end{proof}

\proofdiagram
{$\A$ is unital $C^*$; $a\in\A$ normal}
{$C^*$-identity; spectral radius formula}
{Prove $\|a^{2^n}\|=\|a\|^{2^n}$ by induction}
{$r(a)=\lim_n\|a^{2^n}\|^{1/2^n}=\|a\|$}

\begin{cthejemplo}{(Norm and radius of a normal matrix)}
Let us check the theorem for $A=\operatorname{diag}(2,-i)$ with the
Euclidean operator norm. Since $A$ is diagonal, it commutes with
$A^*$ and is normal. Its spectral values are $2$ and $-i$;
therefore $r(A)=\max\{2,1\}=2$. The norm of a diagonal matrix
in this norm is the largest absolute value of its entries, so
$\|A\|=2=r(A)$. This calculation makes concrete the equality that the
theorem guarantees for every normal element.
\end{cthejemplo}

\begin{cthnota}{(Where the hypotheses are used)}
Normality allows us to pass from $(a^*)^2a^2$ to $(a^*a)^2$.
The $C^*$-identity turns that algebraic equality into an equality of norms;
the spectral radius formula, proved in Part 1 for Banach algebras,
completes the argument.
\end{cthnota}


\begin{cthcorolario}{(The $C^*$-norm is determined algebraically)}
Let $\A$ be a unital $C^*$-algebra. For every $a\in\A$,
\[
 \|a\|=r(a^*a)^{1/2}.
\]
In particular, two norms that make the same unital involutive algebra a
$C^*$-algebra coincide.
\end{cthcorolario}
\origen{Corollary 2.7}

\begin{proof}[Proof]
The element $a^*a$ is self-adjoint and hence normal. Therefore
$r(a^*a)=\|a^*a\|=\|a\|^2$ by the preceding theorem and the $C^*$-identity.
Taking nonnegative square roots gives the stated formula.
If $\|\cdot\|_1$ and $\|\cdot\|_2$ make the same unital involutive algebra a
$C^*$-algebra, the invertibility of $a^*a-\lambda1_\A$ depends only on the
product and the identity. Thus $r(a^*a)$ is the same for both norms, and the
formula gives $\|a\|_1=\|a\|_2$ for every $a\in\A$.
\end{proof}

\proofdiagram
{$\A$ is unital $C^*$; $a\in\A$; $a^*a$ is normal}
{Previous theorem + $C^*$-identity}
{$r(a^*a)=\|a^*a\|=\|a\|^2$}
{$\|a\|=r(a^*a)^{1/2}$; uniqueness of the norm}

\begin{cthejemplo}{(The uniform norm is forced on $C(X)$)}
We want to apply the norm-uniqueness corollary to a function algebra.
Let $X$ be compact Hausdorff and $f\in C(X)$, with pointwise conjugation
as involution. The element $f^*f=|f|^2$ has spectrum $|f|^2(X)$,
by the criterion of pointwise invertibility from Part 1. The corollary gives
\[
 \|f\|=r(f^*f)^{1/2}
 =\bigl(\max_{x\in X}|f(x)|^2\bigr)^{1/2}
 =\|f\|_\infty.
\]
Any $C^*$-norm compatible with these operations must therefore
coincide with the uniform norm.
\end{cthejemplo}

\begin{cthteorema}{(Unital $*$-homomorphisms are contractive)}
Let $\A,\B$ be unital $C^*$-algebras and $\Phi:\A\to\B$ a unital $*$-homomorphism. Then
\[
 \|\Phi(a)\|\le\|a\|\quad(a\in\A),
 \qquad \|\Phi\|=1.
\]
\end{cthteorema}
\origen{Theorem 2.8}

\begin{proof}[Proof]
For each $b\in\A$, the unital homomorphism satisfies
$\spec(\Phi(b))\subset\spec(b)$ and hence
$r(\Phi(b))\le r(b)$. Taking $b=a^*a$,
\[
\|\Phi(a)\|^2
=\|\Phi(a)^*\Phi(a)\|
=\|\Phi(a^*a)\|
=r(\Phi(a^*a))
\le r(a^*a)=\|a\|^2.
\]
Since $\Phi(1_\A)=1_\B$ and both identities have norm one,
$\|\Phi\|\ge1$; contractivity gives equality.
\end{proof}

\proofdiagram
{$\Phi$ is unital and a $*$-homomorphism}
{Spectral inclusion for homomorphisms; normal-element theorem $\|x\|=r(x)$}
{Apply to $a^*a$: $\Phi(a^*a)=\Phi(a)^*\Phi(a)$}
{$\|\Phi(a)\|^2\le\|a\|^2$ and $\|\Phi\|=1$}

\begin{cthejemplo}{(Restriction as a contractive $*$-homomorphism)}
Let $X$ be compact Hausdorff and $Y\subset X$ closed and nonempty.
We want to compare a function defined on $X$ with its restriction to
$Y$. The map $\Phi:C(X)\to C(Y)$,
$\Phi(f)=f|_Y$, preserves sums, products, pointwise conjugation, and the
constant function $1$, so it is a unital $*$-homomorphism.
The contractivity theorem implies
$\|f|_Y\|_\infty\le\|f\|_\infty$, an inequality also obtained by
taking the supremum over a smaller set. Restricting the domain cannot
increase the uniform norm.
\end{cthejemplo}

\begin{cthobservacion}{(Real and imaginary parts)}
Let $\A$ be a $C^*$-algebra. Every $a\in\A$ has a unique expression
\[
 a=x+iy,
 \qquad
 x=\frac{a+a^*}{2},\quad y=\frac{a-a^*}{2i},
\]
with $x=x^*$ and $y=y^*$.
\end{cthobservacion}
\origen{Observation 2.9}

\begin{proof}[Proof]
The given formulas satisfy $x^*=x$, $y^*=y$, and $x+iy=a$. If also
$a=u+iv$ with $u=u^*$ and $v=v^*$, then $a^*=u-iv$; adding and subtracting
these two identities yields $u=(a+a^*)/2=x$ and
$v=(a-a^*)/(2i)=y$. This proves uniqueness.
\end{proof}

\begin{cthejemplo}{(Self-adjoint parts of a non-Hermitian matrix)}
To see the decomposition in the observation in a nontrivial case,
take $A=\begin{pmatrix}1&i\\0&2\end{pmatrix}$.
Calculating the real and imaginary parts gives
\[
 X=\frac{A+A^*}{2}
 =\begin{pmatrix}1&i/2\\-i/2&2\end{pmatrix},
 \qquad
 Y=\frac{A-A^*}{2i}
 =\begin{pmatrix}0&1/2\\1/2&0\end{pmatrix}.
\]
We check that $X^*=X$, $Y^*=Y$, and $A=X+iY$.
Thus even a matrix that is not self-adjoint can be reconstructed from
two matrices that are.
\end{cthejemplo}

\begin{cthlema}{(Vertical estimate for self-adjoint elements)}
Let $\A$ be a unital $C^*$-algebra, $x\in\A$ self-adjoint, and $t\in\mathbb R$. Then
\[
 \|x+it1_\A\|^2\le \|x\|^2+t^2.
\]
\end{cthlema}
\origen{Lemma 2.10}

\begin{proof}[Proof]
\[
\|x+it1_\A\|^2
=\|(x-it1_\A)(x+it1_\A)\|
=\|x^2+t^2 1_\A\|
\le\|x^2\|+t^2
=\|x\|^2+t^2.
\]
\end{proof}

\proofdiagram
{$x\in\A$ self-adjoint; $t\in\mathbb R$}
{$C^*$-identity}
{$\|x+it1_\A\|^2=\|x^2+t^2 1_\A\|$}
{$\le\|x\|^2+t^2$}

\begin{cthejemplo}{(The vertical bound is sharp for real scalars)}
Let us check whether the lemma's inequality can be improved in general.
In the $C^*$-algebra $\mathbb C$, take the self-adjoint element
$s\in\mathbb R$ and $t\in\mathbb R$. Then
$|s+it|^2=(s-it)(s+it)=s^2+t^2$ and $\|s\|^2=s^2$.
The lemma's inequality is an equality in this case; the term
$t^2$ and the coefficient of $\|s\|^2$ are therefore natural.
\end{cthejemplo}

\begin{cthproposicion}{(A vertical bound forces a scalar to be real)}
If $z\in\mathbb C$ and there exists $C\ge0$ such that
\[
 |z+it|^2\le C+t^2\qquad(t\in\mathbb R),
\]
then $z\in\mathbb R$.
\end{cthproposicion}
\origen{Proposition 2.11}

\begin{proof}[Proof]
If $z=x+iy$ with $x,y\in\mathbb R$, then
$|z+it|^2-t^2=x^2+y^2+2yt$. If $y>0$, the right-hand side tends to $+\infty$
as $t\to+\infty$; if $y<0$, this happens as $t\to-\infty$. The existence
of a uniform bound forces $y=0$.
\end{proof}

\proofdiagram
{$z\in\mathbb C$; $C\ge0$; write $z=x+iy$}
{Expand $|z+it|^2-t^2=x^2+y^2+2yt$}
{If $y\neq0$, let $t\to\pm\infty$ with the sign of $y$}
{The bound fails; hence $y=0$}

\begin{cthejemplo}{(A bound that fails for $1+i$)}
To see how Proposition~2.11 detects an imaginary part, test its hypothesis with
$z=1+i$. The calculation gives
\[
 |z+it|^2-t^2=|1+i(1+t)|^2-t^2=2+2t.
\]
As $t\to+\infty$, this expression exceeds every constant $C$. Thus there is
no bound $|1+i+it|^2\le C+t^2$ valid for every $t\in\mathbb R$.
The term linear in $t$ is the concrete sign that $z$ is not real.
\end{cthejemplo}

\begin{cthproposicion}{(Functionals attaining their norm at the identity are real on self-adjoint elements)}
Let $\A$ be a unital $C^*$-algebra and $\varphi:\A\to\mathbb C$ a continuous
linear functional with $\|\varphi\|=\varphi(1_\A)$. Then
\[
 \varphi(x)\in\mathbb R\qquad\text{for every }x\in\A\text{ with }x=x^*.
\]
\end{cthproposicion}
\origen{Proposition 2.12}

\begin{proof}[Proof]
If $\varphi=0$, then $\varphi(x)=0\in\mathbb R$ for every self-adjoint $x$.
Otherwise, $\varphi(1_\A)=\|\varphi\|>0$.
Define $\psi=\varphi/\varphi(1_\A)$; then $\|\psi\|=\psi(1_\A)=1$.
For $t\in\mathbb R$,
\[
 |\psi(x)+it|^2
 =|\psi(x+it1_\A)|^2
 \le\|x+it1_\A\|^2
 \le\|x\|^2+t^2.
\]
The previous proposition, with $z=\psi(x)$, implies that $\psi(x)$ is real;
since $\varphi(1_\A)>0$, so is $\varphi(x)$.
\end{proof}

\proofdiagram
{$\varphi:\A\to\mathbb C$ continuous; $\|\varphi\|=\varphi(1_\A)$; $x=x^*$}
{Vertical estimate lemma + preceding scalar proposition}
{If $\varphi\ne0$: $\psi=\varphi/\varphi(1_\A)$ and $|\psi(x)+it|^2\le\|x+it1_\A\|^2$}
{$\varphi=0$ or $\psi(x)\in\mathbb R$; hence $\varphi(x)\in\mathbb R$}

\begin{cthejemplo}{(An evaluation satisfies the reality hypothesis)}
Take $X=[0,1]$ and $x_0\in X$. The evaluation
$\varphi=\operatorname{ev}_{x_0}:C(X)\to\mathbb C$ satisfies
$\|\varphi\|=\varphi(1)=1$, as checked in Section~1.
For the self-adjoint function $h(t)=t^2$, we obtain
$\varphi(h)=x_0^2\in\mathbb R$.
Proposition~2.12 gives the same conclusion for \emph{every} self-adjoint
element of $C(X)$ and, more generally, for functionals with this norm
condition in any unital $C^*$-algebra.
\end{cthejemplo}

\begin{cthproposicion}{(The spectrum of a self-adjoint element is real)}
Let $\A$ be a unital $C^*$-algebra and $a\in\A$ self-adjoint. Then
\[
 \spec(a)\subset\mathbb R.
\]
\end{cthproposicion}
\origen{Proposition 2.13}

\begin{proof}[Proof]
By Problem 1.9 there is a commutative, involutive, unital Banach subalgebra
$\mathcal C\subset\A$ that contains $a$ and preserves its spectrum. In it,
\[
 \spec_\A(a)=\spec_{\mathcal C}(a)=\{\varphi(a):\varphi\in\Omega(\mathcal C)\}.
\]
Since $\mathcal C$ is closed and inherits the $C^*$-identity, it too is a
$C^*$-algebra. Every character of $\mathcal C$ has norm one and
$\varphi(1_{\mathcal C})=1$, so the previous proposition gives
$\varphi(a)\in\mathbb R$. The spectral description above then shows
that every $\lambda\in\spec_\A(a)$ is real.
\end{proof}

\proofdiagrambranch
{$\A$ is unital $C^*$; $a\in\A$, $a=a^*$}
{Reduce to a commutative subalgebra preserving the spectrum (Problem 1.9)}
{Spectrum via characters + functionals real on self-adjoint elements}
{$\lambda\in\spec(a)\Rightarrow\lambda=\varphi(a)\in\mathbb R$}
{$\spec(a)\subset\mathbb R$}

\begin{cthejemplo}{(Real eigenvalues of a Hermitian matrix)}
Consider $H=\begin{pmatrix}2&i\\-i&3\end{pmatrix}\in M_2(\mathbb C)$.
Conjugate transposition gives $H^*=H$, so Proposition~2.13 predicts
$\spec(H)\subset\mathbb R$. Indeed,
\[
 \det(\lambda I-H)=(\lambda-2)(\lambda-3)-1
 =\lambda^2-5\lambda+5,
 \qquad
 \spec(H)=\left\{\frac{5-\sqrt5}{2},\frac{5+\sqrt5}{2}\right\}.
\]
This finite-dimensional calculation confirms the spectral conclusion
that the proposition extends to self-adjoint elements with no associated matrix.
\end{cthejemplo}

\begin{cthnota}{(Where the hypotheses are used)}
The equality $a=a^*$ allows us to apply Problem 1.9 in its involutive case
and then the proposition on real-valued functionals. The identity is used
in the description of the spectrum by normalized characters.
\end{cthnota}


\begin{cthdefinicion}{(Positive element and order)}
Let $\A$ be a unital $C^*$-algebra. An element $x\in\A$ is \emph{positive},
written $x\ge0$, if $x=x^*$ and $\spec_\A(x)\subset[0,\infty)$.
For self-adjoint $x,y\in\A$, $x\le y$ means that $y-x\ge0$.
\end{cthdefinicion}

\begin{puente}
This is one of the clearest connections between the problems and Section 2:
the original proof that the spectrum is real explicitly uses the exercise in
Problem 1 that constructs a commutative subalgebra containing the element and
preserving its spectrum. The exercise is an essential part of the proof.
\end{puente}

\section{Characters, spectral permanence, and representation}

\begin{cthcorolario}{(Functionals attaining their norm at the identity preserve the involution)}
Let $\A$ be a unital $C^*$-algebra and $\varphi:\A\to\mathbb C$ a continuous
linear functional with $\|\varphi\|=\varphi(1_\A)$. Then, for every $a\in\A$,
\[
 \varphi(a^*)=\overline{\varphi(a)}.
\]
In particular, this holds for the characters of $\A$.
\end{cthcorolario}
\origen{Corollary 2.14}

\begin{proof}[Proof]
Let $a\in\A$ and write $a=x+iy$, where
$x=(a+a^*)/2$ and $y=(a-a^*)/(2i)$ are self-adjoint. By the previous
proposition, $\varphi(x),\varphi(y)\in\mathbb R$, and
\[
 \varphi(a^*)=\varphi(x-iy)=\varphi(x)-i\varphi(y)
 =\overline{\varphi(x)+i\varphi(y)}=\overline{\varphi(a)}.
\]
If $\varphi$ is a character, then it is a continuous functional with
$\varphi(1_\A)=\|\varphi\|=1$, so it satisfies the hypothesis used.
\end{proof}

\proofdiagram
{$\varphi:\A\to\mathbb C$ continuous; $\|\varphi\|=\varphi(1_\A)$; $a=x+iy$ with $x,y$ self-adjoint}
{Proposition 2.12: $\varphi(x),\varphi(y)\in\mathbb R$}
{$\varphi(a^*)=\varphi(x)-i\varphi(y)$}
{$\varphi(a^*)=\overline{\varphi(a)}$}

\begin{cthejemplo}{(An evaluation preserves conjugation)}[unbreakable]
In $C([0,1])$ choose $f(t)=1+it$ and $x\in[0,1]$.
The involution is pointwise conjugation, so
$f^*(t)=1-it$. For the character $\operatorname{ev}_x$,
\[
 \operatorname{ev}_x(f^*)=1-ix
 =\overline{1+ix}
 =\overline{\operatorname{ev}_x(f)}.
\]
This equality illustrates Corollary~2.14 for a concrete character; the
corollary also guarantees it for functionals attaining their norm at the identity.
\end{cthejemplo}

\begin{cthcorolario}{(The Gelfand transform is a $*$-homomorphism)}
If $\A$ is a unital commutative $C^*$-algebra, then
\[
 \Gamma:\A\to C(\Omega(\A)),\qquad a\mapsto\widehat a,
\]
is a unital $*$-homomorphism.
\end{cthcorolario}
\origen{Corollary 2.15}

\begin{proof}[Proof]
For $a,b\in\A$ and $\varphi\in\Omega(\A)$,
$\widehat{a+b}(\varphi)=\widehat a(\varphi)+\widehat b(\varphi)$,
$\widehat{ab}(\varphi)=\widehat a(\varphi)\widehat b(\varphi)$, and
$\widehat{1_\A}(\varphi)=1$, by linearity, multiplicativity, and the unital
property of the character. Also
$\widehat{\lambda a}(\varphi)=\lambda\widehat a(\varphi)$ for
$\lambda\in\mathbb C$. Thus the Gelfand transform is a linear,
multiplicative, unital homomorphism. Every $\varphi\in\Omega(\A)$ satisfies
$\|\varphi\|=\varphi(1_\A)=1$, so the previous corollary gives, for $a\in\A$,
\[
 \widehat{a^*}(\varphi)=\varphi(a^*)
 =\overline{\varphi(a)}=(\widehat a)^*(\varphi).
\]
Since the equality holds for every $\varphi\in\Omega(\A)$, the transform
preserves $*$.
\end{proof}

\proofdiagram
{Gelfand is a unital homomorphism}
{Characters preserve $*$}
{$\widehat{a^*}(\varphi)=\varphi(a^*)=\overline{\varphi(a)}$}
{$\widehat{a^*}=(\widehat a)^*$}

\begin{cthejemplo}{(Conjugation under the Gelfand transform)}
Take $\A=C([0,1])$ and $f(t)=1+it$. Since the characters are evaluations,
$\widehat f(\operatorname{ev}_x)=1+ix$ for each $x\in[0,1]$. Likewise,
\[
 \widehat{f^*}(\operatorname{ev}_x)=1-ix
 =\overline{\widehat f(\operatorname{ev}_x)}.
\]
Corollary~2.15 states this identity simultaneously for all elements and
characters: the Gelfand transform respects the involution.
\end{cthejemplo}

\begin{cthproposicion}{(Invertibility detected by $a^*a$ and $aa^*$)}
For $a$ in a unital $C^*$-algebra $\A$,
\[
 a\in\GL(\A)
 \quad\Longleftrightarrow\quad
 a^*a\in\GL(\A)\text{ and }aa^*\in\GL(\A).
\]
\end{cthproposicion}
\origen{Proposition 2.16}

\begin{proof}[Proof]
If $a\in\GL(\A)$, the identity
$(a^{-1})^*a^*=a^*(a^{-1})^*=1_\A$ shows that $a^*$ is invertible.
Hence its products $a^*a$ and $aa^*$ are invertible. Conversely, take
$y=(a^*a)^{-1}$ and $z=(aa^*)^{-1}$, so that
$y(a^*a)=1_\A$ and $(aa^*)z=1_\A$. Then $ya^*$ is a left inverse of $a$ and
$a^*z$ is a right inverse. Moreover,
\[
 ya^*=ya^*(aa^*z)=(ya^*a)a^*z=a^*z.
\]
Therefore a two-sided inverse exists.
\end{proof}

\proofdiagrambranch
{$a\in\A$; $a^*a\in\GL(\A)$}
{$aa^*\in\GL(\A)$}
{$y=(a^*a)^{-1}$ and $z=(aa^*)^{-1}$}
{$ya^*$ and $a^*z$ coincide by associativity}
{$a\in\GL(\A)$}

\begin{cthejemplo}{(A matrix's inverses detected by its products)}
In $M_2(\mathbb C)$, let
$A=\begin{pmatrix}1&1\\0&1\end{pmatrix}$. We have
\[
 A^*A=\begin{pmatrix}1&1\\1&2\end{pmatrix},
 \qquad
 AA^*=\begin{pmatrix}2&1\\1&1\end{pmatrix},
\]
and both determinants are $1$. By Proposition~2.16, $A$ is invertible;
directly, $A^{-1}=\begin{pmatrix}1&-1\\0&1\end{pmatrix}$.
The check with square matrices shows how the two products provide left
and right inverses in the abstract argument.
\end{cthejemplo}

\begin{cthproposicion}{(Spectral permanence in $C^*$-subalgebras)}
Let $\A$ be a unital $C^*$-algebra and $\B\subset\A$ a unital $C^*$-subalgebra
with $1_\B=1_\A$. For every $b\in\B$,
\[
 b\in\GL(\B)\Longleftrightarrow b\in\GL(\A),
 \qquad
 \spec_\B(b)=\spec_\A(b).
\]
\end{cthproposicion}
\origen{Proposition 2.17}

\begin{proof}[Proof]
If $b\in\GL(\B)$, its inverse $b^{-1}\in\B\subset\A$ satisfies
$bb^{-1}=b^{-1}b=1_\B=1_\A$; hence $b\in\GL(\A)$.
For the other direction, first let $h\in\B$ be self-adjoint and invertible
in $\A$. Since $\spec_\B(h)\subset\mathbb R$, for every $n\ge1$ the element
$h+i n^{-1}1_\B$ is invertible in $\B$. These elements converge to $h$;
inversion is continuous on $\GL(\A)$, so their inverses converge in $\A$
to $h^{-1}$. Since $\B$ is closed, $h^{-1}\in\B$.
For a general $b\in\B\cap\GL(\A)$, apply this argument to $b^*b$ and $bb^*$,
which are invertible in $\A$ when $b$ is, and obtain that both are invertible
in $\B$. The previous proposition then concludes that $b$ is invertible in
$\B$. Finally, for each $\lambda\in\mathbb C$ apply the invertibility
equivalence to $b-\lambda1_\B=b-\lambda1_\A$; it follows that
$\lambda\notin\spec_\B(b)$ if and only if
$\lambda\notin\spec_\A(b)$, that is,
$\spec_\B(b)=\spec_\A(b)$.
\end{proof}

\proofdiagrambranch
{$b\in\B\cap\GL(\A)$}
{$h=b^*b$ or $bb^*$: self-adjoint and invertible in $\A$}
{$\spec_\B(h)\subset\mathbb R$; $h+i n^{-1}1_\B\in\GL(\B)$}
{Continuity of inversion in $\A$ and closure of $\B$: $h^{-1}\in\B$}
{$b^*b,bb^*\in\GL(\B)$; hence $b\in\GL(\B)$}

\begin{cthejemplo}{(Counterexample: the disk algebra does not preserve spectra)}
Problem~1.12 compares the disk algebra $A(\mathbb D)$ with $C(\mathbb T)$.
Take the coordinate function $u(z)=z$ on $\mathbb T$. Since
$|u(z)|=1$ on the circle, $u$ is invertible in $C(\mathbb T)$,
with inverse $\overline u$. However, $\overline u$ does not belong to
$A(\mathbb D)$: a holomorphic inverse of $z$ on the disk would have a pole
at $0$. Hence $0\in\spec_{A(\mathbb D)}(u)$ whereas
$0\notin\spec_{C(\mathbb T)}(u)$.
This closed subalgebra is not stable under pointwise involution and does not
satisfy the $C^*$ hypothesis of Proposition~2.17.
\end{cthejemplo}

\begin{cthnota}{(Where the hypotheses are used)}
Stability of $\B$ under involution allows us to use $b^*b$ and $bb^*$ within it;
the common identity lets us compare invertibility, and closedness of $\B$
lets us retain the limit of the inverses.
\end{cthnota}


\begin{cthproposicion}{(The spectrum of a normal element is the character space it generates)}
Let $\A$ be a unital $C^*$-algebra and $a\in\A$ normal. Denote by
$C^*(a,1_\A)$ the smallest unital $C^*$-subalgebra of $\A$ containing $a$. Then
\[
 \Omega(C^*(a,1_\A))\longrightarrow\spec_\A(a),\qquad
 \varphi\longmapsto\varphi(a),
\]
is a homeomorphism.
\end{cthproposicion}
\origen{Proposition 2.18}

\begin{proof}[Proof]
The $C^*$-algebra $C^*(a,1_\A)$ is commutative because $a$ commutes with
$a^*$. By spectral permanence, the spectrum of $a$ in this subalgebra
is the same as in $\A$. In a unital commutative Banach algebra, every point
of the spectrum of $a$ is $\varphi(a)$ for some character $\varphi$; this gives
surjectivity. The map is continuous because the topology of
$\Omega(C^*(a,1_\A))$ is that of pointwise convergence on the algebra's
elements. If $\varphi,\psi\in\Omega(C^*(a,1_\A))$ and
$\varphi(a)=\psi(a)$, then also
$\varphi(a^*)=\overline{\varphi(a)}=\overline{\psi(a)}=\psi(a^*)$.
Since polynomials in $a$ and $a^*$ are dense in $C^*(a,1_\A)$,
linearity and multiplicativity give equality on those polynomials, and
continuity of the characters extends equality to the whole subalgebra.
Thus $\varphi=\psi$. The maximal ideal space and the spectrum are compact
Hausdorff spaces; a continuous bijection between them has a continuous
inverse and is a homeomorphism.
\end{proof}

\proofdiagram
{$a$ normal}
{$C^*(a,1_\A)$ is commutative; spectral permanence; characters}
{Surjective by the spectrum; injective because $a,a^*$ generate densely}
{Continuous bijection between compact Hausdorff spaces: homeomorphism}

\begin{cthejemplo}{(Interpolation on the spectrum of a diagonal matrix)}
Let $\lambda_1,\ldots,\lambda_n$ be distinct values and
$A=\operatorname{diag}(\lambda_1,\ldots,\lambda_n)$.
Proposition~2.18 identifies the characters of $C^*(A,I_n)$ with these
eigenvalues. Concretely, for arbitrary numbers $c_1,\ldots,c_n$,
Lagrange interpolation produces a polynomial $p$ such that
$p(\lambda_j)=c_j$; then
\[
 p(A)=\operatorname{diag}(c_1,\ldots,c_n).
\]
Thus $C^*(A,I_n)$ is the entire diagonal algebra and is identified with
$C(\{\lambda_1,\ldots,\lambda_n\})$. Each character selects one coordinate.
\end{cthejemplo}

\begin{cthteorema}{(Complex Stone--Weierstrass theorem)}
Let $X$ be compact Hausdorff and $\mathcal D\subset C(X)$ a unital
$*$-subalgebra that separates points. Then $\mathcal D$ is dense in $C(X)$
in the uniform norm.
\end{cthteorema}
\origen{Theorem 2.19. Used as an external tool; its full proof is not reproduced here.}

\begin{cthejemplo}{(Polynomials in the coordinate and its conjugate on a compact set)}
Let $K\subset\mathbb C$ be compact. We want to approximate functions in
$C(K)$ by algebraic expressions in the coordinate $z$. The unital
subalgebra of polynomials in $z$ and $\overline z$ separates points because
the function $z\mapsto z$ already distinguishes them; it is also stable
under conjugation. Stone--Weierstrass Theorem~2.19 therefore says that it
is dense in $C(K)$ in the uniform norm.
On $K=\mathbb T$ this includes trigonometric polynomials, since
$\overline z=z^{-1}$. Problem~2.8 uses precisely this density for the
functional calculus.
\end{cthejemplo}

\begin{cthteorema}{(Commutative Gelfand--Naimark theorem)}
Let $\A$ be a unital commutative $C^*$-algebra. The Gelfand transform
\[
 \Gamma:\A\longrightarrow C(\Omega(\A))
\]
is an isometric $*$-isomorphism.
\end{cthteorema}
\origen{Theorem 2.20}

\begin{proof}[Proof]
For each $a\in\A$, commutativity makes $a$ normal. The description of
the spectrum by characters, together with the norm theorem for normal
elements, gives
\[
 \|\widehat a\|_\infty=r(a)=\|a\|.
\]
Thus $\Gamma$ is isometric and injective. Its image is also closed: a
convergent sequence of images comes, by the isometry, from a Cauchy sequence
in $\A$, which converges by completeness. The image is a unital
$*$-subalgebra of $C(\Omega(\A))$ and separates points: if
$\varphi\neq\psi$, there is $a$ with $\varphi(a)\neq\psi(a)$.
Stone--Weierstrass says that the image is dense. Being closed and dense,
it is all of $C(\Omega(\A))$.
\end{proof}

\proofdiagrambranch
{$\A$ commutative unital $C^*$; $a\in\A$}
{Gelfand is a $*$-homomorphism; $\|a\|=r(a)$; Stone--Weierstrass}
{Isometry: $\|\widehat a\|_\infty=r(a)=\|a\|$}
{The image is closed, unital, $*$-stable, and separates points; hence it is dense}
{$\Gamma$ is surjective and isometric}

\begin{cthejemplo}{(Diagonal matrices as functions on two points)}
Let $\A=\{\operatorname{diag}(\alpha,\beta):\alpha,\beta\in\mathbb C\}$.
Its two characters are evaluations of the entries,
$\varphi_1(\operatorname{diag}(\alpha,\beta))=\alpha$ and
$\varphi_2(\operatorname{diag}(\alpha,\beta))=\beta$.
The Gelfand transform of Theorem~2.20 sends
$\operatorname{diag}(\alpha,\beta)$ to the function
$f:\{\varphi_1,\varphi_2\}\to\mathbb C$ with values $\alpha,\beta$.
The equality $\|\operatorname{diag}(\alpha,\beta)\|=\max\{|\alpha|,|\beta|\}
=\|f\|_\infty$ displays the theorem's isometry in this finite model.
\end{cthejemplo}

\begin{cthcorolario}{(Classification of unital commutative $C^*$-algebras)}
Let $\A$ be a unital commutative $C^*$-algebra. There exists a compact
Hausdorff space $X$ such that $\A$ is isometrically $*$-isomorphic to $C(X)$;
this space $X$ is unique up to homeomorphism.
\end{cthcorolario}
\origen{Corollary 2.21}

\begin{proof}[Proof]
Existence follows from Gelfand--Naimark with $X=\Omega(\A)$.
For uniqueness, let $Y$ be compact Hausdorff and
$\varphi\in\Omega(C(Y))$. If the functions in $\ker\varphi$ had no common
zero, compactness would give a finite subcover of the open sets
$\{y\in Y:f(y)\neq0\}$, with $f\in\ker\varphi$. There would then be
$f_1,\ldots,f_m\in\ker\varphi$ with no common zero. The corollary on the
involution implies that $f_j^*\in\ker\varphi$; hence
$g=\sum_{j=1}^m f_j^*f_j\in\ker\varphi$. But $g(y)>0$ for every $y\in Y$,
so $g$ is invertible in $C(Y)$ with inverse $1/g$. This contradicts
$1=\varphi(1)=\varphi(g)\varphi(g^{-1})=0$.
Thus there is $y_0\in Y$ at which every $f\in\ker\varphi$ vanishes.
For $h\in C(Y)$, the function $h-\varphi(h)1$ belongs to the kernel, so
$h(y_0)=\varphi(h)$. Thus $\varphi=\operatorname{ev}_{y_0}$.
The map $y\mapsto\operatorname{ev}_y$ is continuous because, for each
$h\in C(Y)$, its coordinate $y\mapsto\operatorname{ev}_y(h)=h(y)$ is continuous.
It is injective because $C(Y)$ separates points and surjective by the
preceding argument. Since $Y$ is compact and $\Omega(C(Y))$ is Hausdorff,
it is a homeomorphism.
If $F:\A\to C(Y)$ is a $*$-isomorphism, composition
$\chi\mapsto\chi\circ F$ is a homeomorphism from $\Omega(C(Y))$ to
$\Omega(\A)$, with inverse induced by $F^{-1}$. Therefore
$Y\cong\Omega(\A)=X$. Problem 2.5 develops this identification separately.
\end{proof}

\proofdiagram
{$\A$ commutative unital $C^*$; Gelfand--Naimark}
{$X=\Omega(\A)$}
{$\A\cong C(\Omega(\A))$; the characters of $C(X)$ are evaluations}
{The representation exists and $X$ is unique up to homeomorphism}

\begin{cthejemplo}{(Reconstructing the interval from its algebra)}
Consider $\A=C([0,1])$. The proof of Corollary~2.21 shows that every
character of $\A$ is an evaluation $\operatorname{ev}_x$ at a unique
$x\in[0,1]$. Thus the map $x\mapsto\operatorname{ev}_x$ identifies
$[0,1]$ topologically with $\Omega(\A)$.
If $\A$ were also $*$-isomorphic to $C(Y)$ for a compact Hausdorff space $Y$,
the same corollary would give $Y\cong\Omega(\A)\cong[0,1]$.
The function algebra therefore recovers the space from which it arises.
\end{cthejemplo}

\begin{investigacion}[unbreakable]
\textbf{Use in a construction.} In \cite[Sec.~2, pp.~465--466]{Karlovich2017},
$\SO(\mathbb R_+)$ is a unital commutative $C^*$-subalgebra of
$L^\infty(\mathbb R_+)$: its functions are continuous and bounded, and their
oscillation on $[r,2r]$ tends to zero as $r\to0$ or $r\to\infty$.
Identifying $C(\overline{\mathbb R}_+)$ with its restrictions to
$\mathbb R_+=(0,\infty)$, where $\overline{\mathbb R}_+=[0,\infty]$ is the
compactification, the fibers over $s\in\{0,\infty\}$ are
\[
 M_s(\SO(\mathbb R_+))=
 \{\xi\in\Omega(\SO(\mathbb R_+)):
 \xi|_{C(\overline{\mathbb R}_+)}=\operatorname{ev}_s\}.
\]
Proposition~2.2 of that source relates evaluation of a character in a fiber
to partial limits of slowly oscillating functions.
\textbf{What to recognize:} the characters and maximal ideal space studied
in the course serve here to encode information at the endpoints.
\end{investigacion}

\section{Continuous functional calculus}

\begin{cthteorema}{(Continuous functional calculus for a normal element)}
Let $\A$ be a unital $C^*$-algebra and $a\in\A$ normal. There is a unique
isometric $*$-homomorphism
\[
 C(\spec_\A(a))\longrightarrow\A,
 \qquad f\longmapsto f(a),
\]
that sends the identity function $z\mapsto z$ to $a$ and the constant
function $1$ to $1_\A$. Its image is $C^*(a,1_\A)$.
\end{cthteorema}
\origen{Theorem 2.22}

\begin{proof}[Proof]
Let $\B=C^*(a,1_\A)$. By Gelfand--Naimark,
$\Gamma_\B:\B\to C(\Omega(\B))$ is an isometric $*$-isomorphism.
The preceding proposition gives the homeomorphism
$\Psi:\Omega(\B)\to\spec(a)$, $\Psi(\varphi)=\varphi(a)$. The composition
\[
 C(\spec(a))\xrightarrow{f\mapsto f\circ\Psi} C(\Omega(\B))
 \xrightarrow{\Gamma_\B^{-1}}\B\subset\A
\]
is an isometric $*$-isomorphism from $C(\spec(a))$ onto $\B$.
For the identity function $\mathrm{id}(z)=z$, we have
$(\mathrm{id}\circ\Psi)(\varphi)=\varphi(a)=\widehat a(\varphi)$;
thus its image is $a$. The constant $1$ corresponds to
$\widehat{1_\A}$, so its image is $1_\A$.
If $F:C(\spec(a))\to\A$ is another isometric $*$-homomorphism with these
properties, then $F(\overline{\mathrm{id}})=a^*$ and it agrees with the
preceding construction on all polynomials in $z$ and $\overline z$.
Stone--Weierstrass makes that unital $*$-subalgebra dense in
$C(\spec(a))$; continuity of both maps extends equality to every
function. This gives the stated uniqueness and image.
\end{proof}

\proofdiagrambranch
{$\A$ is unital $C^*$; $a\in\A$ normal}
{$C^*(a,1_\A)$ commutative + Gelfand--Naimark}
{Identify $\Omega(C^*(a,1_\A))$ with $\spec_\A(a)$}
{Transport functions through the homeomorphism and the Gelfand transform}
{Define $f(a)$ uniquely}

\begin{cthejemplo}{(Spectral square root of a diagonal matrix)}
Take the positive normal matrix
$A=\operatorname{diag}(1,4)$ and the function
$f(t)=\sqrt t$ on $\spec(A)=\{1,4\}$. Theorem~2.22 gives
\[
 f(A)=\operatorname{diag}(f(1),f(4))
     =\operatorname{diag}(1,2),
 \qquad f(A)^2=A.
\]
In general, if
$A=U\operatorname{diag}(\lambda_1,\ldots,\lambda_n)U^*$ is normal,
$f(A)=U\operatorname{diag}(f(\lambda_1),\ldots,f(\lambda_n))U^*$.
The example shows that abstract functional calculus preserves the familiar
spectral recipe for matrices.
\end{cthejemplo}

\begin{cthnota}{(Dependence on the problems)}
Normality makes $C^*(a,1_\A)$ commutative; spectral permanence identifies
its spectrum with that of $a$ in $\A$. Problem 2.4 combines this calculus
with the spectrum of continuous functions from Problem 1.10 to obtain
$\spec(f(a))=f(\spec(a))$.
\end{cthnota}


\begin{cthdefinicion}{(Functional calculus by restriction from an ambient set)}
Let $\A$ be a unital $C^*$-algebra, $a\in\A$ normal, and $E\subset\mathbb C$
locally compact with $\spec_\A(a)\subset E$. For $f\in C_b(E)$ define
\[
 f(a):=(f|_{\spec_\A(a)})(a).
\]
This gives a unital $*$-homomorphism $C_b(E)\to\A$. The same formula
defines $f(a)$ if $f:E\to\mathbb C$ is continuous and unbounded, because its
restriction to the compact set $\spec_\A(a)$ is bounded. Only that
restriction matters.
\end{cthdefinicion}
\origen{Definition 2.23}

\begin{cthobservacion}{(Polynomials in $a$ and $a^*$)}
Let $a$ be normal in a unital $C^*$-algebra $\A$. If
$f(z)=z^2-\overline z^{\,3}+1$ for $z\in\mathbb C$, then
\[
 f(a)=a^2-(a^*)^3+1_\A.
\]
If $f(z)=|z|^2=\overline z z$, then $f(a)=a^*a$.
\end{cthobservacion}
\origen{Observation 2.24}

\begin{cthobservacion}{(In $C(X)$, functional calculus is composition)}
If $X$ is compact Hausdorff, $g\in C(X)$, and $f\in C(g(X))$, then
\[
 f(g)=f\circ g.
\]
\end{cthobservacion}
\origen{Observation 2.25}

\begin{proof}[Proof]
Under the identification $\Omega(C(X))\cong X$, the transform of $g$ takes
the value $g(x)$ at $x\in X$, and $\spec(g)=g(X)$. The construction of the
functional calculus transports $f$ through this transform; at each $x$ it
gives $f(g)(x)=f(g(x))$. Two functions in $C(X)$ that agree at every point
are equal, so $f(g)=f\circ g$.
\end{proof}

\begin{cthejemplo}{(The square root of $t^2$ is $|t|$)}
On $X=[-1,1]$ take the positive function $g(t)=t^2$.
Since $\spec_{C(X)}(g)=g(X)=[0,1]$, the continuous function
$f(s)=\sqrt s$ can be applied to $g$ by Theorem~2.22.
Observation~2.25 identifies the result pointwise:
\[
 \sqrt g(t)=f(g(t))=\sqrt{t^2}=|t|.
\]
The positive square root of the element $g$ is $|t|$, not the function $t$,
which changes sign.
\end{cthejemplo}

\begin{cthlema}{(Naturality of functional calculus)}
Let $\A,\B$ be unital $C^*$-algebras, $\Phi:\A\to\B$ a unital $*$-homomorphism,
$a\in\A$ normal, and $f\in C(\spec_\A(a))$. Since
$\spec_\B(\Phi(a))\subset\spec_\A(a)$, we evaluate $f$ at $\Phi(a)$ by
restricting it to $\spec_\B(\Phi(a))$. Then
\[
 \Phi(f(a))=f(\Phi(a)).
\]
\end{cthlema}
\origen{Lemma 2.26}

\begin{proof}[Proof]
The map $f\mapsto\Phi(f(a))$ is the composition of the functional calculus
for $a$ with $\Phi$; it is a unital, contractive $*$-homomorphism by the
preceding contractivity theorem.
The map $f\mapsto f(\Phi(a))$ is restriction
$C(\spec_\A(a))\to C(\spec_\B(\Phi(a)))$ followed by the functional calculus
for $\Phi(a)$, which is normal because $\Phi$ preserves products and
involution. It too is a unital, contractive $*$-homomorphism. Both maps send
the identity function to $\Phi(a)$ and the constant $1$ to $1_\B$; because
they preserve the involution, they also send $\overline z$ to $\Phi(a)^*$.
By linearity and multiplicativity they agree on polynomials in $z$ and
$\overline z$. These are dense by Stone--Weierstrass, and contractivity
allows us to pass to the limit for every $f\in C(\spec_\A(a))$.
\end{proof}

\proofdiagram
{$\Phi:\A\to\B$ unital and $*$; $a\in\A$ normal}
{Spectral inclusion; functional calculus; contractivity}
{The maps $f\mapsto\Phi(f(a))$ and $f\mapsto f(\Phi(a))$ agree on the identity and constants}
{By uniqueness/density, they agree on all of $C(\spec(a))$}

\begin{cthejemplo}{(A square root carried to multiplication operators)}
Let $\A=C([0,1])$ and represent it on $H=L^2([0,1])$ by
$\pi(f)=M_f$, the multiplication operator by $f$.
For $a(t)=t^2$, the positive square root in $\A$ is
$\sqrt a(t)=t$. By Lemma~2.26,
\[
 \pi(\sqrt a)=M_t=\sqrt{M_{t^2}}=\sqrt{\pi(a)}.
\]
Indeed, $M_t$ is positive and $M_t^2=M_{t^2}$.
We obtain the same operator by first taking the square root of the function
or by taking it after representing the algebra.
\end{cthejemplo}

\begin{cthteorema}{(A unital $*$-monomorphism is isometric)}
Let $\A,\B$ be unital $C^*$-algebras. If $\Phi:\A\to\B$ is a unital,
injective $*$-homomorphism, then
\[
 \|\Phi(a)\|=\|a\|\qquad(a\in\A).
\]
\end{cthteorema}
\origen{Theorem 2.27}

\begin{proof}[Proof]
We already know that $\|\Phi(a)\|\le\|a\|$. Suppose the inequality is strict
for some $a\neq0$ and put $x=a^*a$. Then $x$ is self-adjoint and
$\|\Phi(x)\|=\|\Phi(a)\|^2<\|a\|^2=\|x\|$.
Since $x$ is self-adjoint and $\|x\|=r(x)$, one of the points $\pm\|x\|$
belongs to $\spec(x)\subset[-\|x\|,\|x\|]$.
Moreover, $\spec(\Phi(x))\subset[-\|\Phi(x)\|,\|\Phi(x)\|]$.
Write $c=\|\Phi(x)\|$ and $d=\|x\|$, with $0\le c<d$, and define
\[
 f(t)=\max\!\left\{0,\min\!\left\{1,\frac{|t|-c}{d-c}\right\}\right\}
 \qquad(t\in\mathbb R).
\]
This function is continuous, vanishes on $[-c,c]$, and equals $1$ at both
points $\pm d$. Thus $f(\Phi(x))=0$, whereas $f(x)\neq0$ because the functional
calculus is isometric and $f$ does not vanish on $\spec(x)$. By naturality,
$0=f(\Phi(x))=\Phi(f(x))$, contradicting injectivity.
\end{proof}

\proofdiagrambranch
{$\Phi:\A\to\B$ unital, $*$, and injective}
{$a\in\A$; $x=a^*a$ is self-adjoint}
{If $\|\Phi(x)\|<\|x\|$, separate the two spectra by a continuous function}
{Naturality: $0=f(\Phi(x))=\Phi(f(x))$; injectivity contradicts $f(x)\neq0$}
{$\|\Phi(a)\|=\|a\|$}

\begin{cthejemplo}{(The diagonal representation of $C(\{1,2\})$)}
Let $X=\{1,2\}$ and $H=\ell^2(X)\cong\mathbb C^2$.
The unital representation $\pi:C(X)\to\mathcal B(H)$ assigns to
$f$ the matrix $M_f=\operatorname{diag}(f(1),f(2))$.
It is faithful, since $M_f=0$ implies $f(1)=f(2)=0$, and
\[
 \|\pi(f)\|=\max\{|f(1)|,|f(2)|\}=\|f\|_\infty.
\]
This concrete calculation agrees with the general conclusion of
Theorem~2.27: a faithful representation preserves the norm.
\end{cthejemplo}

\section{Recognizing self-adjoint elements, unitaries, projections, and positive elements by their spectra}

\begin{cthproposicion}{(Spectral classification of normal elements)}
Let $\A$ be a unital $C^*$-algebra and $a\in\A$ normal. Then
\begin{enumerate}[label=(\roman*)]
\item $a=a^*$ if and only if $\spec(a)\subset\mathbb R$;
\item $a$ is unitary if and only if $\spec(a)\subset\mathbb T$;
\item $a$ is a projection if and only if $\spec(a)\subset\{0,1\}$.
\end{enumerate}
\end{cthproposicion}
\origen{Proposition 2.28; corresponds directly to Problem 2.6}

\begin{proof}[Proof]
If $a=a^*$, we already know that $\spec(a)\subset\mathbb R$. Conversely,
if the spectrum is real, the functions $z$ and $\overline z$ agree on
$\spec(a)$, so functional calculus gives $a=a^*$.

If $a$ is unitary, the identity $a^*a=1_\A$ translates, by the isometric
functional calculus, into $|z|^2=1$ on $\spec(a)$; hence
$\spec(a)\subset\mathbb T$. Conversely, if
$\spec(a)\subset\mathbb T$, then $|z|^2=1$ on the spectrum and
$a^*a=(\overline z z)(a)=1_\A$; normality also gives $aa^*=1_\A$.

If $a$ is a projection, $a^2-a=0$. The isometry of functional calculus
applied to $z\mapsto z^2-z$ gives $z^2-z=0$ throughout $\spec(a)$;
therefore $\spec(a)\subset\{0,1\}$. Conversely, if the spectrum is
contained in $\{0,1\}$, then $z^2-z=0$ and $\overline z-z=0$ on the
spectrum; functional calculus gives $a^2=a$ and $a^*=a$.
\end{proof}

\proofdiagrambranch
{$\A$ is unital $C^*$; $a\in\A$ normal}
{Continuous functional calculus}
{Translate $a^*=a$, $a^*a=1$, $a^2=a$ into identities of functions on $\spec(a)$}
{$\overline z=z$, $|z|^2=1$, $z^2=z$ on the spectrum}
{The spectral inclusions characterize the three classes}

\begin{cthejemplo}{(Counterexamples: a small spectrum without normality)}
In $M_2(\mathbb C)$, take
$N=\begin{pmatrix}0&1\\0&0\end{pmatrix}$ and
$J=\begin{pmatrix}1&1\\0&1\end{pmatrix}$.
Both are triangular, so
$\spec(N)=\{0\}\subset\mathbb R\cap\{0,1\}$ and
$\spec(J)=\{1\}\subset\mathbb T$.
However, $N^*\neq N$ and $N^2=0\neq N$;
moreover, $J^*J=\begin{pmatrix}1&1\\1&2\end{pmatrix}\neq I$.
Problem~2.6(d) develops these cases. They show that the converse conclusions
of Proposition~2.28 require normality.
\end{cthejemplo}

\begin{cthproposicion}{(Four unitaries linearly generate every element)}
Let $\A$ be a unital $C^*$-algebra. Every $a\in\A$ is a linear combination
of at most four unitary elements.
\end{cthproposicion}
\origen{Proposition 2.29}

\begin{proof}[Proof]
Write $a=x_1+ix_2$ with $x_j$ self-adjoint. Let
$r_j=\max\{1,\|x_j\|\}$ and $h_j=x_j/r_j$, so $\|h_j\|\le1$.
Since $h_j=h_j^*$, its spectrum lies in $[-1,1]$. The continuous function
$f(t)=t+i\sqrt{1-t^2}$ satisfies
$|f(t)|^2=t^2+(1-t^2)=1$ and $f(t)+\overline{f(t)}=2t$ on that interval.
By functional calculus, $u_j=f(h_j)$ satisfies
$u_j^*u_j=u_ju_j^*=(|f|^2)(h_j)=1_\A$ and is unitary. Moreover,
$u_j+u_j^*=(f+\overline f)(h_j)=2h_j$. Therefore
\[
 a=\frac{r_1}{2}(u_1+u_1^*)+\frac{i r_2}{2}(u_2+u_2^*),
\]
the asserted combination.
\end{proof}

\proofdiagram
{$\A$ is unital $C^*$; $a=x_1+ix_2$ with $x_j=x_j^*$}
{Scale $h_j=x_j/r_j$ with $r_j=\max\{1,\|x_j\|\}$}
{$u_j=f(h_j)$, $f(t)=t+i\sqrt{1-t^2}$; $x_j=r_j(u_j+u_j^*)/2$}
{$a$ is a combination of $u_1,u_1^*,u_2,u_2^*$}

\begin{cthejemplo}{(A self-adjoint element as the average of two unitaries)}
In $M_2(\mathbb C)$ consider
$H=\operatorname{diag}(1/2,-1/2)$.
It is self-adjoint and $\|H\|\le1$. The construction in
Proposition~2.29 produces
\[
 U=H+i\sqrt{I-H^2}
 =\operatorname{diag}\!\left(
 \frac12+i\frac{\sqrt3}{2},
 -\frac12+i\frac{\sqrt3}{2}\right).
\]
Each diagonal entry has modulus one, so $U$ is unitary, and
$H=(U+U^*)/2$. This case exhibits the step that allows us to express
any element as a combination of at most four unitaries.
\end{cthejemplo}

\begin{cthproposicion}{(Positive square root)}
Let $\A$ be a unital $C^*$-algebra and $a\in\A$ normal with
$\spec_\A(a)\subset[0,\infty)$. Then there exists a positive self-adjoint
element $y\in C^*(a,1_\A)$ such that
\[
 y^2=a.
\]
The root obtained by functional calculus is denoted $y=\sqrt a$.
\end{cthproposicion}
\origen{Proposition 2.30}

\begin{proof}[Proof]
The function $t\mapsto\sqrt t$ is continuous and real on $\spec(a)$.
Define $y=\sqrt{\,\cdot\,}(a)$. Then $y^*=y$ and
$y^2=(t\mapsto t)(a)=a$.
The identification $C^*(a,1_\A)\cong C(\spec_\A(a))$ shows that the
transform of $y$ is $t\mapsto\sqrt t\ge0$; by the description of the
spectrum in a function algebra and spectral permanence,
$\spec_\A(y)\subset[0,\infty)$.
\end{proof}

\proofdiagram
{$\A$ is unital $C^*$; $a\in\A$ normal and $\spec_\A(a)\subset[0,\infty)$}
{Continuous functional calculus}
{Apply $f(t)=\sqrt t$}
{$y=f(a)$ is self-adjoint, $\spec(y)\subset[0,\infty)$, and $y^2=a$}

\begin{cthejemplo}{(Modulus of a nilpotent matrix)}
Let $A=\begin{pmatrix}0&2\\0&0\end{pmatrix}$ in $M_2(\mathbb C)$.
Although $A^2=0$, the product
\[
 A^*A=\operatorname{diag}(0,4)
\]
is positive. Proposition~2.30 allows us to take its positive square root,
which is calculated directly here:
\[
 |A|=(A^*A)^{1/2}=\operatorname{diag}(0,2).
\]
This $|A|$ retains the magnitude of the action of $A$ and is the object used
in Problem~2.7 for the polar decomposition. In general, applying the
proposition to $a^*a$ requires the standard fact $a^*a\ge0$.
\end{cthejemplo}

\begin{investigacion}
\textbf{Direct example.} In \cite[Sec.~1, p.~20]{BFK2006}, the operators
$(U_z\varphi)(t)=\varphi(zt)$, with $z,t\in\mathbb T$ and
$\varphi\in L^2(\mathbb T)$, form a unitary representation of the group
$\mathbb T$ on $L^2(\mathbb T)$; rotations preserve arc-length measure.
\textbf{What to recognize:} the abstract unitary here implements a concrete
geometric symmetry and participates in the nonlocal extension of the algebra
studied in that source.
\end{investigacion}

\section*{Historical note}
\origen{Historical note of the source chapter, pp.~8--10}

\paragraph{The original axioms.}
Gelfand and Naimark's 1943 work sought to represent $C^*$-algebras,
including noncommutative ones, by bounded operators on Hilbert spaces.
The source notes recall that the original definition included, besides
$\|a^*a\|=\|a\|^2$, isometry of the involution and invertibility of
$1+a^*a$. It was later proved that these additional conditions follow from
the $C^*$-identity.

\paragraph{$B^*$ and $C^*$.}
For a time, $B^*$ denoted the axiomatic structure and $C^*$ denoted
adjoint-closed subalgebras of bounded operators on Hilbert space. The
general Gelfand--Naimark representation theorem identifies the two
descriptions; the commutative theorem proved here is the course's step
toward that idea.

\paragraph{Rigidity and functional calculus.}
Theorem~2.6, Corollary~2.7, and Theorem~2.27 in the source notes show three
forms of norm rigidity. Spectral permanence contrasts with the disk algebra
of Problem~1.12, which is a closed subalgebra of $C(\mathbb T)$ but not a
$C^*$-subalgebra. Functional calculus transports continuous functions on
the spectrum to elements of the algebra without first choosing a
representation on Hilbert space; Stone--Weierstrass explains why functions
of $a$ and $a^*$ suffice.

\paragraph{Topology and positivity.}
The commutative classification recovers the compact space $X$ from the
characters of $C(X)$. This viewpoint motivated the later language of
noncommutative topology. The square root in Proposition~2.30 of the source
notes opens the systematic study of positive elements and order, while
Proposition~2.29 there anticipates the importance of unitaries for the
geometry of the algebra.

\chapter{Assignment 1: worked solutions and connections to the theory}

\begin{context}
After completing this assignment, you will be able to explain which continuity properties support
constructions by closure, decide in which algebra a spectrum is computed, and construct the commutative
subalgebra used to study a self-adjoint element. Each solution identifies the results it uses and how its
conclusion will be used later.
\end{context}

\section{Problem 1.1: continuity of multiplication and inversion}

Let $\A$ be a Banach algebra. In part (b), assume in addition that it has an identity.

\subsection*{1.1(a) Multiplication is jointly continuous}

\noindent\textbf{Statement.} If $x_n\to x$ and $y_n\to y$ in $\A$, prove that $x_ny_n\to xy$; equivalently, prove joint continuity of multiplication $\A\times\A\to\A$.

\begin{proof}[Solution]
Since $x_n\to x$, the sequence $(x_n)$ is bounded: there is $M$ with $\|x_n\|\le M$. Thus
\[
\|x_ny_n-xy\|
\le \|x_n\|\,\|y_n-y\|+\|x_n-x\|\,\|y\|
\le M\|y_n-y\|+\|y\|\|x_n-x\|\to0.
\]
Hence multiplication $\A\times\A\to\A$ is jointly continuous.
\end{proof}

\proofdiagram
{$x_n\to x$, $y_n\to y$}
{$\|ab\|\le\|a\|\|b\|$}
{Write $x_ny_n-xy=x_n(y_n-y)+(x_n-x)y$}
{$\|x_ny_n-xy\|\to0$}

\begin{cthejemplo}{(Product of convergent operator sequences)}
In $\mathcal B(H)$ we check Problem 1.1(a) with $T_n=T+n^{-1}I$ and $S_n=S+n^{-1}I$.
Then $T_nS_n-TS=n^{-1}(T+S)+n^{-2}I$, whence
$\|T_nS_n-TS\|\le n^{-1}(\|T\|+\|S\|)+n^{-2}\to0$.
This calculation shows that the product converges even when both factors vary simultaneously.
\end{cthejemplo}

\noindent\textbf{Dependencies.} The submultiplicative inequality in the definition of a Banach algebra and the boundedness of every convergent sequence. The estimate separates the variation of the two factors.

\noindent\textbf{How this problem is used later.} It allows passage to the closure of the algebras generated in Problem 1.3, passage to the limit in the normality relation in Problem 1.7, and control of polynomials in $x_n,x_n^*$ in Problem 2.8.

\subsection*{1.1(b) Inversion is continuous on $\GL(\A)$}

For this part, $\A$ is a unital Banach algebra and $\GL(\A)$ carries the norm inherited from
$\A$.

\noindent\textbf{Statement.} Prove that the map $\GL(\A)\to\GL(\A)$, $a\mapsto a^{-1}$, is continuous in the norm of $\A$.

\begin{proof}[Solution]
First, if $x\to1$ and $\|1-x\|<1$, the Neumann series gives
\[
 x^{-1}=\sum_{n=0}^\infty(1-x)^n,
 \qquad
 \|x^{-1}-1\|
 \le \sum_{n=1}^\infty\|1-x\|^n
 =\frac{\|1-x\|}{1-\|1-x\|}\to0.
\]
Therefore inversion is continuous at $1$. If $a_n\to a\in\GL(\A)$, then
$a^{-1}a_n\to1$ and, for all sufficiently large $n$, $a_n$ is invertible. Moreover,
\[
 a_n^{-1}=(a^{-1}a_n)^{-1}a^{-1}\longrightarrow a^{-1}.
\]
\end{proof}

\proofdiagram
{$a_n\to a\in\GL(\A)$}
{Continuity at $1$ by the Neumann series}
{$a_n=a(a^{-1}a_n)$ and $a^{-1}a_n\to1$}
{$a_n^{-1}=(a^{-1}a_n)^{-1}a^{-1}\to a^{-1}$}

\begin{cthejemplo}{(Inverse of a small perturbation)}
In $\mathcal B(H)$, we ask when a perturbation of the identity remains invertible.
If $E\in\mathcal B(H)$ satisfies $\|E\|<1$, the Neumann series gives
\[
 (I+E)^{-1}=\sum_{k=0}^{\infty}(-E)^k,
 \qquad \|(I+E)^{-1}-I\|\le\frac{\|E\|}{1-\|E\|}.
\]
The bound tends to zero with $\|E\|$: this is the concrete mechanism behind continuity in Problem 1.1(b).
\end{cthejemplo}

\noindent\textbf{Dependencies.} The Neumann series from Part 1 gives an inverse for $1-h$ if $\|h\|<1$; openness of $\GL(\A)$ allows the formula to be applied near any invertible element.

\noindent\textbf{How this problem is used later.} The proof of spectral permanence in Section 2 uses this continuity by going from $(b+i/n)^{-1}$ to $b^{-1}$ within a closed subalgebra.

\Needspace{14\baselineskip}
\section{Problem 1.2: commutation and inverses}

Let $\A$ be a unital complex algebra, let $a,b\in\A$, and assume that $a$ is invertible. In part (b),
assume in addition that $\A$ has an involution.

Recall that $x$ and $y$ \emph{$*$-commute} when
\[
 xy=yx,
 \qquad x^*y=yx^*.
\]

\subsection*{1.2(a) Commuting with $a$ is equivalent to commuting with $a^{-1}$}

\noindent\textbf{Statement.} Prove that $ab=ba$ if and only if $a^{-1}b=ba^{-1}$.

\begin{proof}[Solution]
If $ab=ba$, multiply first on the left and then on the right by $a^{-1}$:
\[
 b=a^{-1}ba \quad\Longrightarrow\quad ba^{-1}=a^{-1}b.
\]
Conversely, if $a^{-1}b=ba^{-1}$, multiply first on the left and then on the
right by $a$: from $b=aba^{-1}$ we obtain $ba=ab$. Thus $a$ commutes with $b$.
\end{proof}

\proofdiagram
{$ab=ba$, $a$ invertible}
{Multiply by $a^{-1}$ left and right}
{$a^{-1}b=ba^{-1}$}
{Equivalence, using $(a^{-1})^{-1}=a$}

\begin{cthejemplo}{(Commutation of diagonal matrices and their inverses)}
In $M_2(\mathbb C)$ we can see the equivalence of Problem 1.2(a) without resorting to limits.
Let's take $A=\operatorname{diag}(2,3)$ and $B=\operatorname{diag}(5,7)$.
As $A^{-1}=\operatorname{diag}(1/2,1/3)$,
\[
 AB=BA=\operatorname{diag}(10,21),\qquad
 A^{-1}B=BA^{-1}=\operatorname{diag}(5/2,7/3).
\]
Taking the inverse changes the entries of $A$ but preserves its commutation with $B$.
\end{cthejemplo}

\noindent\textbf{Dependencies.} Only associativity, the identity, and the existence of $a^{-1}$ are needed; no norm is involved.

\noindent\textbf{How this problem is used later.} In Problem 1.9 this identity is applied to $\lambda1-c$ to show that its inverse commutes with the chosen subalgebra.

\subsection*{1.2(b) $*$-commutation and inverses}

\noindent\textbf{Statement.} Prove that $a$ and $b$ $*$-commute if and only if $a^{-1}$ and $b$ $*$-commute.

\begin{proof}[Solution]
From $ab=ba$ and subpart (a), $a^{-1}b=ba^{-1}$. Additionally $a^*b=ba^*$; since $a^*$ is invertible,
$(a^*)^{-1}b=b(a^*)^{-1}$. Since $(a^*)^{-1}=(a^{-1})^*$, both conditions of
$*$-commutation hold. Conversely, if $a^{-1}$ and $b$ $*$-commute, apply the same argument
to the invertible element $a^{-1}$: its inverse $a$ and its adjoint $a^*$ commute with $b$.
\end{proof}

\Needspace{14\baselineskip}
\proofdiagram
{$a$ and $b$ $*$-commute}
{Subpart (a) applied to $a$ and $a^*$}
{$a^{-1}$ and $(a^{-1})^*=(a^*)^{-1}$ commute with $b$}
{$a^{-1}$ and $b$ $*$-commute}

\begin{cthejemplo}{(The inverse of a normal matrix)}
To check $*$-commutation in Problem 1.2(b), use the normal invertible matrix
$A=\operatorname{diag}(2,i)$ in $M_2(\mathbb C)$.
Then $A^{-1}=\operatorname{diag}(1/2,-i)$ and
\[
 A^{-1}(A^{-1})^*=(A^{-1})^*A^{-1}=\operatorname{diag}(1/4,1).
\]
Thus $A^{-1}$ remains normal: the key step is that its adjoint is the inverse of $A^*$.
\end{cthejemplo}

\noindent\textbf{Dependencies.} Problem 1.2(a) applies to both $a$ and $a^*$, along with $(a^{-1})^*=(a^*)^{-1}$. The definition of $*$-commutation sets the two relationships that need to be verified.

\noindent\textbf{How this problem is used later.} It controls inverses in the involutive version of Problem 1.9 and explains why the inverse of an invertible normal element is again normal.

\section{Problem 1.3: algebras generated by a set}

Let $\A$ be a Banach algebra and $S\subset\A$. We define
\[
\mathcal B(S)=\overline{\spn\{s_1\cdots s_k:k\ge1,\ s_j\in S\}}.
\]
If $\A$ has an involution, we put $S^*=\{s^*:s\in S\}$ and
$\mathcal B^*(S)=\mathcal B(S\cup S^*)$.

\subsection*{1.3(a) Minimality of $\mathcal B(S)$ and $\mathcal B^*(S)$}

\noindent\textbf{Statement.} Prove that $\mathcal B(S)$ is the smallest Banach subalgebra of $\A$ that contains $S$ and, when $\A$ is involutive, that $\mathcal B^*(S)$ is the minimal involutive Banach subalgebra containing $S$.

\begin{proof}[Solution]
The set of linear combinations of words in $S$ is a subalgebra; its closure is still a subalgebra by continuity of multiplication. Thus $\mathcal B(S)$ is a
Banach subalgebra containing $S$. If $C\subset\A$ is any Banach subalgebra with $S\subset C$,
then $C$ contains all the words $s_1\cdots s_k$, their linear combinations and, being
closed, its closure. Therefore $\mathcal B(S)\subset C$.

For $\mathcal B^*(S)$, apply the same argument to $S\cup S^*$. The generated algebra is stable under
$*$: the adjoint of a word is another word in $S\cup S^*$, and the involution is continuous. Hence it is the smallest involutive Banach subalgebra containing $S$.
\end{proof}

\pagebreak
\proofdiagram
{$S\subset\A$}
{Finite products, linear combinations and closure}
{Every closed subalgebra containing $S$ contains every word, linear combination, and limit}
{$\mathcal B(S)$ is the minimal Banach subalgebra containing $S$}

\begin{cthejemplo}{(Adding the unit changes the generated algebra)}
In $M_2(\mathbb C)$, the operator $T=\operatorname{diag}(0,2)$ distinguishes
$\mathcal B(\{T\})$ from the unital algebra generated by $T$.
As $T^k=\operatorname{diag}(0,2^k)$ for $k\ge1$,
\[
 \mathcal B(\{T\})=\{\operatorname{diag}(0,c):c\in\mathbb C\}.
\]
Adding $I$ gives all diagonal matrices. The distinction follows from the definition using words of positive length in Problem 1.3(a).
\end{cthejemplo}

\noindent\textbf{Dependencies.} The closure of linear combinations of words is complete by being closed in $\A$; Problem 1.1(a) allows products to pass to the closure. The continuity of $*$ gives the involutive case.

\noindent\textbf{How this problem is used later.} It constructs $C^*(a,1)$ concretely, the commutative subalgebra used in Section 2 for Gelfand--Naimark and functional calculus. The same generation operation appears in singular operator algebras.

\begin{investigacion}[title={Generation of a real operator algebra}]
To recognize Problem 1.3(a) in a research construction, consider the \emph{real} operator algebra in $\mathcal B(\mathbb R L^p(\Gamma,w))$ studied by Karlovich.
Its generators include $M_a$ with $a\in\SO^\diamond(\Gamma)$, $S_\Gamma$ and
real-linear conjugation $C$. For example, the word $M_aS_\Gamma C$ belongs to
the generated algebra, as do real linear combinations of words and their limits in
norm \cite[Eq.~(1.4), p.~3]{Karlovich2025}.
The generation operation is that of Problem 1.3(a) on real scalars. The presence of $C$
places this example outside the complex functional calculus of Section 2.
\end{investigacion}

\subsection*{1.3(b) commutation propagates to the generated algebra}

\noindent\textbf{Statement.} If $a\in\A$ commutes with each $s\in S$, prove commutation with every $x\in\mathcal B(S)$. In the involutive case, if $a$ $*$-commutes with each $s\in S$, prove the same conclusion for $\mathcal B^*(S)$ and for $*$-commutation.

\begin{proof}[Solution]
If $a$ commutes with each $s\in S$, then by induction it commutes with each word
$s_1\cdots s_k$. By linearity, it commutes with every linear combination of words. If
$x_n\to x$ and $ax_n=x_na$, continuity of multiplication gives $ax=x a$. Thus it commutes with all of
$\mathcal B(S)$.

In the involutive case, if $a$ $*$-commutes with each $s\in S$, then $a$ and $a^*$ commute with
$s$ and, taking adjoints, also with $s^*$. By the above argument, $a$ and $a^*$ commute with
all $\mathcal B^*(S)$, which is exactly to say that $a$ $*$-commutes with all its elements.
\end{proof}

\proofdiagram
{$as=sa$ for all $s\in S$}
{Commutation is preserved under products and sums}
{Pass to limits by continuity}
{$ax=xa$ for all $x\in\mathcal B(S)$}

\begin{cthejemplo}{(The commutant of diagonal matrices)}
In $M_2(\mathbb C)$ we want to follow the commutation from a generator to its algebra.
Let $D=\operatorname{diag}(1,2)$ and $A=\operatorname{diag}(3,4)$.
For any polynomial $p$,
\[
 Ap(D)=\operatorname{diag}(3p(1),4p(2))=p(D)A.
\]
The equality persists in polynomial limits in norm, as stated in Problem 1.3(b).
\end{cthejemplo}

\noindent\textbf{Dependencies.} The identity $as=sa$ is preserved when forming words and linear combinations; Problem 1.1(a) allows you to go to the limit. For $*$-commutation, the adjoint relations also apply.

\noindent\textbf{How this problem is used later.} Propagation of commutation justifies adding an inverse in Problem 1.9 without losing commutativity and helps analyze commutants of operator families.

\subsection*{1.3(c) If $S$ commutes pairwise, $\mathcal B(S)$ is commutative}

\noindent\textbf{Statement.} If $st=ts$ for all $s,t\in S$, prove that $\mathcal B(S)$ is commutative.

\begin{proof}[Solution]
Two words $s_1\cdots s_k$ and $t_1\cdots t_m$ can be rearranged factor by factor because all
generators commute. Therefore the words commute. By bilinearity the result extends to linear combinations of words, and by continuity to their closure.
\end{proof}

\proofdiagram
{The elements of $S$ commute pairwise}
{Reorder factors in two words}
{The words commute; then their linear combinations too}
{$\mathcal B(S)$ is commutative}

\begin{cthejemplo}{(Algebra of a nilpotent)}
The commutativity of an algebra generated by a single element does not require that it be normal.
In $M_2(\mathbb C)$, let's take $N=\begin{pmatrix}0&1\\0&0\end{pmatrix}$.
Since $N^2=0$, every word of positive length is $N$ or $0$, and
$\mathcal B(\{N\})=\mathbb C N$.
Two elements $cN,dN$ have product $0$ in either order. This illustrates Problem 1.3(c).
\end{cthejemplo}

\noindent\textbf{Dependencies.} The construction of Problem 1.3(a) is used: if the generators commute pairwise, any two words can be reordered. Problem 1.1(a) extends the identity to its limits.

\noindent\textbf{How this problem is used later.} It allows us to recognize the commutativity of $C^*(a,1)$ when $a$ is normal, a necessary step to represent that algebra as continuous functions in Section 2.

\subsection*{1.3(d) $*$-commutation pairwise}

\noindent\textbf{Statement.} In the involutive case, if each pair of elements of $S$ $*$-commutes, prove that $\mathcal B^*(S)$ is an involutive and commutative Banach subalgebra.

\begin{proof}[Solution]
Taking adjoints of $st=ts$ and $s^*t=ts^*$ also yields
$st^*=t^*s$ and $s^*t^*=t^*s^*$. Therefore all the elements of $S\cup S^*$ commute pairwise.
Applying part (c) to $S\cup S^*$ shows that $\mathcal B^*(S)$ is commutative; by construction,
it is an involutive Banach algebra.
\end{proof}

\proofdiagram
{Each pair $s,t\in S$ $*$-commutes}
{Take adjoints from relations}
{$S\cup S^*$ commutes pairwise}
{$\mathcal B^*(S)$ is commutative and involutive}

\begin{cthejemplo}{(A normal element generates a commutative involutive algebra)}
In $M_2(\mathbb C)$, the matrix $A=\operatorname{diag}(1,i)$ shows why
normality allows using functional calculus.
Its adjoint is $A^*=\operatorname{diag}(1,-i)$, so $AA^*=A^*A=I$.
Thus words in $A$, $A^*$, and $I$ are diagonal and commute, as does their closure.
By Problem 1.3(d), $C^*(A,I)$ is commutative, as required to apply Section 2.
\end{cthejemplo}

\noindent\textbf{Dependencies.} The definition of $*$-commutation applied to each pair of generators gives commutation in $S\cup S^*$; then Problem 1.3(c) is applied.

\noindent\textbf{How this problem is used later.} If $a$ is normal, applying this part to $S=\{a,1\}$ gives the commutative $C^*$-subalgebra $C^*(a,1)$, starting point of continuous functional calculus.

\section{Problem 1.4: the quotient by an ideal is a Banach algebra}

Let $\A$ be a Banach algebra and $J\subset\A$ be a closed two-sided ideal.

\noindent\textbf{Statement.} Prove that $(x+J)(y+J)=xy+J$ is well defined and that $\A/J$, with $\|x+J\|=\inf_{j\in J}\|x-j\|$, is a Banach algebra.

\begin{proof}[Solution]
If $x'=x+j_1$ and $y'=y+j_2$ with $j_1,j_2\in J$, then
\[
 x'y'-xy=xj_2+j_1y+j_1j_2\in J,
\]
so $(x+J)(y+J)=xy+J$ is well defined.

For the norm, take $j_1,j_2\in J$: 
\[
(x-j_1)(y-j_2)=xy-\underbrace{(xj_2+j_1y-j_1j_2)}_{\in J},
\]
so
\[
\|xy+J\|\le\|x-j_1\|\,\|y-j_2\|.
\]
Taking the infimum over $j_1$ and $j_2$ gives
$\|(x+J)(y+J)\|\le\|x+J\|\|y+J\|$. Finally, since $J$ is closed, the quotient space
$\A/J$ is complete. This is a Banach algebra.
\end{proof}

\proofdiagrambranch
{$J$ closed two-sided ideal}
{$\|x+J\|=\inf_{b\in J}\|x-b\|$}
{Check independence of representatives and submultiplicativity}
{Use that the quotient by a closed subspace is Banach}
{$\A/J$ is a Banach algebra}

\begin{cthejemplo}{(A class in the Calkin algebra)}
To see what the quotient in Problem 1.4 identifies, work in $X=\ell^2(\mathbb N)$.
Let $P$ be the rank one projection onto $\mathbb C e_1$; then $P\in\K(X)$ and
\[
 (I+P)+\K(X)=I+\K(X).
\]
Two different operators represent the same class because their difference is compact.
This is the basic equality that allows us to work in the Calkin algebra.
\end{cthejemplo}

\noindent\textbf{Dependencies.} A two-sided ideal ensures that the product of classes does not depend on representatives; closedness of $J$ ensures completeness of the quotient. The quotient norm and submultiplicativity in $\A$ are used.

\noindent\textbf{How this problem is used later.} By taking $\A=\mathcal B(X)$ and $J=\K(X)$ we obtain the Calkin algebra. The equivalence between Fredholm and invertibility in that quotient is Atkinson's theorem, an additional result that this exercise allows us to formulate but does not prove.

\begin{investigacion}[title={Fredholm operators in the Calkin algebra}]
To place Problem 1.4 in an application, let $X$ be an infinite-dimensional Banach space and
$\pi:\mathcal B(X)\to\mathcal B(X)/\K(X)$ the quotient map.
The class $\pi(A)$ is invertible exactly when there exists $B\in\mathcal B(X)$ with
$AB-I,BA-I\in\K(X)$; both equalities follow by multiplying classes in the quotient.
Atkinson's theorem adds that this condition is equivalent to $A$ being Fredholm: range
closed, finite-dimensional kernel, and finite-dimensional quotient $X/\operatorname{Ran}(A)$
\cite[p.~2]{Karlovich2025}. Problem 1.4 provides the algebra in which this equivalence is formulated; Atkinson's theorem is an additional result.
\end{investigacion}

\section{Problem 1.5: \texorpdfstring{$C_0(X)$}{C0(X)} is a Banach subalgebra of \texorpdfstring{$C_b(X)$}{Cb(X)}}

Let $X$ be a locally compact and Hausdorff space.

\noindent\textbf{Statement.} Define
\begin{multline*}
C_0(X)=\{f\in C_b(X):\text{ for every }\varepsilon>0\text{ there is a compact }K\subset X\\
\text{ such that }|f(x)|<\varepsilon\text{ outside }K\}.
\end{multline*}
Prove that it is a Banach subalgebra of $C_b(X)$.

\begin{proof}[Solution]
Any continuous function that vanishes at infinity is bounded: for $\varepsilon=1$ it is less than $1$ outside a compact set, and over that compact it reaches a finite maximum.
Also, if $f,h\in C_0(X)$, we choose compact $K_f,K_h$ outside which
$|f|,|h|<\varepsilon/2$; outside $K_f\cup K_h$ we have $|f+h|<\varepsilon$.
Multiplication by scalars is verified by choosing the threshold
$\varepsilon/(1+|c|)$ for $cf$. Thus $C_0(X)$ is a linear subspace of $C_b(X)$.

If $f\in C_0(X)$ and $g\in C_b(X)$, given $\varepsilon>0$ we choose a compact $K$ such that
$|f(x)|<\varepsilon/(1+\|g\|_\infty)$ outside $K$. So
$|f(x)g(x)|<\varepsilon$ outside $K$, so $fg\in C_0(X)$. This even shows that $C_0(X)$
is an ideal of $C_b(X)$.

To show closedness, let $f_n\in C_0(X)$ and $f_n\to f$ uniformly. Given $\varepsilon>0$,
we choose $n$ with $\|f-f_n\|_\infty<\varepsilon/2$ and a compact $K$ with
$|f_n(x)|<\varepsilon/2$ outside $K$. Then $|f(x)|<\varepsilon$ outside $K$.
Therefore $f\in C_0(X)$. A closed subspace of a Banach space is itself a Banach space.

Pointwise conjugation preserves $C_0(X)$ and
$\|f^*f\|_\infty=\sup_x|f(x)|^2=\|f\|_\infty^2$; therefore $C_0(X)$ is also a
$C^*$-algebra, although it has no unit when $X$ is not compact.
\end{proof}

\proofdiagram
{$f\in C_0(X)$, $g\in C_b(X)$}
{Definition of vanishing at infinity}
{$fg$ goes small out of a compact; uniform limits preserve the property}
{$C_0(X)$ is a closed ideal of $C_b(X)$, then Banach}

\begin{cthejemplo}{(A function that vanishes at infinity on the real line)}
In $X=\mathbb R$ we want to compare a vanishing function with the unit of $C_b(X)$.
For $f(x)=1/(1+x^2)$ and $\varepsilon>0$, if $|x|>\varepsilon^{-1/2}$ then
$|f(x)|<\varepsilon$; therefore $f\in C_0(\mathbb R)$.
The constant function $1$ does not satisfy that condition for $\varepsilon=1/2$ outside of any
compact. Problem 1.5 places $f$ in the closed ideal $C_0(\mathbb R)$ and leaves the unit outside of it.
\end{cthejemplo}

\noindent\textbf{Dependencies.} The definition of vanishing at infinity allows a compact to be chosen for each tolerance. The $C_b(X)$ uniform norm controls both products and limits.

\noindent\textbf{How this problem is used later.} If $X$ is not compact, $C_0(X)$ is the non-unital commutative model; for compact $X$, it coincides with $C(X)$. The general classification of $C_0(X)$ requires the non-unital version of Gelfand--Naimark, while Section 2 demonstrates the unital version.

\section{Problem 1.6: spectrum and involution}

Let $\A$ be a unital involutive Banach algebra and $x\in\A$.

\noindent\textbf{Statement.} For each $x\in\A$, prove $\spec(x^*)=\overline{\spec(x)}$.

\begin{proof}[Solution]
An element $a$ is invertible if and only if $a^*$ is, since from $ab=ba=1$ we obtain
$b^*a^*=a^*b^*=1$. Now
\[
(\lambda1-x^*)^*=\overline\lambda1-x.
\]
Therefore
\[
\lambda1-x^*\notin\GL(\A)
\Longleftrightarrow
\overline\lambda1-x\notin\GL(\A),
\]
which is equivalent to
\[
 \spec(x^*)=\overline{\spec(x)}.
\]
\end{proof}

\proofdiagram
{$\lambda1-x^*$}
{involution preserves invertibility: $a^{-1}\mapsto(a^*)^{-1}$}
{$(\lambda1-x^*)^*=\overline\lambda1-x$}
{$\lambda\in\spec(x^*)\Longleftrightarrow\overline\lambda\in\spec(x)$}

\begin{cthejemplo}{(Spectrum of a non-normal adjoint)}
The symmetry of Problem 1.6 does not require normality. Let's consider
$A=\begin{pmatrix}1&1\\0&i\end{pmatrix}$ in $M_2(\mathbb C)$.
Since $A$ is triangular, $\spec(A)=\{1,i\}$; its adjoint is lower triangular with diagonal
$(1,-i)$, so $\spec(A^*)=\{1,-i\}=\overline{\spec(A)}$.
Conjugation of the spectrum follows from the involution, even though $A$ does not commute with $A^*$.
\end{cthejemplo}

\noindent\textbf{Dependencies.} The spectrum definition of Part 1 and the identity $(a^*)^{-1}=(a^{-1})^*$ allow invertibility to be carried through the involution.

\noindent\textbf{How this problem is used later.} It shows a spectral symmetry that holds before imposing the identity $C^*$. To conclude that the entire spectrum of a self-adjoint is real, Section 2 adds arguments proper to $C^*$-algebras.

\section{Problem 1.7: normal elements form a closed set}

Let $\A$ be a $C^*$-algebra and $x_n\to x$ in norm, with each $x_n$ being normal.

\noindent\textbf{Statement.} Prove that the set $\{x\in\A:x^*x=xx^*\}$ of normal elements is closed in norm.

\begin{proof}[Solution]
As $\|x_n^*-x^*\|=\|x_n-x\|\to0$ and multiplication is jointly continuous,
\[
 x_n^*x_n\to x^*x,
 \qquad
 x_nx_n^*\to xx^*.
\]
But both terms are the same for each $n$. By uniqueness of limits, $x^*x=xx^*$.
\end{proof}

\proofdiagram
{$x_n\to x$, $x_n^*x_n=x_nx_n^*$}
{The involution is isometric; jointly continuous product}
{$x_n^*\to x^*$ and both products converge}
{$x^*x=xx^*$}

\begin{cthejemplo}{(Limit of normal matrices)}
In $M_2(\mathbb C)$ we can observe the closure of the set of normals without changing dimension.
Let $A_n=\operatorname{diag}(1+1/n,i-1/n)$; each $A_n$ is normal and
$\|A_n-\operatorname{diag}(1,i)\|=1/n$.
The limit $A=\operatorname{diag}(1,i)$ also commutes with its adjoint.
This is a concrete instance of the passage to the limit proved in Problem 1.7.
\end{cthejemplo}

\noindent\textbf{Dependencies.} The isometric involution gives $x_n^*\to x^*$ and Problem 1.1(a) allows going to the limit in the two products of the normality equation.

\noindent\textbf{How this problem is used later.} It is the first step of Problem 2.8: the limit $x$ is still normal, so $f(x)$ is defined by continuous functional calculus.

\section{Problem 1.8: spectral inclusion under homomorphisms}

Let $\A,\B$ be unital Banach algebras and $\psi:\A\to\B$ be a unital complex homomorphism.

\noindent\textbf{Statement.} For each $x\in\A$, prove $\spec_\B(\psi(x))\subset\spec_\A(x)$ and $r_\B(\psi(x))\le r_\A(x)$.

\begin{proof}[Solution]
If $\lambda\notin\spec_\A(x)$, then $(\lambda1-x)^{-1}$ exists. Applying $\psi$,
\[
 \psi(\lambda1-x)\,\psi((\lambda1-x)^{-1})
 =\psi((\lambda1-x)^{-1})\,\psi(\lambda1-x)=1_\B,
\]
so $\lambda1-\psi(x)$ is invertible. This proves
$\spec_\B(\psi(x))\subset\spec_\A(x)$. Taking maxima of $|\lambda|$ we obtain
$r(\psi(x))\le r(x)$.
\end{proof}

\proofdiagram
{$\psi:\A\to\B$ unital homomorphism}
{Unital homomorphisms send inverses to inverses}
{$\lambda1-\psi(x)=\psi(\lambda1-x)$}
{$\spec(\psi(x))\subset\spec(x)$ and $r(\psi(x))\le r(x)$}

\begin{cthejemplo}{(Strict spectral inclusion in a quotient)}
To show that the inclusion of Problem 1.8 can be strict, let's take
$\A=\mathbb C\oplus\mathbb C$ with supremum norm and $J=\{0\}\oplus\mathbb C$.
The class of $a=(1,2)$ in $\A/J\cong\mathbb C$ is $1$.
As $\lambda1-a=(\lambda-1,\lambda-2)$ is invertible exactly when
$\lambda\notin\{1,2\}$,
\[
 \spec_{\A/J}(a+J)=\{1\}\subsetneq\{1,2\}=\spec_\A(a).
\]
The quotient may lose spectral points, according to Problem 1.8.
\end{cthejemplo}

\noindent\textbf{Dependencies.} A unital homomorphism preserves the inverse of $\lambda1-x$ when it exists; the definition of spectrum converts this algebraic observation into an inclusion.

\noindent\textbf{How this problem is used later.} In Section 2 we apply inclusion to $a^*a$ to prove that unital $*$-homomorphisms are contractive.

\section{Problem 1.9: a spectrum-preserving commutative subalgebra}

Let $\A$ be a unital Banach algebra and $S\subset\A$ be a family that commutes pairwise. For the
involutive version, we further assume that $\A$ is involutive and the elements of $S$
$*$-commute pairwise.

\noindent\textbf{Statement.} Prove that there exists a unital and commutative Banach subalgebra $C\subset\A$ that contains $S$ and satisfies $\spec_C(s)=\spec_\A(s)$ for every $s\in S$. In the involutive case with $*$-commutation pairwise, require in addition that $C$ be involutive.

\begin{proof}[Solution]
Let $\mathscr F$ be the family of unital commutative Banach subalgebras $C\subset\A$ containing $S$, ordered by inclusion. It is non-empty: it contains the closure of the unital algebra generated by
$S$. If $\{C_i\}$ is a chain, the closure of $\bigcup_i C_i$ is again a unital commutative Banach subalgebra containing $S$. By Zorn's lemma there exists a maximal element $C$.

We assert that $C$ is inverse-closed in $\A$. If $c\in C$ is invertible in $\A$, then
from $cd=dc$ for $d\in C$ it follows that $c^{-1}d=dc^{-1}$. Therefore the closed unital subalgebra
generated by $C\cup\{c^{-1}\}$ is still commutative. Maximality forces
$c^{-1}\in C$. Consequently, for $s\in S$,
\[
\lambda1-s\in\GL(\A)\Longleftrightarrow\lambda1-s\in\GL(C),
\]
and the spectra match.

In the involutive case, we start from the family of involutive, commutative, unital Banach subalgebras containing $S$. The assumption of $*$-commutation guarantees that the family is non-empty by
Problem 1.3(d), applied to $S\cup\{1\}$. If $c\in C$ is invertible in $\A$, both $c^{-1}$ and $(c^{-1})^*$ commute with
$C$ and each other: $c,c^*\in C$ commute and their inverses as well. They can be added while keeping
commutativity and involution. Maximality again forces $c^{-1}\in C$.
\end{proof}

\proofdiagrambranch
{$S$ commutes pairwise}
{Family of closed, commutative, unital subalgebras containing $S$}
{Zorn produces a maximal; the inverses of its elements also commute with it}
{Maximality forces the inverses to belong to the subalgebra}
{$\spec_\A(s)=\spec_C(s)$ for $s\in S$}

\begin{cthejemplo}{(A spectrally faithful subalgebra of matrices)}
In $M_2(\mathbb C)$, the self-adjoint $A=\operatorname{diag}(1,2)$ allows building
explicitly the subalgebra of Problem 1.9.
Let's take $C$ as all diagonal matrices. The inverse of
$\lambda I-A=\operatorname{diag}(\lambda-1,\lambda-2)$, when it exists, still belongs to $C$:
\[
 (\lambda I-A)^{-1}=\operatorname{diag}((\lambda-1)^{-1},(\lambda-2)^{-1}).
\]
Therefore $\spec_C(A)=\{1,2\}=\spec_{M_2(\mathbb C)}(A)$; the spectrum is conserved because
the necessary inverses belong to $C$.
\end{cthejemplo}

\noindent\textbf{Dependencies.} Problem 1.3 is used to generate a commutative subalgebra, Problem 1.2 to commute with inverses and Zorn's lemma to choose a maximal subalgebra. The involutive version requires $*$-commutation.

\noindent\textbf{How this problem is used later.} Section 2 uses this reduction to study a self-adjoint using characters without changing its spectrum. In the $C^*$ environment, spectral permanence then offers a more direct path to $C^*(a,1)$.

\pagebreak
\section{Problem 1.10: spectrum of a continuous function}

Let $X$ be compact Hausdorff and $f\in C(X)$. Then we will compare with $C_b(X)$ for a locally compact non-compact space.

\noindent\textbf{Statement.} Prove that $\spec_{C(X)}(f)=f(X)$ for each $f\in C(X)$ and decide, using an example, whether the equality still holds in $C_b(X)$ for $X$ only locally compact.

\begin{proof}[Solution]
If $\lambda\notin f(X)$, the function $\lambda-f$ has no zeros. Since $X$ is compact,
$|\lambda-f|$ has positive minimum; therefore
$(\lambda-f)^{-1}\in C(X)$ and $\lambda\notin\spec(f)$. If $\lambda=f(x_0)$, then
$\lambda-f$ vanishes at $x_0$ and cannot be invertible. Therefore $\spec(f)=f(X)$.

If $X$ is only locally compact and we use $C_b(X)$, the statement may fail because a nonvanishing function can approach zero without reaching zero. For example, in $X=(0,1)$, $f(x)=x$ has
$f(X)=(0,1)$, but
\[
 \spec_{C_b((0,1))}(f)=[0,1].
\]
Indeed, for $\lambda\in(0,1)$ the function $\lambda-f$ vanishes; for $\lambda=0,1$,
their reciprocals are unbounded. If $\lambda\notin[0,1]$, however,
$|\lambda-f(x)|\ge\operatorname{dist}(\lambda,[0,1])>0$ for all $x$, so its reciprocal belongs to $C_b((0,1))$. This proves the announced equality.
\end{proof}

\proofdiagram
{$f\in C(X)$, $X$ compact}
{Pointwise invertibility in an algebra of functions}
{$\lambda1-f$ invertible $\Longleftrightarrow$ $\lambda-f$ does not vanish}
{$\spec(f)=f(X)$}

\begin{cthejemplo}{(A spectrum shaped like a circle)}
To read Problem 1.10 geometrically, consider $f(z)=2+z$ in $C(\mathbb T)$.
Function $\lambda-f$ vanishes on $\mathbb T$ exactly when
$\lambda=2+z$ for some $|z|=1$; outside those values it has a continuous reciprocal.
Therefore
\[
 \spec_{C(\mathbb T)}(f)=\{\lambda\in\mathbb C:|\lambda-2|=1\}.
\]
The spectrum is the image of $f$, a circle with center $2$ and radius $1$.
\end{cthejemplo}

\noindent\textbf{Dependencies.} In $C(X)$, $\lambda1-f$ is invertible exactly when $\lambda-f$ does not vanish; compactness of $X$ ensures that its reciprocal is continuous and bounded.

\noindent\textbf{How this problem is used later.} Identifies the spectrum with a range of values in the Gelfand--Naimark representation and is the specific step that Problem 2.4 uses to demonstrate spectral mapping.

\pagebreak
\section{Problem 1.11: \texorpdfstring{$\mathcal B(H)$}{B(H)} and the unilateral shift}

Let $H$ be a Hilbert space and $T\in\mathcal B(H)$. Subsections (b)--(d) change the assumptions
on $H$ as stated in each part.

\subsection*{1.11(a) $\mathcal B(H)$ is a $C^*$-algebra}

\noindent\textbf{Statement.} Prove that $\mathcal B(H)$, with the operator norm and the involution $T\mapsto T^*$, is a $C^*$-algebra.

\begin{proof}[Solution]
Part 1 already established that bounded operators on a Banach space form a
Banach algebra; therefore $\mathcal B(H)$ is complete. The adjoint rules give an involution.
The norm of the adjoint coincides with that of $T$: by the characterization of the norm of a vector of
a Hilbert space using inner products,
\[
 \|T^*\|=\sup_{\|\eta\|=\|\xi\|=1}|\langle T^*\eta,\xi\rangle|
 =\sup_{\|\eta\|=\|\xi\|=1}|\langle\eta,T\xi\rangle|=\|T\|.
\]
Submultiplicativity then gives $\|T^*T\|\le\|T^*\|\|T\|=\|T\|^2$.
For the reverse inequality, if $\|\xi\|=1$,
\[
\|T\xi\|^2=\langle T^*T\xi,\xi\rangle\le\|T^*T\|.
\]
Taking the supremum over $\xi$ gives $\|T\|^2\le\|T^*T\|$.
\end{proof}

\proofdiagram
{$T\in\mathcal B(H)$}
{adjoint properties}
{$\|T^*T\|\le\|T\|^2$ and $\|T\xi\|^2=\langle T^*T\xi,\xi\rangle$}
{$\|T^*T\|=\|T\|^2$}

\begin{cthejemplo}{(The identity $C^*$ in a diagonal matrix)}
In $H=\mathbb C^2$ we want to check the equality of norms that distinguishes
$\mathcal B(H)$ as $C^*$-algebra.
For $T=\operatorname{diag}(2,i)$, there is
$T^*T=\operatorname{diag}(4,1)$, then $\|T^*T\|=4$ and $\|T\|=2$.
Thus $\|T^*T\|=\|T\|^2$, the identity verified generally in Problem 1.11(a).
\end{cthejemplo}

\noindent\textbf{Dependencies.} Part 1 presents $\mathcal B(H)$ as a Banach algebra; the properties of the adjoint and the operator norm allow the $C^*$-identity to be verified.

\noindent\textbf{How this problem is used later.} Provides the concrete model from Section 2 and the operator space in which Problem 2.1 isometrically represents $C_b(X)$.

\subsection*{1.11(b) In finite dimension, spectrum = eigenvalues}

\noindent\textbf{Statement.} If $\dim H<\infty$, prove that $\spec(T)$ matches the set of eigenvalues of $T$.

\begin{proof}[Solution]
In finite dimension a linear operator is invertible if and only if its kernel is zero. So
$\lambda I-T$ is noninvertible exactly when $\xi\neq0$ exists with
$(\lambda I-T)\xi=0$, that is, $T\xi=\lambda\xi$.
\end{proof}

\proofdiagram
{$\dim H<\infty$}
{In finite dimension: injective $\Longleftrightarrow$ surjective}
{$\lambda I-T$ non-invertible $\Longleftrightarrow\Ker(\lambda I-T)\neq0$}
{$\lambda\in\spec(T)\Longleftrightarrow\lambda$ eigenvalue}

\begin{cthejemplo}{(Eigenvalues of a triangular matrix)}
In dimension two, the determinant allows the spectrum and
eigenvectors of Problem 1.11(b).
For $T=\begin{pmatrix}1&1\\0&2\end{pmatrix}$,
$\det(\lambda I-T)=(\lambda-1)(\lambda-2)$.
Also $Te_1=e_1$ and $T(1,1)=(2,2)$, so $1$ and $2$ are eigenvalues.
There are no other spectral points: in finite dimension the singularity of $\lambda I-T$
matches the existence of an eigenvector.
\end{cthejemplo}

\noindent\textbf{Dependencies.} In finite dimension, injectivity and surjectivity of $\lambda I-T$ are equivalent by rank and nullity. The definition of spectrum from Part 1 applies.

\noindent\textbf{How this problem is used later.} Allows you to interpret the matrix examples in Section 2 as spectrum calculations using eigenvalues; Problem 1.11(c)--(d) show what changes in infinite dimension.

\subsection*{1.11(c) Approximate eigenvectors produce spectral points}

\noindent\textbf{Statement.} If $\lambda\in\mathbb C$ and there exist unit vectors $\xi_n\in H$ such that $\|(\lambda I-T)\xi_n\|\to0$, prove that $\lambda\in\spec(T)$.

\begin{proof}[Solution]
If $\lambda\notin\spec(T)$, then $R=(\lambda I-T)^{-1}$ is bounded. Therefore
\[
1=\|\xi_n\|=\|R(\lambda I-T)\xi_n\|
\le\|R\|\,\|(\lambda I-T)\xi_n\|\to0,
\]
contradiction. Then $\lambda\in\spec(T)$.
\end{proof}

\proofdiagram
{$\|\xi_n\|=1$, $\|(\lambda I-T)\xi_n\|\to0$}
{Assume $\lambda I-T$ invertible}
{$\xi_n=(\lambda I-T)^{-1}(\lambda I-T)\xi_n$}
{$1\le\|(\lambda I-T)^{-1}\|\|(\lambda I-T)\xi_n\|\to0$, contradiction}

\begin{cthejemplo}{(Approximate spectrum without eigenvector)}
In $\ell^2(\mathbb N)$ we want to see that an approximate spectral point can lack
eigenvector. Let's define $T(\xi_n)=(\xi_n/n)$ and take $\lambda=0$.
Each $e_n$ is a unit vector and $\|Te_n\|=1/n\to0$, but $T\xi=0$ implies
$\xi_n/n=0$ for all $n$ i.e. $\xi=0$.
By Problem 1.11(c), $0\in\spec(T)$; the calculation distinguishes approximate spectrum from point spectrum.
\end{cthejemplo}

\noindent\textbf{Dependencies.} If $\lambda I-T$ is invertible, its inverse is bounded. The norm estimate of the composition contradicts the existence of almost null unit vectors.

\noindent\textbf{How this problem is used later.} Defines the approximate point spectrum and prepares the comparison with the unilateral shift of Problem 1.11(d), whose spectral zero is not an approximate eigenvalue.

\subsection*{1.11(d) Unilateral shift shows the difference}

\noindent\textbf{Statement.} Let $S(x_1,x_2,\ldots)=(0,x_1,x_2,\ldots)$ be
$\ell^2(\mathbb N)$. Prove that $0\in\spec(S)$ and that there are no unit vectors
$\xi_n$ with $\|S\xi_n\|\to0$.

\begin{proof}[Solution]
$S$ is not surjective: no vector with non-zero first coordinate is in its range. So $S$ is not invertible and $0\in\spec(S)$. But
\[
 \|S\xi\|^2=0+\sum_{n\ge1}|\xi_n|^2=\|\xi\|^2.
\]
Therefore, for every unit vector $\xi$, $\|S\xi\|=1$; a sequence of unit vectors cannot exist with
$\|S\xi_n\|\to0$.
\end{proof}

\proofdiagram
{$S(x_1,x_2,\dots)=(0,x_1,x_2,\dots)$}
{$S$ is isometry but is not surjective}
{$0\in\spec(S)$; however $\|S\xi\|=\|\xi\|$}
{No unit vectors with $\|S\xi_n\|\to0$}

\begin{cthejemplo}{(An isometry that is not surjective)}
To distinguish conservation of norm and invertibility in Problem 1.11(d), let's take the first
$\ell^2(\mathbb N)$ base vector.
The shift satisfies $Se_1=e_2$ and $\|Se_1\|=\|e_1\|=1$, but no vector can be mapped to $e_1$, because the first coordinate of $S\xi$ is always zero.
Isometry is not surjective; that is why $0$ belongs to its spectrum without having unit vectors that are nearly annihilated.
\end{cthejemplo}

\noindent\textbf{Dependencies.} The unilateral shift is an isometry, but its range omits vectors with a non-zero first coordinate. The definition of spectrum detects this lack of surjectivity.

\noindent\textbf{How this problem is used later.} Problem 2.10 will determine the entire spectral disk using eigenvectors of adjoint. Isometry without surjectivity explains why $0$ can be spectral without being an approximate eigenvalue.

\begin{investigacion}[title={Two shifts in different spaces}]
To recognize the scope of Problem 1.11(d), let's compare the shift $S$ of $\ell^2(\mathbb N)$ with
$V_\alpha$ in Haseman's formulation. The first one satisfies $Se_1=e_2$; the second
composes boundary functions, $(V_\alpha f)(t)=f(\alpha(t))$
\cite[pp.~463--464, ecs.~(1.1)--(1.2)]{Karlovich2017}.
Both move data, but act in different spaces and geometries. Problem 1.11(d)
determines $0\in\spec(S)$; it does not determine by analogy any point of $\spec(V_\alpha)$.
\end{investigacion}

\section{Problem 1.12: the disk algebra and failure of spectral permanence in general}

Let $\mathbb D=\{z:|z|<1\}$, $\mathbb T=\{z:|z|=1\}$ and
\[
 A(\mathbb D)=\{f\in C(\overline{\mathbb D}): f\text{ is holomorphic on }\mathbb D\}.
\]

\subsection*{1.12(a) $A(\mathbb D)$ is a Banach algebra}

\noindent\textbf{Statement.} Prove that $A(\mathbb D)$ is a unital Banach algebra with the norm of the supremum in $\overline{\mathbb D}$.

\begin{proof}[Solution]
Let $(f_n)$ be Cauchy in the supremum norm. It converges uniformly on $\overline{\mathbb D}$ to a
continuous function $f$. In each $K\subset\mathbb D$ compact, the convergence is uniform; by
Weierstrass' theorem, $f$ is holomorphic in $\mathbb D$. So $f\in A(\mathbb D)$ and the space is
complete. The sums and products of holomorphic functions are holomorphic, and the constant function
$1$ belongs to $A(\mathbb D)$; operations and norm are inherited from
$C(\overline{\mathbb D})$.
\end{proof}

\proofdiagram
{Cauchy sequence in supremum norm}
{$C(\overline{\mathbb D})$ is Banach; Weierstrass theorem for holomorphic limits}
{The uniform limit is continuous and the limit in interior compacts is holomorphic}
{$A(\mathbb D)$ is closed, unital and Banach}

\begin{cthejemplo}{(Polynomial approximation in disk algebra)}
In $A(\mathbb D)$ we are interested in approximating holomorphic functions without leaving the uniform norm.
For $f\in A(\mathbb D)$, the dilations $f_r(z)=f(rz)$ converge uniformly to $f$
in $\overline{\mathbb D}$ when $r\uparrow1$, by uniform continuity of $f$.
Each $f_r$ is holomorphic on a disk of radius $1/r>1$, so its Taylor polynomials
converge uniformly on $\overline{\mathbb D}$.
This explains why the polynomials form a dense subalgebra of the Banach algebra in Problem 1.12(a).
\end{cthejemplo}

\noindent\textbf{Dependencies.} The completeness of $C(\overline{\mathbb D})$ and the Weierstrass theorem on local uniform limits of holomorphic functions guarantee that $A(\mathbb D)$ is closed.

\noindent\textbf{How this problem is used later.} It makes it possible to treat the boundary restriction as an inclusion of Banach algebras in Problem 1.12(b) and to compare spectra in Problem 1.12(c).

\subsection*{1.12(b) The boundary restriction is isometric}

\noindent\textbf{Statement.} Prove that the restriction $f\mapsto f|_{\mathbb T}$ is a unital isometric homomorphism from $A(\mathbb D)$ into $C(\mathbb T)$.

\begin{proof}[Solution]
By the maximum modulus principle,
\[
 \|f\|_{C(\overline{\mathbb D})}=\max_{|z|=1}|f(z)|=\|f|_{\mathbb T}\|_{C(\mathbb T)}.
\]
Thus the restriction is isometric and injective. For $f,g\in A(\mathbb D)$ and $z\in\mathbb T$,
\[
 (f+g)|_{\mathbb T}(z)=f(z)+g(z),\qquad
 (fg)|_{\mathbb T}(z)=f(z)g(z),\qquad
 1|_{\mathbb T}(z)=1.
\]
Therefore it is a unital homomorphism of Banach algebras.
\end{proof}

\proofdiagram
{$f\in A(\mathbb D)$}
{Maximum modulus principle for holomorphic functions}
{$\max_{\overline{\mathbb D}}|f|=\max_{\mathbb T}|f|$}
{$f\mapsto f|_{\mathbb T}$ is an isometric embedding}

\begin{cthejemplo}{(The disk norm is attained on the boundary)}
In $A(\mathbb D)$ we check Problem 1.12(b) with $f(z)=1+z^2$.
For $|z|\le1$, $|f(z)|\le2$, and equality is achieved in $z=\pm1\in\mathbb T$.
So $\|f\|_{C(\overline{\mathbb D})}=\|f|_{\mathbb T}\|_{C(\mathbb T)}=2$:
The boundary already determines the norm of this function.
\end{cthejemplo}

\noindent\textbf{Dependencies.} The maximum modulus principle identifies the norm on the closed disk with the norm on the boundary; therefore the restriction is isometric and injective.

\noindent\textbf{How this problem is used later.} Provides an isometric and unital inclusion in $C(\mathbb T)$ for Problem 1.12(c), where the same function will have different inverses available depending on the algebra.

\subsection*{1.12(c) The same element changes spectrum when changing algebra}

\noindent\textbf{Statement.} For $h(z)=z$, prove $\spec_{A(\mathbb D)}(h)=\overline{\mathbb D}$ and $\spec_{C(\mathbb T)}(h|_{\mathbb T})=\mathbb T$.

\begin{proof}[Solution]
If $|\lambda|>1$, function $(\lambda-z)^{-1}$ is continuous in $\overline{\mathbb D}$ and holomorphic
on $\mathbb D$, so $\lambda-h$ is invertible in $A(\mathbb D)$. If $|\lambda|<1$,
$\lambda-z$ vanishes at $z=\lambda$; if $|\lambda|=1$, any putative inverse would agree on
$\mathbb D$ with $(\lambda-z)^{-1}$, which becomes unbounded when approaching $\lambda$. Therefore
\[
 \spec_{A(\mathbb D)}(h)=\overline{\mathbb D}.
\]

In $C(\mathbb T)$, $\lambda-h$ is invertible exactly when it does not vanish on $\mathbb T$; hence
\[
 \spec_{C(\mathbb T)}(h)=\mathbb T.
\]
\end{proof}

\proofdiagrambranch
{$h(z)=z$ in $A(\mathbb D)$}
{$h|_{\mathbb T}$ in $C(\mathbb T)$}
{Invertibility of $\lambda-h$ depends on whether its reciprocal belongs to the algebra under consideration}
{Invertibility fails in $A(\mathbb D)$ for $\lambda\in\overline{\mathbb D}$; in $C(\mathbb T)$ it fails only on $\mathbb T$}
{$\spec_{A(\mathbb D)}(h)=\overline{\mathbb D}$, $\spec_{C(\mathbb T)}(h)=\mathbb T$}

\begin{cthejemplo}{(The inverse depends on the ambient algebra)}
To compare the two spectra of Problem 1.12(c), let's take $h(z)=z$ in $A(\mathbb D)$ and
its restriction in $C(\mathbb T)$.
In the first algebra, $0-h$ has no holomorphic inverse because $1/z$ has a pole at
$0$; in the second, $h^{-1}(z)=\overline z$ is continuous over $\mathbb T$.
So $0\in\spec_{A(\mathbb D)}(h)$ and $0\notin\spec_{C(\mathbb T)}(h|_{\mathbb T})$.
The same boundary function has an inverse only in the larger algebra. Additionally,
$\overline z$ does not come from $A(\mathbb D)$: if a function $g$ had those boundary values, $zg(z)=1$ in $\mathbb T$ and the maximum modulus principle would give $zg(z)=1$ in all
$\mathbb D$, a contradiction at $z=0$. That is why the inclusion is not of $C^*$-subalgebras.
\end{cthejemplo}

\noindent\textbf{Dependencies.} Spectrum definition and inclusion of Problem 1.12(b) are used; Problem 1.10 calculates the spectrum at $C(\mathbb T)$. In $A(\mathbb D)$ the reciprocal of $\lambda-z$ must also be holomorphic on the disk.

\noindent\textbf{How this problem is used later.} It is the counterexample that motivates the $C^*$-subalgebra hypothesis in spectral permanence from Section 2: a closed and unital Banach subalgebra can change the spectrum.

\chapter{Assignment 2: worked solutions on functional calculus}

\begin{context}
After completing this assignment, you will be able to verify a concrete representation of $C_b(X)$, reconstruct
spectral characterizations of normal elements, and distinguish norm convergence from strong operator
convergence. The dependencies given after each solution show which hypotheses justify each step.
\end{context}

\section{Problem 2.1: \texorpdfstring{$C_b(X)$}{Cb(X)} as a concrete \texorpdfstring{$C^*$}{C-star}-algebra}

Let $X$ be locally compact and Hausdorff.


Let $H=\ell^2(X)$ and, for $f\in C_b(X)$, define
\[
 (M_f\xi)(x)=f(x)\xi(x).
\]
For an arbitrary index set, the sum in $\|\xi\|_2^2$ is understood as
\[
 \|\xi\|_2^2=\sup_{F\subset X\text{ finite}}\sum_{x\in F}|\xi(x)|^2.
\]
We establish the lower norm bound when $X\ne\varnothing$.

\noindent\textbf{Statement.} Prove that $M:C_b(X)\to\mathcal B(\ell^2(X))$, $f\mapsto M_f$, is an isometric $*$-homomorphism.

\begin{proof}[Solution]
If $X=\varnothing$, $C_b(X)$ contains only the empty function and $\ell^2(X)=\{0\}$; therefore
$\mathcal B(\ell^2(X))=\{0\}$ and the unique map $f\mapsto M_f$ preserves operations and
norms. Assume henceforth that $X\ne\varnothing$.

For $\xi\in\ell^2(X)$,
\[
\|M_f\xi\|_2^2
=\sum_{x\in X}|f(x)|^2|\xi(x)|^2
\le\|f\|_\infty^2\|\xi\|_2^2,
\]
so $\|M_f\|\le\|f\|_\infty$. For the reverse inequality, given
$\varepsilon>0$ we choose $x_0$ with $|f(x_0)|>\|f\|_\infty-\varepsilon$ and take the vector
$e_{x_0}$ supported at $x_0$. Then
$\|M_fe_{x_0}\|=|f(x_0)|$, whence $\|M_f\|=\|f\|_\infty$.

Additionally,
\[
M_{f+g}=M_f+M_g,
\qquad M_{fg}=M_fM_g,
\qquad M_1=I.
\]
Finally,
\[
\langle M_f\xi,\eta\rangle
=\sum_x f(x)\xi(x)\overline{\eta(x)}
=\langle \xi,M_{\bar f}\eta\rangle,
\]
so $M_f^*=M_{\bar f}$. Therefore $M$ is an isometric $*$-homomorphism.
\end{proof}

\proofdiagram
{$f\in C_b(X)$}
{Norm in $\ell^2(X)$; pointwise conjugation}
{$\|M_f\|=\|f\|_\infty$, $M_{fg}=M_fM_g$, $M_{\bar f}=M_f^*$}
{$f\mapsto M_f$ is an isometric $*$-homomorphism}

\begin{cthejemplo}{(A diagonal operator on the space of sequences)}
To see the representation of Problem 2.1 over a known space, let's take $X=\mathbb N$ and
$f(n)=1/(n+1)$.
In $H=\ell^2(\mathbb N)$ we have $M_fe_n=(n+1)^{-1}e_n$; also,
\[
 \|M_f\|=\sup_{n\ge1}\|M_fe_n\|=\frac12=\|f\|_\infty.
\]
The bounded function appears as a diagonal operator, and its norm is preserved by
the isometric $*$-homomorphism of Problem 2.1.
\end{cthejemplo}

\noindent\textbf{Dependencies.} The definition of $C_b(X)$ from Part 1, the norm on $\ell^2(X)$, and the Hilbert space characterization of the adjoint. Here $X$ serves as an index set; its topology is not used to define $\ell^2(X)$.

\noindent\textbf{How this problem is used later.} It exhibits an isometric $*$-monomorphism into $\mathcal B(H)$, a concrete model of the homomorphisms and representations studied in Section 2.

\section{Problem 2.2: spectra of idempotents, products, and unitaries}

Let $\A$ be a unital Banach algebra. For the conclusion about unitaries in part (c), $\A$
is also a $C^*$-algebra.

\subsection*{2.2(a) Spectrum of an idempotent}

\noindent\textbf{Statement.} If $e\in\A$ and $e^2=e$, prove $\spec(e)\subset\{0,1\}$.

\begin{proof}[Solution]
If $\lambda\neq0,1$, set
\[
 b=\frac1\lambda 1+\frac{1}{\lambda(\lambda-1)}e.
\]
Since $e(1-e)=(1-e)e=0$ and $e^2=e$, we can write
\[
 \lambda1-e=\lambda(1-e)+(\lambda-1)e,
 \qquad b=\lambda^{-1}(1-e)+(\lambda-1)^{-1}e.
\]
On multiplication, the cross terms vanish and
$(\lambda1-e)b=b(\lambda1-e)=(1-e)+e=1$. Thus every
$\lambda\notin\{0,1\}$ belongs to the resolvent set and
$\spec(e)\subset\{0,1\}$.
\end{proof}

\pagebreak
\proofdiagram
{$e^2=e$}
{Solve algebraically $(\lambda1-e)^{-1}$}
{For $\lambda\notin\{0,1\}$, construct $\lambda^{-1}1+[\lambda(\lambda-1)]^{-1}e$}
{$\spec(e)\subset\{0,1\}$}

\begin{cthejemplo}{(Spectrum of a matrix projection)}
In $M_2(\mathbb C)$ we want to identify the two points allowed by Problem 2.2(a).
For $P=\operatorname{diag}(1,0)$, $P^2=P$ and
$\lambda I-P=\operatorname{diag}(\lambda-1,\lambda)$.
This matrix is invertible exactly when $\lambda\ne0,1$; therefore
$\spec(P)=\{0,1\}$. If $P$ is replaced by $0$ or $I$, the same calculation leaves only
$\{0\}$ or $\{1\}$, respectively.
\end{cthejemplo}

\noindent\textbf{Dependencies.} The equation $e^2=e$ allows us to check an explicit inverse of $\lambda1-e$ for $\lambda\ne0,1$; the definition of spectrum in Part 1 applies.

\noindent\textbf{How this problem is used later.} This gives the forward implication of the spectral characterization of projections in Problem 2.6(c) and in Proposition 2.28.

\subsection*{2.2(b) Spectral radius of a commutative product}

\noindent\textbf{Statement.} If $x,y\in\A$ commute, prove $r(xy)\le r(x)r(y)$.

\begin{proof}[Solution]
Commutativity allows us to write $(xy)^n=x^ny^n$. So
\[
\|(xy)^n\|^{1/n}
\le\|x^n\|^{1/n}\|y^n\|^{1/n}.
\]
Passing to the limit and using the spectral radius formula,
$r(xy)\le r(x)r(y)$.
\end{proof}

\proofdiagram
{$xy=yx$}
{Spectral radius formula}
{$(xy)^n=x^ny^n$ and $\|(xy)^n\|\le\|x^n\|\|y^n\|$}
{$r(xy)\le r(x)r(y)$}

\begin{cthejemplo}{(The spectral radius inequality can be strict)}
In $C([0,1])$ we can check how much the bound of Problem 2.2(b) can be improved.
Let's take $f(t)=t$ and $g(t)=1-t$; they commute because multiplication is pointwise.
By Problem 1.10, $r(f)=\|f\|_\infty=1$ and $r(g)=1$, while
\[
 r(fg)=\max_{0\le t\le1}t(1-t)=\frac14<r(f)r(g)=1.
\]
The inequality is strict because the maxima of $f$ and $g$ occur at opposite endpoints.
\end{cthejemplo}

\noindent\textbf{Dependencies.} The commutation allows $(xy)^n=x^ny^n$ and the spectral radius formula from Part 1 converts the estimate of norms to the desired inequality.

\noindent\textbf{How this problem is used later.} It allows estimating products within commutative subalgebras such as $C^*(x,1)$. The classification of Problem 2.6 uses functional calculus and does not depend on this inequality.

\subsection*{2.2(c) Spectrum of an invertible with controlled norm}

\noindent\textbf{Statement.} If $z\in\GL(\A)$ and $\|z\|,\|z^{-1}\|\le1$, prove $\spec(z)\subset\mathbb T$; deduce the same inclusion for every unitary element of a unital $C^*$-algebra.

\begin{proof}[Solution]
If $\lambda\in\spec(z)$, then $|\lambda|\le r(z)\le\|z\|\le1$. Since $z$ is invertible,
$0\notin\spec(z)$. Additionally
\[
 \lambda\in\spec(z)\quad\Longleftrightarrow\quad
 \lambda^{-1}\in\spec(z^{-1}),
\]
because
$\lambda1-z=-\lambda z(\lambda^{-1}1-z^{-1})$. Therefore
$|\lambda|^{-1}\le r(z^{-1})\le1$, that is $|\lambda|\ge1$. Hence $|\lambda|=1$.

If $u$ is unitary in a $C^*$-algebra, $u^{-1}=u^*$ and
$\|u\|=\|u^*\|=1$, so the above result gives
$\spec(u)\subset\mathbb T$.
\end{proof}

\proofdiagrambranch
{$z\in\GL(\A)$ and $\|z\|\le1$}
{$\|z^{-1}\|\le1$}
{$\lambda\in\spec(z)\Rightarrow |\lambda|\le1$ and $\lambda^{-1}\in\spec(z^{-1})$}
{$|\lambda|^{-1}\le1$}
{$|\lambda|=1$, then $\spec(z)\subset\mathbb T$}

\begin{cthejemplo}{(Spectrum of a unitary diagonal matrix)}
To see the unitary conclusion of Problem 2.2(c) with two spectral values, let's take
$U=\operatorname{diag}(1,i)$ in $M_2(\mathbb C)$.
We have $U^*U=I$, $U^{-1}=U^*$ and $\|U\|=\|U^{-1}\|=1$.
Moreover, $\det(\lambda I-U)=(\lambda-1)(\lambda-i)$, so
$\spec(U)=\{1,i\}\subset\mathbb T$.
The norm bounds force both points to be on the unit circle.
\end{cthejemplo}

\noindent\textbf{Dependencies.} The bound $r(z)\le\|z\|$, the relationship between the spectra of an invertible and its inverse, and both hypotheses $\|z\|,\|z^{-1}\|\le1$ are used.

\noindent\textbf{How this problem is used later.} This gives the unitary implication of Problem 2.6(b) and Proposition 2.28; for a unitary $u$, the identity $C^*$ ensures $\|u\|=\|u^{-1}\|=1$.

\begin{investigacion}[title={Unitary rotations on $L^2(\mathbb T)$}]
To apply Problem 2.2(c) in the Hilbert space model of \cite{BFK2006}, fix $z\in\mathbb T$ and
define $(U_z\varphi)(t)=\varphi(zt)$ on $L^2(\mathbb T)$.
The invariance of the Haar measure under rotations gives $\|U_z\varphi\|_2=\|\varphi\|_2$;
also $U_{z^{-1}}U_z=U_zU_{z^{-1}}=I$. Thus $U_z$ is unitary and
$\|U_z\|=\|U_z^{-1}\|=1$.
By Problem 2.2(c), $\spec(U_z)\subset\mathbb T$. The conclusion uses those properties
verified on $L^2(\mathbb T)$.
\end{investigacion}

\pagebreak
\section{Problem 2.3: products of \texorpdfstring{$C^*$}{C-star}-algebras and the ideal \texorpdfstring{$c_0$}{c0}}

Let $I$ be a set of indices and $(A_i)_{i\in I}$ be a family of $C^*$-algebras. In part
(b), take $I=\mathbb N$.

\subsection*{2.3(a) Bounded product of a family of $C^*$-algebras}

\noindent\textbf{Statement.} Prove that the product of families $(a_i)_{i\in I}$ with $a_i\in A_i$ and $\sup_i\|a_i\|<\infty$ is a $C^*$-algebra with pointwise operations and supremum norm.

\begin{proof}[Solution]
Let
\[
\prod_{i\in I}^{\ell^\infty} A_i
=\{(a_i): a_i\in A_i,\ \sup_i\|a_i\|<\infty\}.
\]
Pointwise operations preserve uniform boundedness; in particular,
\[
 \sup_i\|a_i+b_i\|\le\sup_i\|a_i\|+\sup_i\|b_i\|,
 \qquad
 \sup_i\|a_ib_i\|\le\bigl(\sup_i\|a_i\|\bigr)
                        \bigl(\sup_i\|b_i\|\bigr).
\]
With the norm $\|(a_i)\|=\sup_i\|a_i\|$, this is a normed algebra. If
$(a^{(n)})$ is Cauchy, each coordinate $a_i^{(n)}$ converges in $A_i$ to an element $a_i$.
Fix $\varepsilon>0$. There is $N$ such that
$\sup_i\|a_i^{(n)}-a_i^{(m)}\|<\varepsilon$ for $n,m\ge N$. Letting $m\to\infty$
coordinate by coordinate, we get $\|a_i^{(n)}-a_i\|\le\varepsilon$ for all $i$ and $n\ge N$.
In particular $\sup_i\|a_i\|\le\|a^{(N)}\|+\varepsilon<\infty$, and the same estimate shows
$\|a^{(n)}-a\|\le\varepsilon$ for $n\ge N$. Therefore the bounded product is complete.
Defining $(a_i)^*=(a_i^*)$,
\[
\|(a_i)^*(a_i)\|
=\sup_i\|a_i^*a_i\|
=\sup_i\|a_i\|^2
=\|(a_i)\|^2.
\]
This is a $C^*$-algebra.
\end{proof}

\proofdiagram
{$\{A_i\}_{i\in I}$ $C^*$-algebras}
{Coordinatewise operations; supremum norm}
{Coordinate completeness and $\|(a_i)^*(a_i)\|=\sup_i\|a_i^*a_i\|$}
{The bounded product $\prod_i A_i$ is a $C^*$-algebra}

\begin{cthejemplo}{(A bounded product of matrices)}
In the product of Problem 2.3(a), each coordinate can be a different matrix.
Let's take $A_n=M_2(\mathbb C)$ and $a_n=\operatorname{diag}(1,1/n)$ for $n\ge1$.
The family $a=(a_n)$ belongs to the bounded product because $\sup_n\|a_n\|=1$; moreover,
\[
 a_n^*a_n=\operatorname{diag}(1,1/n^2),
 \qquad \|a^*a\|=1=\|a\|^2.
\]
The $C^*$-identity for the product is verified coordinate by coordinate, as in Problem 2.3(a).
\end{cthejemplo}

\noindent\textbf{Dependencies.} Coordinate-by-coordinate operations, completeness of each $A_i$, and the $C^*$-identity are applied on each coordinate. The supremum norm restricts the product to uniformly bounded families.

\noindent\textbf{How this problem is used later.} This constructs the ambient algebra of Problem 2.3(b), where families tending to zero form a closed ideal.

\subsection*{2.3(b) The sum $c_0$ is a closed ideal}

\noindent\textbf{Statement.} For $I=\mathbb N$, prove that the sequences $(a_n)$ of the bounded product satisfying $\|a_n\|\to0$ form a closed two-sided ideal.

\begin{proof}[Solution]
Let
\[
\bigoplus_{n\in\mathbb N}^{c_0}A_n
=\{(a_n)\in\prod^{\ell^\infty}A_n:\|a_n\|\to0\}.
\]
If $\|a_n\|,\|c_n\|\to0$, then
$\|a_n+c_n\|\le\|a_n\|+\|c_n\|\to0$ and
$\|\alpha a_n\|=|\alpha|\|a_n\|\to0$; therefore the $c_0$-sum is a linear subspace.
If $b=(b_n)$ is in the bounded product and $a=(a_n)$ is in the $c_0$-sum, then
\[
\|b_na_n\|\le\|b\|_\infty\|a_n\|\to0,
\quad
\|a_nb_n\|\le\|a_n\|\|b\|_\infty\to0,
\]
so it is a two-sided ideal. If $a^{(k)}\to a$ in supremum norm and each $a^{(k)}$ tends to zero in
coordinates, fix a positive $\varepsilon$ and choose $k$ with
$\|a-a^{(k)}\|_\infty<\varepsilon/2$ and then $N$ with
$\|a_n^{(k)}\|<\varepsilon/2$ for $n\ge N$. Then $\|a_n\|<\varepsilon$ for $n\ge N$.
Therefore the ideal is closed.
Furthermore $(a_n^*)$ belongs to the same ideal because $\|a_n^*\|=\|a_n\|\to0$.
\end{proof}

\proofdiagram
{$a=(a_n)$ with $\|a_n\|\to0$}
{Coordinatewise product and supremum norm}
{$\|b_na_n\|\le\|b\|_\infty\|a_n\|\to0$; uniform limit preserves condition}
{$\bigoplus_n A_n$ is a closed two-sided ideal}

\begin{cthejemplo}{(The classical ideal of null sequences)}
To recognize the construction of Problem 2.3(b), let's put $A_n=\mathbb C$ for each $n$.
The sequence $a_n=1/n$ belongs to $c_0$, while $b_n=(-1)^n$ is in $\ell^\infty$.
Their product satisfies $|b_na_n|=1/n\to0$, and the same goes for $a_nb_n$.
Thus the ideal $c_0\subset\ell^\infty$ is the scalar case of the two-sided ideal property proved in Problem 2.3(b).
\end{cthejemplo}

\noindent\textbf{Dependencies.} The bounded product of Problem 2.3(a) provides the supremum norm. Coordinate-by-coordinate estimates prove the ideal property, and the uniform criterion proves closedness.

\noindent\textbf{How this problem is used later.} This allows us to form the quotient of the product by the ideal $c_0$ using Problem 1.4; interpreting that quotient requires further results about quotients of $C^*$-algebras.

\section{Problem 2.4: continuous spectral mapping theorem}

Let $\A$ be a unital $C^*$-algebra, let $x\in\A$ be normal, and let $f:\spec(x)\to\mathbb C$ be continuous.

\noindent\textbf{Statement.} Prove the continuous spectral mapping theorem: $\spec_\A(f(x))=f(\spec_\A(x))$.

\begin{proof}[Solution]
Work in $B=C^*(x,1)$, which is commutative. By spectral permanence, the spectrum of
$f(x)$ is the same whether computed in $B$ or in $\A$. The functional calculus is a $*$-isomorphism
\[
 C(\spec(x))\longrightarrow B.
\]
Under this isomorphism, $f(x)$ corresponds to the function $f$. By Problem 1.10, the
spectrum of $f$ in $C(\spec(x))$ is its range. Therefore
\[
 \spec_\A(f(x))=\spec_B(f(x))=f(\spec(x)).
\]
\end{proof}

\proofdiagrambranch
{$x$ normal}
{$f\in C(\spec(x))$}
{Functional calculus + spectral permanence in $C^*(x,1)$}
{Identify $C^*(x,1)\cong C(\spec(x))$}
{$\spec(f(x))=f(\spec(x))$}

\begin{cthejemplo}{(Squaring a self-adjoint element using its spectrum)}
In $M_2(\mathbb C)$, take $x=\operatorname{diag}(-2,3)$ to check Problem 2.4.
The element is self-adjoint and $\spec(x)=\{-2,3\}$.
As $x^2=\operatorname{diag}(4,9)$,
\[
 \spec(x^2)=\{4,9\}=\{\lambda^2:\lambda\in\spec(x)\}.
\]
The square transforms each spectral value by $t\mapsto t^2$ and produces a nonnegative spectrum.
\end{cthejemplo}

\noindent\textbf{Dependencies.} The functional calculus for normal $x$ identifies $C^*(x,1)$ with $C(\spec(x))$; the spectral permanence of Section 2 and Problem 1.10 allow the spectrum of $f(x)$ to be calculated in the ambient algebra.

\noindent\textbf{How this problem is used later.} This transfers equalities and spectral restrictions from functions to operators; Problem 2.6 applies it to recognize self-adjoint elements, unitary elements, and normal projections.

\section{Problem 2.5: \texorpdfstring{$X\cong\Omega(C(X))$}{X is homeomorphic to Omega(C(X))}}

Let $X$ be a compact Hausdorff space and, for $x\in X$, let
$\operatorname{ev}_x(f)=f(x)$ be the evaluation map on $C(X)$.

\noindent\textbf{Statement.} Prove that $x\mapsto\operatorname{ev}_x$ defines a homeomorphism $X\to\Omega(C(X))$.

\begin{proof}[Solution]
For $f,g\in C(X)$, we have
$\operatorname{ev}_x(fg)=f(x)g(x)$ and $\operatorname{ev}_x(1)=1$;
also $|\operatorname{ev}_x(f)|\le\|f\|_\infty$. Thus each evaluation is a continuous character
and the map is well defined. It is injective because
if $x\neq y$, there is $f\in C(X)$ with $f(x)\neq f(y)$.

Now let $\varphi\in\Omega(C(X))$ and set $I=\ker\varphi$, a maximal ideal. We claim that there is
$x_0\in X$ such that $I\subset I_{x_0}:=\{f:f(x_0)=0\}$. Otherwise, for each $x\in X$ there would be
$f_x\in I$ with $f_x(x)\neq0$. The open sets
$U_x=\{y:|f_x(y)|>0\}$ would cover $X$; by compactness we could choose
$f_1,\dots,f_n\in I$ with $\sum_j|f_j|^2>0$ throughout $X$. Then
\[
 g:=\sum_{j=1}^n f_j^*f_j\in I
\]
would be a strictly positive function and therefore invertible in $C(X)$, impossible for a proper ideal. Thus there is $x_0$ with $I\subset I_{x_0}$. Since $I_{x_0}$ is proper, maximality gives
$I=I_{x_0}$. For each $f\in C(X)$, we have $f-f(x_0)1\in I=\ker\varphi$; since
$\varphi(1)=1$, it follows that $\varphi(f)=f(x_0)$. Therefore
$\varphi=\operatorname{ev}_{x_0}$.

The map $x\mapsto\operatorname{ev}_x$ is continuous for the weak-$*$ topology because, for each
$f\in C(X)$, the coordinate $x\mapsto\operatorname{ev}_x(f)=f(x)$ is continuous. As $X$ is compact
and $\Omega(C(X))$ is Hausdorff, a continuous bijection is a homeomorphism.
\end{proof}

\proofdiagrambranch
{$x\mapsto\operatorname{ev}_x$}
{Continuous functions separate points in compact Hausdorff spaces}
{For surjectivity, $\ker\varphi$ is a maximal ideal and must be $I_{x_0}$ for some $x_0$}
{Continuous bijection between compact Hausdorff spaces}
{$X\cong\Omega(C(X))$}

\begin{cthejemplo}{(A character given by a point on the circle)}
In $C(\mathbb T)$ we want to make the homeomorphism of Problem 2.5 visible.
Let's set $w=e^{i\theta}\in\mathbb T$ and consider the evaluation
$\operatorname{ev}_w(f)=f(w)$. For functions $h(z)=z$ and $g(z)=1+z$,
\[
 \operatorname{ev}_w(h)=w,\qquad \operatorname{ev}_w(g)=1+w.
\]
The value of the character on $h$ identifies the point $w$; Problem 2.5 ensures that every character in
$C(\mathbb T)$ arises from such an evaluation.
\end{cthejemplo}

\noindent\textbf{Dependencies.} Characters and maximal ideals from Part 1, point separation from $C(X)$, and compactness from $X$ are used. The common-zero-set argument uses $\sum_j|f_j|^2$.

\noindent\textbf{How this problem is used later.} This completes the uniqueness of the compact space in commutative Gelfand--Naimark: a $*$-isomorphism $C(X)\cong C(Y)$ induces $X\cong Y$.

\section{Problem 2.6: spectral characterizations and the role of normality}

Let $\A$ be a unital $C^*$-algebra and let $x\in\A$ be normal. Part (d) asks us to show that
normality cannot be omitted from the three converses.

\subsection*{2.6(a) Self-adjoint $\Longleftrightarrow$ real spectrum for normal elements}

\noindent\textbf{Statement.} Prove that $x=x^*$ if and only if $\spec(x)\subset\mathbb R$.

\begin{proof}[Solution]
The direct implication has already been proven in Section 2. For the converse, by functional calculus,
$x$ corresponds to the function $z$ and $x^*$ to $\overline z$ on $\spec(x)$. If the spectrum is real,
the functions agree, so $x=x^*$.
\end{proof}

\proofdiagram
{$x$ normal}
{Real spectrum of a self-adjoint element; functional calculus}
{If $\spec(x)\subset\mathbb R$, then $z=\overline z$ on the spectrum}
{$x=x^*$}

\begin{cthejemplo}{(Real spectrum without self-adjointness)}
The converse of Problem 2.6(a) requires normality; a triangular matrix allows you to see it.
For $N=\begin{pmatrix}0&1\\0&0\end{pmatrix}$,
$\det(\lambda I-N)=\lambda^2$, then $\spec(N)=\{0\}\subset\mathbb R$.
But $N^*=\begin{pmatrix}0&0\\1&0\end{pmatrix}\ne N$.
Additionally $N^*N=\operatorname{diag}(0,1)\ne\operatorname{diag}(1,0)=NN^*$.
A real spectrum does not force self-adjointness without normality, as shown in Problem 2.6(d).
\end{cthejemplo}

\noindent\textbf{Dependencies.} The reality of the spectrum of a self-adjoint element was proved in Section 2; the converse uses functional calculus in $C^*(x,1)$, available because $x$ is normal.

\noindent\textbf{How this problem is used later.} This identifies exactly which spectral condition characterizes self-adjointness. Problem 2.6(d) shows with a nilpotent matrix why the normality hypothesis is necessary.

\subsection*{2.6(b) Unitary $\Longleftrightarrow$ spectrum in $\mathbb T$ for normal elements}

\noindent\textbf{Statement.} Prove that $x$ is unitary if and only if $\spec(x)\subset\mathbb T$.

\begin{proof}[Solution]
If $x$ is unitary, Problem 2.2(c) gives $\spec(x)\subset\mathbb T$. Conversely, if $x$ is
normal and the spectrum lies in $\mathbb T$, then
\[
 x^*x=(\overline z z)(x)=(|z|^2)(x)=1.
\]
As $x$ is normal, $xx^*=x^*x=1$, so it is unitary.
\end{proof}

\proofdiagram
{$x$ normal}
{Problem 2.2(c); functional calculus}
{$\spec(x)\subset\mathbb T\Rightarrow |z|^2=1$ on spectrum}
{$x^*x=(\overline z z)(x)=1$ and, by normality, $xx^*=1$}

\begin{cthejemplo}{(Spectrum on the unit circle without unitarity)}
To show the limitation of the converse in Problem 2.6(b), let's take
$J=\begin{pmatrix}1&1\\0&1\end{pmatrix}$.
Its characteristic polynomial is $(\lambda-1)^2$, so
$\spec(J)=\{1\}\subset\mathbb T$, but
$J^*J=\begin{pmatrix}1&1\\1&2\end{pmatrix}\ne I$.
The spectrum is on the circle and the matrix is non-unitary: Problem 2.6(d) explains the role of normality.
\end{cthejemplo}

\noindent\textbf{Dependencies.} Problem 2.2(c) gives the forward implication. In the converse, normality allows functional calculus and the identity $x^*x=(|z|^2)(x)$.

\noindent\textbf{How this problem is used later.} It allows recognizing unitary elements from their spectrum, for example in algebras of operators with unitary shifts when their hypotheses have been verified.

\subsection*{2.6(c) Projection $\Longleftrightarrow$ spectrum on $\{0,1\}$ for normals}

\noindent\textbf{Statement.} Prove that $x$ is a projection if and only if $\spec(x)\subset\{0,1\}$.

\begin{proof}[Solution]
If $x$ is a projection, it is idempotent and Problem 2.2(a) gives
$\spec(x)\subset\{0,1\}$. For the converse, if $x$ is normal and its spectrum is contained in
$\{0,1\}$, functional calculus gives
\[
 x^2-x=(z^2-z)(x)=0,
 \qquad
 x^*-x=(\overline z-z)(x)=0.
\]
So $x$ is a projection.
\end{proof}

\proofdiagram
{$x$ normal}
{Spectrum of an idempotent; functional calculus}
{$z^2-z=0$ and $\overline z-z=0$ on $\{0,1\}$}
{$x^2=x$ and $x^*=x$}

\begin{cthejemplo}{(An idempotent with the expected spectrum that is not a projection)}
In Problem 2.6(c), self-adjointness is part of the definition of a projection. Let's consider
$E=\begin{pmatrix}1&1\\0&0\end{pmatrix}$.
A calculation gives $E^2=E$ and, by triangularity, $\spec(E)=\{0,1\}$; however,
$E^*=\begin{pmatrix}1&0\\1&0\end{pmatrix}\ne E$.
Thus, the permitted spectrum and idempotence are not enough without the normality used by the
converse of Problem 2.6(c).
\end{cthejemplo}

\noindent\textbf{Dependencies.} Problem 2.2(a) gives the forward implication. For the converse, normality allows $z^2=z$ and $\overline z=z$ to be transported from the spectrum to identities of $x$.

\noindent\textbf{How this problem is used later.} This connects the spectral classification of Section 2 with Cauchy projections only when, in addition to idempotence, the self-adjointness of those projections is verified.

\subsection*{2.6(d) Where is normality used?}

\noindent\textbf{Statement.} Display in $M_2(\mathbb C)$ three non-normal elements that invalidate the converses of parts (a)--(c) without normality, and locate the step of each proof that uses that hypothesis.

\begin{proof}[Solution]
Let's take
\[
 N=\begin{pmatrix}0&1\\0&0\end{pmatrix},
 \qquad J=I+N=\begin{pmatrix}1&1\\0&1\end{pmatrix},
 \qquad E=\begin{pmatrix}1&1\\0&0\end{pmatrix}.
\]
As they are triangular, $\spec(N)=\{0\}\subset\mathbb R$,
$\spec(J)=\{1\}\subset\mathbb T$ and $\spec(E)=\{0,1\}$.
However, $N\ne N^*$,
\[
 J^*J=\begin{pmatrix}1&1\\1&2\end{pmatrix}\ne I.
\]
$E$ is not a projection because, although $E^2=E$, it has $E^*\ne E$.
All three are non-normal: indeed,
$N^*N=\operatorname{diag}(0,1)\ne\operatorname{diag}(1,0)=NN^*$ and
$J^*J-JJ^*=N^*N-NN^*\ne0$; also,
$E^*E=\begin{pmatrix}1&1\\1&1\end{pmatrix}$ and
$EE^*=\operatorname{diag}(2,0)$ are different.

Normality was used in parts (a)--(c) so that $x$ and $x^*$ generate the
commutative $C^*$-algebra $C^*(x,1)$ and continuous functional calculus can be applied.
In (b) the equality $xx^*=x^*x$ was also used after obtaining $x^*x=1$.
\end{proof}

\proofdiagram
{$N$, $J=I+N$ and $E$ are non-normal}
{Continuous functional calculus was only built for normals}
{Functional calculus transfers equality of functions on the spectrum to equality of elements}
{Without normality, the spectrum does not determine these algebraic relationships}

\begin{cthejemplo}{(Two operators with the same spectrum)}
To see what the spectrum loses without normality, let's compare $I$ with
$J=\begin{pmatrix}1&1\\0&1\end{pmatrix}$ in $M_2(\mathbb C)$.
Both have spectrum $\{1\}$, but $I^*I=I$, while
$J^*J=\begin{pmatrix}1&1\\1&2\end{pmatrix}\ne I$.
Only $I$ is unitary. The concrete comparison explains why Problem 2.6(d) requires normality
before retrieving algebraic properties from the spectrum.
\end{cthejemplo}

\noindent\textbf{Dependencies.} Parts 2.6(a)--(c) use functional calculus on $C^*(x,1)$; that subalgebra is commutative because $x$ is normal. Jordan blocks provide counterexamples.

\noindent\textbf{How this problem is used later.} This delimits the scope of Proposition 2.28 and prevents us from inferring self-adjointness, unitarity or projection from the spectrum of a non-normal operator.

\pagebreak
\section{Problem 2.7: Unitary part of polar decomposition}

Let $\A$ be a unital $C^*$-algebra. We define $|a|$ for every $a\in\A$ and assume that $a$ is
invertible when constructing the unitary factor.


We define $|a|=(a^*a)^{1/2}$ by positive functional calculus.

\noindent\textbf{Statement.} For $a\in\GL(\A)$, prove that $u=a|a|^{-1}$ is unitary, where $|a|=(a^*a)^{1/2}$.

\begin{proof}[Solution]
If $a$ is invertible, then $a^*a$ has inverse $a^{-1}(a^*)^{-1}$.
Since $|a|^2=a^*a$ is invertible, so is $|a|$: the elements $|a|$ and
$(|a|^2)^{-1}$ commute, and $|a|(|a|^2)^{-1}$ is a two-sided inverse of $|a|$.
We can, therefore, define $u=a|a|^{-1}$.
As $|a|$ is positive, $|a|^*=|a|$ and
\[
 u^*=|a|^{-1}a^*.
\]
Then
\[
 u^*u=|a|^{-1}a^*a|a|^{-1}
 =|a|^{-1}|a|^2|a|^{-1}=1.
\]
Moreover,
\[
 uu^*=a|a|^{-2}a^*
 =a(a^*a)^{-1}a^*
 =a\,a^{-1}(a^*)^{-1}a^*=1.
\]
Thus $u$ is unitary and $a=u|a|$.
\end{proof}

\proofdiagram
{$a\in\GL(\A)$}
{Positive root $|a|=(a^*a)^{1/2}$}
{Define $u=a|a|^{-1}$ and calculate $u^*u$, $uu^*$}
{$u^*u=uu^*=1$; $u$ is unitary}

\begin{cthejemplo}{(Polar decomposition of a diagonal matrix)}
To apply the formula in Problem 2.7, let's take the invertible matrix
$A=\operatorname{diag}(2i,3)$.
Then $A^*A=\operatorname{diag}(4,9)$,
$|A|=\operatorname{diag}(2,3)$ and
\[
 U=A|A|^{-1}=\operatorname{diag}(i,1),\qquad U^*U=UU^*=I.
\]
The equality $A=U|A|$ separates a unitary part from a positive part, as stated in Problem 2.7.
\end{cthejemplo}

\noindent\textbf{Dependencies.} The positive root $|a|=(a^*a)^{1/2}$ comes from the functional calculus of Section 2. Invertibility of $a$ implies invertibility of both $a^*a$ and $|a|$.

\noindent\textbf{How this problem is used later.} This constructs the polar decomposition $a=u|a|$ of an invertible element and clarifies where invertibility is used to obtain a unitary factor in the same algebra.

\Needspace{15\baselineskip}
\section{Problem 2.8: continuity of functional calculus under convergence of normal elements}

Let $\A$ be a unital $C^*$-algebra and let $(x_n)$ be a sequence of normal elements
converging in norm to $x$. Fix a compact set $K\subset\mathbb C$ containing $\spec(x_n)$
for every $n$ and $\spec(x)$, and let $f\in C(K)$.

\noindent\textbf{Statement.} Prove that $x$ is normal and that $f(x_n)\to f(x)$ in norm for all $f\in C(K)$.

\begin{proof}[Solution]
By Problem 1.7, $x$ is normal. Let $f\in C(K)$ and $\varepsilon>0$. Stone--Weierstrass provides
a polynomial $p(z,\overline z)$ such that
\[
 \sup_{z\in K}|f(z)-p(z,\overline z)|<\varepsilon.
\]
Functional calculus is contractive, so for all $n$,
\[
\|f(x_n)-p(x_n,x_n^*)\|<\varepsilon,
\qquad
\|f(x)-p(x,x^*)\|<\varepsilon.
\]
Since $x_n\to x$, we also have $x_n^*\to x^*$; continuity of addition and multiplication gives
$p(x_n,x_n^*)\to p(x,x^*)$. Therefore, for large $n$,
\[
\|f(x_n)-f(x)\|
\le 2\varepsilon+\|p(x_n,x_n^*)-p(x,x^*)\|<3\varepsilon.
\]
As $\varepsilon$ is arbitrary, $f(x_n)\to f(x)$.
\end{proof}

\proofdiagrambranch
{$x_n\to x$ and each $x_n$ is normal}
{$\spec(x_n),\spec(x)\subset K$ compact}
{Closedness of normal elements; Stone--Weierstrass approximates $f$ by polynomials in $z,\bar z$}
{Polynomials are continuous in $x_n,x_n^*$; control error uniformly}
{$f(x_n)\to f(x)$}

\begin{cthejemplo}{(Continuity of the positive square root)}
To see Problem 2.8 in positive matrices, let's set
$x_n=\operatorname{diag}(1+1/n,4)$ and $x=\operatorname{diag}(1,4)$.
All these spectra lie in $K=[1,5]$ and $x_n\to x$ in norm.
With $f(t)=\sqrt t$,
\[
 f(x_n)=\operatorname{diag}(\sqrt{1+1/n},2)
 \longrightarrow\operatorname{diag}(1,2)=f(x).
\]
The calculation shows that the positive square root varies continuously under the assumptions of Problem 2.8.
\end{cthejemplo}

\noindent\textbf{Dependencies.} Problem 1.7 guarantees that $x$ is normal; Stone--Weierstrass approximates $f$ uniformly by polynomials in $z,\overline z$, and functional calculus controls the error.

\noindent\textbf{How this problem is used later.} Shows the stability in norm of $f(x_n)$, in particular of positive roots when the spectra are in one compact subset of $[0,\infty)$.

\Needspace{17\baselineskip}
\section{Problem 2.9: monotone convergence of self-adjoint operators}

Let $H$ be a Hilbert space. A net is indexed by a directed set $\Lambda$.
For self-adjoint $R,U\in\mathcal B(H)$, we write $R\le U$ when $U-R$ is positive,
that is, when $\langle(U-R)\xi,\xi\rangle\ge0$ for all $\xi\in H$.


Let $(T_\lambda)_{\lambda\in\Lambda}$ be a bounded net of self-adjoint operators with
$T_{\lambda_1}\le T_{\lambda_2}$ if $\lambda_1\le\lambda_2$.
Here ``pointwise convergence'' means that, for each $\xi\in H$,
$\|T_\lambda\xi-T\xi\|\to0$; it is strong operator convergence.

\noindent\textbf{Statement.} Prove that there exists a self-adjoint $T\in\mathcal B(H)$ such that $\|T_\lambda\xi-T\xi\|\to0$ for each $\xi\in H$.

\begin{proof}[Solution]
Let $M=\sup_\lambda\|T_\lambda\|$. For each $\xi\in H$, the numbers
$q_\lambda(\xi)=\langle T_\lambda\xi,\xi\rangle$ form an increasing net bounded by
$M\|\xi\|^2$, so they converge to a limit $q(\xi)$. Using the polarization identity we obtain
a bounded sesquilinear form
\[
 B(\xi,\eta)=\lim_\lambda\langle T_\lambda\xi,\eta\rangle.
\]
Indeed, $|\langle T_\lambda\xi,\eta\rangle|\le M\|\xi\|\|\eta\|$, and the same bound passes to the
limit.
With the convention that the inner product is linear in the first variable, the polarization identity
expresses this limit using four values of $q$:
\[
 B(\xi,\eta)=\frac14\bigl(q(\xi+\eta)-q(\xi-\eta)
   +i q(\xi+i\eta)-i q(\xi-i\eta)\bigr).
\]
By the Riesz representation theorem there is a unique $T\in\mathcal B(H)$ with
$B(\xi,\eta)=\langle T\xi,\eta\rangle$, and the bound of $B$ implies $\|T\|\le M$.
Since each $T_\lambda$ is self-adjoint, so is $T$, and
$T_\lambda\le T$.

Set $A_\lambda=T-T_\lambda\ge0$. Then
$\|A_\lambda\|\le 2M$. Since
$\spec(A_\lambda)\subset[0,\|A_\lambda\|]$, the function
$t\mapsto t(\|A_\lambda\|-t)$ is nonnegative on that spectrum. Positive functional calculus
gives $0\le A_\lambda^2\le\|A_\lambda\|A_\lambda\le2M A_\lambda$. Therefore
\[
\|(T-T_\lambda)\xi\|^2
=\langle A_\lambda^2\xi,\xi\rangle
\le2M\langle A_\lambda\xi,\xi\rangle.
\]
The last term tends to zero by definition of $T$. So
$T_\lambda\xi\to T\xi$ in norm for each $\xi$.
\end{proof}

\proofdiagrambranch
{Monotonicity: increasing $\langle T_\lambda\xi,\xi\rangle$}
{Uniform boundedness $\sup\|T_\lambda\|<\infty$}
{Define a limit form and use Riesz to obtain $T$}
{Use $0\le(T-T_\lambda)^2\le\|T-T_\lambda\|(T-T_\lambda)$}
{$T_\lambda\xi\to T\xi$ for each $\xi$}

\begin{cthejemplo}{(Strong convergence of increasing matrices)}
To see the limit in Problem 2.9, take in $\mathbb C^2$
$T_n=\operatorname{diag}(1-1/n,2-1/n)$, $n\ge1$.
We have $T_{n+1}-T_n=(1/n-1/(n+1))I\ge0$ and $\|T_n\|\le2$.
With $T=\operatorname{diag}(1,2)$,
\[
 \|(T-T_n)\xi\|=\frac1n\|\xi\|\longrightarrow0
 \qquad(\xi\in\mathbb C^2).
\]
The convergence here is even in operator norm; Problem 2.9 ensures strong convergence
for bounded nets on any Hilbert space.
\end{cthejemplo}

\noindent\textbf{Dependencies.} Convergence of bounded monotone nets of real numbers, polarization identity, Riesz representation, and positive functional calculus are used to justify $0\le A^2\le\|A\|A$ when $A\ge0$.

\noindent\textbf{How this problem is used later.} This gives strong convergence: $\|T_\lambda\xi-T\xi\|\to0$ for each $\xi\in H$. This result prepares techniques for operator topologies, but is not a consequence of the convergence in norm of Problem 2.8.

\section{Problem 2.10: spectrum of the unilateral shift}


Let $S:\ell^2(\mathbb N)\to\ell^2(\mathbb N)$ be given by
\[
S(x_1,x_2,\dots)=(0,x_1,x_2,\dots).
\]

\noindent\textbf{Statement.} Prove that the spectrum of the unilateral shift is $\spec(S)=\overline{\mathbb D}$.

\begin{proof}[Solution]
Since $S$ is an isometry, $\|S\|=1$ and by the spectral bound
$\spec(S)\subset\overline{\mathbb D}$. Now let $|\lambda|<1$. The vector
\[
 v_\lambda=(1,\overline\lambda,\overline\lambda^2,\ldots)
\]
belongs to $\ell^2$ and satisfies
\[
 S^*v_\lambda=\overline\lambda\,v_\lambda.
\]
Thus $\overline\lambda I-S^*$ is not invertible. If $\lambda I-S$ were invertible, then its
adjoint $\overline\lambda I-S^*$ would also be invertible, contradiction. Therefore
$\lambda\in\spec(S)$ for all $|\lambda|<1$. Since the spectrum is closed,
$\overline{\mathbb D}\subset\spec(S)$. Along with the opposite inclusion,
\[
 \spec(S)=\overline{\mathbb D}.
\]
\end{proof}

\proofdiagrambranch
{$\|S\|=1\Rightarrow\spec(S)\subset\overline{\mathbb D}$}
{For $|\lambda|<1$, $S^*$ has eigenvector of value $\overline\lambda$}
{If $\lambda I-S$ were invertible, its adjoint would be too}
{$\mathbb D\subset\spec(S)$; the spectrum is closed}
{$\spec(S)=\overline{\mathbb D}$}

\begin{cthejemplo}{(A spectral disk with no eigenvalues)}
The unilateral shift of Problem 2.10 illustrates the difference between spectrum and eigenvalues.
If $S\xi=\lambda\xi$, the first coordinate gives $0=\lambda\xi_1$; for
$\lambda\ne0$, the subsequent coordinates successively force $\xi_n=0$.
For $\lambda=0$, injectivity of $S$ likewise gives $\xi=0$.
Thus $S$ has no eigenvalues, although Problem 2.10 proves
$\spec(S)=\overline{\mathbb D}$.
\end{cthejemplo}

\noindent\textbf{Dependencies.} The spectral bound of Part 1 gives $\spec(S)\subset\overline{\mathbb D}$. Problem 1.11(d) already shows $0\in\spec(S)$; the eigenvectors of $S^*$ give all the interior points and the closure of the spectrum adds the boundary.

\noindent\textbf{How this problem is used later.} This contrasts the spectrum of a non-normal operator with the characterization of unitary operators in Problem 2.6(b). It does not allow attributing the same spectral disk to the composition shifts of the Haseman problem.

\chapter{Research Examples: How the Theory Reappears in Haseman Problems and Singular Integral Operators}

\section{From the boundary condition to an operator}

In Karlovich \cite[pp.~463--464, eqs.~(1.1)--(1.2)]{Karlovich2017}, $\Gamma$ is an
open star of logarithmic spirals and $1<p<\infty$. Take $G\in\SO(\Gamma)$,
$g\in L^p(\Gamma)$, and a diffeomorphism $\alpha:\Gamma\to\Gamma$ that preserves
the orientation of each branch and satisfies the source's assumptions. The Haseman
condition is
\[
 \Phi^+(\alpha(t))=G(t)\Phi^-(t)+g(t),\qquad t\in\Gamma.
\]
Here $S_\Gamma\in\mathcal B(L^p(\Gamma))$ is the Cauchy singular integral operator,
$V_\alpha f=f\circ\alpha$, and $M_Gf=Gf$. If $\varphi$ is the Cauchy density, the
Sokhotski--Plemelj formulas give
\[
 \Phi^+=P_+\varphi,\qquad \Phi^-=-P_-\varphi,\qquad
 P_\pm=\tfrac12(I\pm S_\Gamma).
\]
The sign of the negative trace explains the plus sign in the operator equation:
\[
 \bigl(V_\alpha P_++M_GP_-\bigr)\varphi=g,\qquad
 T:=V_\alpha P_++M_GP_-\in\mathcal B(L^p(\Gamma)).
\]
The course supplies the language for asking whether $T$ is invertible and for
studying its spectrum. The boundary reduction and Fredholm criterion are later
results that depend on the article's geometric and symbol assumptions.

\begin{cthejemplo}{(The identity as a case of the Haseman equation)}
On $X=L^p(\Gamma)$ consider the formula $T=V_\alpha P_++M_GP_-$ obtained above.
We check which operator results when $\alpha=\mathrm{id}_\Gamma$ and $G=1$.
Then $V_\alpha=I_X$ and $M_G=I_X$, while the definition of $P_\pm$ gives
\[
 P_++P_- =\tfrac12(I_X+S_\Gamma)+\tfrac12(I_X-S_\Gamma)=I_X.
\]
Thus $T=P_++P_-=I_X$. The Plemelj reduction here does not require
$S_\Gamma^2=I_X$.
\end{cthejemplo}

\begin{center}
\begin{tikzpicture}[node distance=5mm and 8mm]
\node[theorybox,text width=3.35cm] (b) {Haseman boundary\\$\Phi^+\circ\alpha=G\Phi^-+g$};
\node[theorybox,text width=3.35cm,right=of b] (p) {Sokhotski--Plemelj\\$\Phi^+=P_+\varphi$, $\Phi^-=-P_-\varphi$};
\node[researchbox,text width=3.35cm,right=of p] (t) {$T=V_\alpha P_++M_GP_-$\\$T\varphi=g$};
\draw[arr] (b)--(p);
\draw[arr] (p)--(t);
\end{tikzpicture}
\end{center}

\section{From the compact ideal to the Fredholm property}

Let $X$ be an infinite-dimensional complex Banach space. The closed two-sided ideal
$\mathcal K(X)\triangleleft\mathcal B(X)$ yields the Calkin algebra
$\mathcal Q(X)=\mathcal B(X)/\mathcal K(X)$. For $A\in\mathcal B(X)$,
\[
 A\text{ Fredholm}
 \quad\Longleftrightarrow\quad
 A+\mathcal K(X)\in\GL(\mathcal Q(X)).
\]
Karlovich states this equivalence in the introduction to his 2025 article
\cite[p.~2]{Karlovich2025}. Problem 1.4 explains why a quotient by a closed ideal
remains a Banach algebra; the Fredholm equivalence adds a theorem specific to operators.

\begin{cthejemplo}{(Invertible modulo compact operators without an exact inverse)}
Atkinson's criterion distinguishes an exact inverse from an inverse in the Calkin algebra.
To see this on $X=\ell^2$, let $P$ be the rank-one orthogonal projection onto
$\mathbb Ce_1$ and set $A=I_X-P$. Then
$\ker A=\operatorname{Ran}P\ne\{0\}$, so $A$ is not invertible. However,
$P\in\mathcal K(X)$ and $A+\mathcal K(X)=I_X+\mathcal K(X)$: its class is the
identity of $\mathcal Q(X)$. Moreover,
$\operatorname{Ran}(A)=(\mathbb Ce_1)^\perp$ is closed, while the kernel and
cokernel each have dimension one. Thus $A$ is Fredholm.
\end{cthejemplo}

If the \emph{Fredholm essential spectrum} is defined by
\[
 \sigma_{\mathrm{ess}}(A)
 :=\sigma_{\mathcal Q(X)}\bigl(A+\mathcal K(X)\bigr),
\]
then $0\notin\sigma_{\mathrm{ess}}(A)$ is equivalent to $A$ being Fredholm.
This does not say that $0\notin\sigma_{\mathcal B(X)}(A)$; the example $I_X-P$
shows the difference.

\begin{center}
\begin{tikzpicture}[node distance=5mm and 7mm]
\node[theorybox,text width=2.65cm] (ideal) {$\mathcal K(X)\triangleleft\mathcal B(X)$\\closed ideal};
\node[theorybox,text width=2.65cm,right=of ideal] (quot) {$\mathcal Q(X)$\\quotient};
\node[theorybox,text width=2.65cm,right=of quot] (class) {$A+\mathcal K(X)$\\invertible class};
\node[researchbox,text width=2.65cm,right=of class] (fred) {$A$ Fredholm\\Calkin criterion};
\draw[arr] (ideal)--(quot);
\draw[arr] (quot)--(class);
\draw[arr] (class)--(fred);
\end{tikzpicture}
\end{center}

On a Hilbert space $H$, $\mathcal B(H)$ and its quotient by $\mathcal K(H)$ are
$C^*$-algebras. For a general Banach space, the preceding argument needs only
the Banach algebra structure. Symbol calculations and localization seek to turn
invertibility of the class into verifiable conditions for a \emph{specific operator
algebra}; simply writing down the quotient does not produce a criterion.

\Needspace{17\baselineskip}
\section{Cauchy projections: the model case and the Haseman case}

The comparison identifies which assumption turns the formula $(I\pm S_\Gamma)/2$
into a projection and which part of the formula is used in the Haseman equation.

In the model $\Gamma=\mathbb R$ or $\mathbb T$ on $L^2(\Gamma)$, let $S_\Gamma$
be the Cauchy singular integral operator. Bastos--Fernandes--Karlovich recall
$S_\Gamma^2=I$ and that
\[
 P_\Gamma^\pm=\tfrac12(I\pm S_\Gamma)
\]
are orthogonal projections \cite[§2, p.~22]{BFK2006}. The identity
$S_\Gamma^2=I$ gives $P_\Gamma^\pm{}^2=P_\Gamma^\pm$ and
$P_\Gamma^+P_\Gamma^-=0$; self-adjointness of $S_\Gamma$ also gives
$(P_\Gamma^\pm)^*=P_\Gamma^\pm$. This is the literal connection with
projections $p=p^*=p^2$ in a $C^*$-algebra.

\begin{center}
\begin{tikzpicture}[node distance=5mm and 7mm]
\node[theorybox,text width=2.6cm] (s) {$S_\Gamma^2=I$\\$S_\Gamma^*=S_\Gamma$};
\node[theorybox,text width=2.6cm,right=of s] (def) {$P_\Gamma^\pm=(I\pm S_\Gamma)/2$};
\node[theorybox,text width=2.6cm,right=of def] (proj) {$P_\Gamma^\pm{}^2=P_\Gamma^\pm$\\$(P_\Gamma^\pm)^*=P_\Gamma^\pm$};
\node[researchbox,text width=2.6cm,right=of proj] (split) {$L^2(\Gamma)=\operatorname{Ran}P_\Gamma^+\oplus\operatorname{Ran}P_\Gamma^-$};
\draw[arr] (s)--(def);
\draw[arr] (def)--(proj);
\draw[arr] (proj)--(split);
\end{tikzpicture}
\end{center}

On the open star of Karlovich 2017, the same \emph{formula}
$P_\pm=(I\pm S_\Gamma)/2$ appears in Sokhotski--Plemelj and in
$T=V_\alpha P_++M_GP_-$, but idempotence from the preceding model cannot
be imported. If $X$ is a Banach space, $S\in\mathcal B(X)$, and
$P_\pm=(I_X\pm S)/2$, then
\[
 P_\pm^2-P_\pm=\tfrac14(S^2-I_X),\qquad
 P_+P_-=\tfrac14(I_X-S^2),\qquad P_++P_-=I_X.
\]
Indeed, the Mellin symbol of $S_\Gamma$ in the one-branch case with $p=2$
and geometric parameter $\delta=0$ is $\varepsilon_1\tanh(\pi x)$, where
$\varepsilon_1\in\{-1,1\}$ records orientation; its square is not
identically one \cite[Theorem~5.3, p.~476]{Karlovich2017}.
Thus, in this case $P_\pm$ are not algebraic projections.
The correct connection to Haseman is the \emph{decomposition of traces}
by Plemelj, not an orthogonal decomposition of the space.

\begin{cthnota}{(Where the assumptions enter)}
The identity $S_\Gamma^2=I$ is exactly what gives idempotence;
$S_\Gamma^*=S_\Gamma$ is needed for orthogonality. The model
$L^2(\mathbb R)$ or $L^2(\mathbb T)$ satisfies both properties in the cited
source. The 2017 Haseman problem acts on $L^p(\Gamma)$ over an open star and
uses only the Plemelj formula here.
\end{cthnota}

\Needspace{25\baselineskip}
\section{Unitary shifts and composition shifts}

We compare a unitary rotation with the composition operator that appears in
Haseman. The underlying space and measure determine which property can be asserted.

In \cite[§1, p.~20]{BFK2006}, the circle $\mathbb T$ acts by a unitary representation;
for $z,t\in\mathbb T$ and $\varphi\in L^2(\mathbb T)$ one writes
\[
 U:\mathbb T\to\mathcal B(L^2(\mathbb T)),
 \qquad
 (U_z\varphi)(t)=\varphi(zt).
\]
Rotation preserves Haar measure $m$: substituting $u=zt$ gives
\[
 \|U_z\varphi\|_2^2=\int_{\mathbb T}|\varphi(zt)|^2\,dm(t)
 =\int_{\mathbb T}|\varphi(u)|^2\,dm(u)=\|\varphi\|_2^2.
\]
Since $U_z^{-1}=U_{z^{-1}}$, $U_z$ is unitary. The Haseman problem instead
involves $V_\alpha f=f\circ\alpha$ on $L^p(\Gamma)$; Karlovich imposes
regularity and derivative-control assumptions on the shift on each branch
\cite[§2, p.~467]{Karlovich2017}. These conditions allow one to work with a
bounded operator, but do not assert that it is unitary.

\begin{center}
\begin{tikzpicture}[node distance=10mm]
\node[theorybox] (shift) {Common idea\\composition operator / shift};
\node[researchbox,above right=of shift] (unit) {$U_z$ on $L^2(\mathbb T)$\\unitary};
\node[researchbox,below right=of shift] (hasa) {$V_\alpha$ on $L^p(\Gamma)$\\bounded under shift assumptions};
\node[theorybox,right=22mm of unit] (uspec) {Unitary $\Rightarrow\spec(U_z)\subset\mathbb T$};
\node[theorybox,right=22mm of hasa] (fred) {Combined with $P_\pm$ and $M_G$ in the Haseman operator};
\draw[arr] (shift)--(unit);
\draw[arr] (shift)--(hasa);
\draw[arr] (unit)--(uspec);
\draw[arr] (hasa)--(fred);
\end{tikzpicture}
\end{center}

\section{Slowly oscillating functions and fibers of characters}

For $a\in C_b(\mathbb R_+)$, write
$\operatorname{osc}(a,[r,2r])=\sup_{s,t\in[r,2r]}|a(s)-a(t)|$.
Karlovich defines $\SO(\mathbb R_+)$ by requiring this oscillation to
tend to zero both as $r\to0$ and as $r\to\infty$
\cite[§2, pp.~465--466]{Karlovich2017}. For $I=[r,2r]$ and functions
$a,b\in C_b(\mathbb R_+)$,
\[
\begin{aligned}
 \operatorname{osc}(a+b,I)&\le\operatorname{osc}(a,I)+\operatorname{osc}(b,I),\\
 \operatorname{osc}(ab,I)&\le\|a\|_\infty\operatorname{osc}(b,I)
                         +\|b\|_\infty\operatorname{osc}(a,I),\\
 \operatorname{osc}(\overline a,I)&=\operatorname{osc}(a,I).
\end{aligned}
\]
If $a_n\to a$ uniformly, then
$\operatorname{osc}(a,I)\le2\|a-a_n\|_\infty+
\operatorname{osc}(a_n,I)$; first fix a large $n$ and then let $r$
approach the endpoint. Thus the class is a closed subalgebra, contains
the constants, and is stable under pointwise involution. The identity
$\|\overline a a\|_\infty=\|a\|_\infty^2$ shows that it is a unital
commutative $C^*$-algebra.

\begin{cthejemplo}{(Slow oscillation without an ordinary limit)}
We check that slow oscillation at the endpoints need not imply an
ordinary limit. On $\mathbb R_+$ consider the function
\[
 a(r)=\sin\!\bigl(\log(1+|\log r|)\bigr),\qquad r>0.
\]
It belongs to $\SO(\mathbb R_+)$. If $s,t\in[r,2r]$, then
$|\log s-\log t|\le\log2$. Put $h(u)=\log(1+|u|)$.
On the interval $I_r=[\log r,\log r+\log2]$, when it is away from $0$,
the mean value theorem gives
\[
 |h(\log s)-h(\log t)|
 \le \frac{\log2}{1+\min_{u\in I_r}|u|}.
\]
The denominator tends to infinity both as $r\to0$ and as $r\to\infty$.
Since $|\sin v-\sin w|\le|v-w|$, the oscillation of $a$ on $[r,2r]$
tends to zero at both endpoints. To see that there is no ordinary
limit, define
\[
 r_n^+=\exp\!\bigl(e^{\pi/2+2\pi n}-1\bigr),\qquad
 r_n^-=\exp\!\bigl(e^{3\pi/2+2\pi n}-1\bigr).
\]
Both sequences tend to infinity, while $a(r_n^+)=1$ and $a(r_n^-)=-1$.
Their reciprocals tend to zero and give the same two values, because
$a(r)=a(1/r)$. This function shows why boundary characters of
$\SO(\mathbb R_+)$ cannot simply be
replaced by evaluations $a(0)$ and $a(\infty)$.
\end{cthejemplo}

The maximal ideal space $\Omega(\SO(\mathbb R_+))$ contains evaluations
at $r>0$ and fibers $M_0,M_\infty$ over the endpoints
\cite[§2, p.~466]{Karlovich2017}. The Gelfand transform
$a\mapsto\widehat a$, $\widehat a(\xi)=\xi(a)$, describes values of these
characters without pretending that $a(0)$ and $a(\infty)$ exist. In
symbol calculations, these fibers are part of the parameter space;
the Gelfand map of the coefficients and the full symbol of an operator
are different maps.

\section{Generators, quotients, and symbol maps}

Karlovich \cite[eq.~(1.4), p.~3]{Karlovich2025} fixes a star-shaped
curve with cusps $\Gamma$ from (1.1)--(1.2), an admissible power weight
$w$ from §2.1, with $1<p<\infty$ and exponent
$-1/p<\alpha_{\mathrm{wt}}<1-1/p$, and considers the Banach algebra
generated by
\[
 \mathfrak B=
 \operatorname{alg}_{\mathbb R}
 \{M_a,S_\Gamma,C:\ a\in\SO^\diamond(\Gamma)\}
 \subset\mathcal B(\mathbb RL^p(\Gamma,w)).
\]
Here $M_a$ multiplies by $a$, $S_\Gamma$ is the Cauchy singular operator,
and $Cf=\overline f$ is complex conjugation. The notation
$\mathbb RL^p(\Gamma,w)$ means that the $L^p$ space is considered over
the real scalars: for $f\in L^p(\Gamma,w)$,
$C(if)=\overline{if}=-i\overline f=-iCf$, whereas
$C(rf)=rCf$ for $r\in\mathbb R$. Thus $C$ is linear over $\mathbb R$
but not over $\mathbb C$. This algebra therefore does not automatically
become a complex $C^*$-algebra, and the spectral results for complex
Banach algebras in Part 1 cannot simply be applied. The notation $\SO^\diamond(\Gamma)$
belongs to the 2025 article and does not denote the same algebra
$\SO(\mathbb R_+)$ as in the 2017 article. The 2025 article calls the
weight exponent $\alpha$; the subscript here avoids confusing it with
the shift $\alpha$ of Karlovich 2017.

\begin{center}
\begin{tikzpicture}[node distance=5mm and 8mm]
\node[researchbox,text width=3.25cm] (gen) {$M_a,S_\Gamma$ complex-linear\\$C$ real-linear};
\node[theorybox,text width=3.25cm,right=of gen] (alg) {$\mathfrak B$\\real algebra;\\norm closed};
\node[researchbox,text width=3.25cm,right=of alg] (sym) {Quotient and symbols\\criterion in a fixed algebra};
\draw[arr] (gen)--(alg);
\draw[arr] (alg)--(sym);
\end{tikzpicture}
\end{center}

Problem 1.3 teaches what it means to generate and close an algebra.
Problem 1.4 teaches how to pass to a quotient. In the 2025 article,
the additional step is to construct matrix-valued symbol maps and
prove that they detect the Fredholm property. Theorem~6.2 treats the
complex subalgebra
$\mathfrak A=\operatorname{alg}_{\mathbb C}
\{M_a,S_\Gamma,CS_\Gamma C:a\in\SO^\diamond(\Gamma)\}$;
the complex matrix representation of Theorem~7.1 then allows the
criterion for $N\in\mathfrak B$ to be formulated using the symbol map
$Y$ of Theorem~7.2 \cite[pp.~25--26]{Karlovich2025}. The latter
requires membership in $\mathfrak B$ and the article's assumptions
on the curve, weight, and coefficients; it does not by itself certify
a Haseman operator with a shift.

\begin{cthnota}{(Four related but distinct maps)}
An algebra homomorphism preserves sums and products. A
$*$-homomorphism also preserves the involution. The Gelfand transform
$a\mapsto\widehat a$ of a commutative algebra evaluates at its
characters. A Fredholm symbol map is constructed for a particular
operator algebra, often after passing modulo compact operators, and
may take matrix values. These constructions share algebraic
operations, but their domains, codomains, and assumptions cannot be
interchanged.
\end{cthnota}

\Needspace{20\baselineskip}
\section{The full route from the problem to a criterion}

In the diagram, blue marks the language of the course, green the
research object, and the arrow toward symbols presupposes that the
corresponding calculus has been constructed for the particular
algebra. Here $X=L^p(\Gamma)$, so $T+\mathcal K(X)$ is a class in the
Calkin algebra associated with the same underlying space. The
criterion of Karlovich 2017 also requires the assumptions of its
Theorem 7.1 \cite[pp.~487--489]{Karlovich2017}.

\begin{center}
\begin{tikzpicture}[node distance=5mm and 7mm]
\node[researchbox,text width=3.25cm] (h) {[research]\\Haseman boundary};
\node[researchbox,text width=3.25cm,right=of h] (p) {[research]\\Sokhotski--Plemelj};
\node[researchbox,text width=3.25cm,right=of p] (t) {[research]\\$T=V_\alpha P_++M_GP_-$};
\node[theorybox,text width=3.25cm,below=of t] (b) {[course]\\$T\in\mathcal B(L^p(\Gamma))$};
\node[theorybox,text width=3.25cm,left=of b] (q) {[course + bridge]\\$T+\mathcal K(X)$ in $\mathcal Q(X)$};
\node[researchbox,text width=3.25cm,left=of q] (f) {[bridge]\\invertibility of the class\\$\Longleftrightarrow$ Fredholm};
\node[researchbox,text width=3.25cm,below=of f] (s) {[research]\\symbols and localization};
\node[researchbox,text width=3.25cm,right=of s] (c) {[research]\\Fredholm criterion\\under article assumptions};
\draw[arr] (h)--(p);
\draw[arr] (p)--(t);
\draw[arr] (t)--(b);
\draw[arr] (b)--(q);
\draw[arr] (q)--(f);
\draw[arr] (f)--(s);
\draw[arr] (s)--(c);
\end{tikzpicture}
\end{center}

The last row describes how the Fredholm property is \emph{checked} in
the literature. It does not claim that every Fredholm operator
automatically has the particular symbol of the article, or that the
2017 criterion applies to a 2025 operator with cusps.

\Needspace{19\baselineskip}
\section{A reading path for the research manuscript}

When an operator construction appears in the manuscript, it is
helpful to ask the following questions in order:
\begin{enumerate}
\item In which Banach algebra does the operator live?
\item What is the identity, and what are the relevant generators?
\item Are we working exactly or modulo an ideal, usually $\K(X)$?
\item Is the subalgebra under consideration commutative, or does it only
      have a commutative central subalgebra?
\item Do characters, a maximal ideal space, or a Gelfand transform appear?
\item Is there a complex $C^*$-algebra structure (typically on $L^2$)
      that permits the use of normal elements, unitaries, projections,
      or functional calculus?
\item Which part of the Fredholm criterion is equivalent to
      invertibility in a quotient or of a symbol?
\end{enumerate}

These questions distinguish the common language from the assumptions
that authorize each theorem. For example, a published criterion for
an operator algebra on a curve with cusps does not automatically apply
to an operator with a shift on a different geometry.

\chapter{Dependency Map: Results, Assignments, and Research}

\section{What an arrow means in this map}

\begin{context}
A solid arrow means that the source result is used in proving or constructing the target
result. A dashed arrow marks an application or conceptual comparison: it does not claim
that the target result follows from the first. The assumptions written along each path are
part of the connection.
\end{context}

\section{Paths used in proofs}

\begin{dependencias}
\textbf{Continuity and inverses.}
\[
\boxed{\text{Neumann series}}
\longrightarrow
\boxed{\text{1.1(b): continuous inversion}}
\longrightarrow
\boxed{\text{Prop. 2.17: spectral permanence}}.
\]
The last arrow uses $(b+i/n)^{-1}\to b^{-1}$ in the ambient algebra and the fact that the
$C^*$-subalgebra is closed. Problem 1.1(a) supplies continuity of multiplication through
another route for Problems 1.3, 1.7, and 2.8.
\end{dependencias}

\begin{dependencias}
\textbf{Commutative reduction.}
\[
\begin{gathered}
\boxed{\text{Problem 1.2: commuting inverses}}
\quad+\quad
\boxed{\text{Problem 1.3: generated algebra}}\\[2mm]
\Downarrow\\[2mm]
\boxed{\text{Problem 1.9: inverse-closed subalgebra}}\\[2mm]
\Downarrow\\[2mm]
\boxed{\text{Prop. 2.13: real spectrum}}.
\end{gathered}
\]
The last proof also uses the $C^*$ identity, the fact that characters are real on
self-adjoint elements, and the description of the spectrum through characters.
Problem 1.9 supplies the commutative subalgebra that preserves the spectrum;
by itself it does not prove that the spectrum is real.
\end{dependencias}

\begin{dependencias}
\textbf{From functions to functional calculus.}
\[
\begin{gathered}
\boxed{\text{Problem 1.3(d): }C^*(a,1)\text{ commutative for normal $a$}}\\[2mm]
\quad+\quad
\boxed{\text{Prop. 2.17: spectral permanence}}
\quad+\quad
\boxed{\text{Thm. 2.20: Gelfand--Naimark}}\\[2mm]
\Downarrow\\[2mm]
\boxed{\text{Thm. 2.22: functional calculus}}.
\end{gathered}
\]
Problem 1.10 computes the spectrum of a continuous function and, together with functional
calculus and spectral permanence, gives Problem 2.4. Problems 2.2(a,c) give the forward
directions of Problem 2.6(c,b); the converses use normality and functional calculus.
\end{dependencias}

\begin{dependencias}
\textbf{Continuity as the operator varies.}
\[
\begin{gathered}
\boxed{\text{Problem 1.7: norm limit of normal elements is normal}}
\quad+\quad
\boxed{\text{Stone--Weierstrass}}\\[2mm]
\quad+\quad
\boxed{\text{isometric functional calculus}}\\[2mm]
\Downarrow\\[2mm]
\boxed{\text{Problem 2.8: }f(x_n)\to f(x)}.
\end{gathered}
\]
Convergence here is in norm. Problem 2.9 studies a different situation: a bounded,
increasing net of self-adjoint operators converges strongly, that is, vector by vector
in the norm of $H$.
\end{dependencias}

\section{Coverage of the exercises}

\small
\begin{longtable}{>{\bfseries}p{1.3cm}p{4.3cm}p{8.2cm}}
\toprule
Problem & Result proved & Exact use or scope of the connection\\
\midrule
\endhead
1.1 & Continuous multiplication and inversion & Multiplication supports Problems 1.3 and 1.7; inversion is used in Prop. 2.17, and multiplication again in Problem 2.8.\\
1.2 & Commutation with inverses and adjoints & Used when adjoining inverses in the construction for Problem 1.9.\\
1.3 & Generated algebra and commutation & For normal $a$, constructs commutative $C^*(a,1)$, used in Proposition 2.18 and Theorem 2.22.\\
1.4 & Quotient by a closed ideal & Makes $\mathcal B(X)/\K(X)$ a Banach algebra; the Fredholm--invertibility equivalence is the external Atkinson theorem.\\
1.5 & $C_0(X)$ closed in $C_b(X)$ & When $X$ is noncompact, it is nonunital; the general classification of $C_0(X)$ requires the nonunital version of the theorem in Section 2.\\
1.6 & $\spec(x^*)=\overline{\spec(x)}$ & Symmetry holds in involutive algebras; reality of a self-adjoint element's spectrum also needs the $C^*$ argument of Prop. 2.13.\\
1.7 & Normality is closed & The first assumption needed in Problem 2.8 to define $f(x)$ at the limit.\\
1.8 & Spectral inclusion under homomorphisms & Applied to $a^*a$, it enters the proof that $*$-homomorphisms are contractive (Thm. 2.8).\\
1.9 & Inverse-closed commutative subalgebra & Explicitly used in the proof of Prop. 2.13 to preserve a self-adjoint element's spectrum.\\
1.10 & $\spec(f)=f(X)$ for compact $X$ & Computes the spectrum of $f$ in the proof of Problem 2.4; illustrates the $C(X)$ model of Gelfand--Naimark.\\
1.11 & $\mathcal B(H)$ and unilateral shift & Part (a) gives a concrete model for Problem 2.1; (c,d) distinguish approximate point spectrum from spectrum, and Problem 2.10 completes the calculation of $\spec(S)$.\\
1.12 & Disc algebra & Shows that a unital, closed inclusion of Banach algebras can change the spectrum; the inclusion is not one of $C^*$-subalgebras.\\
\midrule
2.1 & $C_b(X)\hookrightarrow\mathcal B(\ell^2(X))$ & Directly verifies an isometric $*$-homomorphism; a concrete example of Thm. 2.27.\\
2.2 & Idempotents, products, unitaries & Parts (a,c) give forward directions of Prop. 2.28; (b) estimates radii of commuting products and is not needed for that proposition.\\
2.3 & Bounded product and $c_0$ direct sum & Constructs a $C^*$-algebra and its closed ideal; Problem 1.4 allows the quotient to be viewed as a Banach algebra.\\
2.4 & Continuous spectral mapping & Depends on Thm. 2.22, Prop. 2.17, and Problem 1.10; applies to functions of normal elements.\\
2.5 & $X\cong\Omega(C(X))$ & Gives uniqueness of the compact space in Cor. 2.21.\\
2.6 & Self-adjoints, unitaries, and projections & Reconstructs Prop. 2.28 through functional calculus and uses counterexamples to show why normality is needed.\\
2.7 & Unitary factor of an invertible element & Uses the positive root from Prop. 2.30 and invertibility of $a^*a$ to obtain $a=u|a|$.\\
2.8 & Continuity of $f(x_n)$ & Combines Problem 1.7, Stone--Weierstrass, and functional calculus; convergence is in norm.\\
2.9 & Monotone convergence of self-adjoint operators & Uses positive order, polarization, and Riesz; concludes strong convergence without claiming norm convergence.\\
2.10 & Spectrum of the unilateral shift & Extends Problem 1.11(d) to $\spec(S)=\overline{\mathbb D}$; $S$ is not normal.\\
\bottomrule
\end{longtable}
\normalsize

\section{Two carefully qualified bridges to Haseman}

\begin{center}
\begin{tikzpicture}[node distance=8mm]
\node[taskbox,text width=2.6cm] (q) {Problem 1.4\\quotient by an ideal};
\node[theorybox,text width=2.7cm,right=of q] (calkin) {$\K(X)$ closed\\$\mathcal B(X)/\K(X)$};
\node[researchbox,text width=2.6cm,right=of calkin] (class) {Class of $A$\\$A+\K(X)$};
\node[flowend,text width=2.7cm,right=of class] (fred) {Atkinson\\invertible $\Leftrightarrow$ Fredholm};
\draw[arr] (q)--(calkin);
\draw[arr] (calkin)--(class);
\draw[arr] (class)--(fred);
\end{tikzpicture}
\end{center}

\noindent For infinite-dimensional Banach $X$, the first arrow uses that $\K(X)$ is a closed
ideal. The last is Atkinson's theorem, not a consequence of Problem 1.4. For operators
in the Haseman research line, this formulation poses the Fredholm problem in a quotient;
see \cite[p.~2]{Karlovich2025}. As a concrete check, on $X=\ell^2$ let $P$ project onto
the first standard basis vector $e_1$, and set $A=I_X-P$. Then $A e_1=0$, but
$A+\K(X)=I_X+\K(X)$ because $P$ is compact. The calculation shows what information
the quotient retains and why Problem 1.4 prepares the formulation of Atkinson's theorem.

\begin{center}
\begin{tikzpicture}[node distance=9mm and 8mm]
\node[theorybox,text width=3.0cm] (model) {Model $L^2(\Gamma_0)$,\\$\Gamma_0=\mathbb R$ or $\mathbb T$\\$S_{\Gamma_0}^2=I$, $S_{\Gamma_0}^*=S_{\Gamma_0}$};
\node[taskbox,text width=3.0cm,right=of model] (orth) {$P_{\Gamma_0}^\pm=(I\pm S_{\Gamma_0})/2$\\orthogonal projections};
\draw[arr] (model)--(orth);
\node[researchbox,text width=3.0cm,below=of model] (plem) {Karlovich's open star\\Sokhotski--Plemelj};
\node[researchbox,text width=3.0cm,right=of plem] (trace) {$P_\pm=(I\pm S_\Gamma)/2$\\trace components};
\node[flowend,text width=3.0cm,right=of trace] (has) {$T=V_\alpha P_+ + M_GP_-$\\Haseman equation};
\draw[arr] (plem)--(trace);
\draw[arr] (trace)--(has);
\draw[mergearr,dash pattern=on 4pt off 2pt] (orth.south)--(trace.north);
\end{tikzpicture}
\end{center}

\noindent The dashed arrow compares only the form of the formulas. In the model of
Bastos--Fernandes--Karlovich \cite[§2, p.~22]{BFK2006}, $P_{\Gamma_0}^\pm$ are
orthogonal projections. On the open star of Karlovich 2017, boundary reduction to $T$
follows from Sokhotski--Plemelj \cite[pp.~463--464]{Karlovich2017}; according to
Theorem 5.3 of that source, $S_\Gamma^2=I$ does not hold there in general. Thus the
operators $P_\pm$ in that equation are not called idempotents or projections here.
For the check $\alpha=\mathrm{id}$ and $G=1$, the Haseman formula gives
$T=P_++P_-=I$ by cancellation of the two $S_\Gamma$ terms. This calculation uses
the Plemelj formula without requiring $S_\Gamma^2=I$.

\begin{reflection}
To check an arrow, locate the line of the proof that uses its source. For example, identify
where Problem 1.9 enters the reality of the spectrum and where Problem 1.10 enters Problem
2.4. Then verify in the $L^2(\Gamma_0)$ model how $S_{\Gamma_0}^2=I$ gives
idempotents and how $S_{\Gamma_0}^*=S_{\Gamma_0}$ makes them orthogonal projections.
Compare those assumptions with the assumptions for the open star.
\end{reflection}

\chapter*{Index of Connections to the Research Literature}
\addcontentsline{toc}{chapter}{Index of Connections to the Research Literature}

This table helps the reader return to a source after studying the course concept.
The word \emph{related} does not claim that the advanced object generalizes the
elementary result. In the operator rows, $X$ denotes a complex Banach space and
$H$ a Hilbert space.

\footnotesize
\begin{longtable}{>{\raggedright\arraybackslash}p{2.5cm}>{\raggedright\arraybackslash}p{5.5cm}>{\raggedright\arraybackslash}p{3.6cm}>{\raggedright\arraybackslash}p{2.0cm}}
\toprule
Course concept & Appearance and what to recognize & Precise source & Type\\
\midrule
\endfirsthead
\toprule
Course concept & Appearance and what to recognize & Precise source & Type\\
\midrule
\endhead
Banach algebra &
$T=V_\alpha P_++M_GP_-\in\mathcal B(L^p(\Gamma))$:
sum and composition of bounded operators. &
Karlovich 2017, pp.~463--464, (1.1)--(1.2) \cite{Karlovich2017}. &
Direct example\\
Commutative $C^*$-algebra &
$C_b(\Gamma)$ and $\SO(\Gamma)$ are concrete coefficient algebras
with pointwise conjugation. &
Karlovich 2017, p.~464 \cite{Karlovich2017}. &
Direct example\\
Ideal and quotient &
$J=\mathcal K(X)$ in $\A=\mathcal B(X)$ gives the Calkin algebra;
invertibility of the class characterizes the Fredholm property. &
Karlovich 2025, p.~2 \cite{Karlovich2025}. &
Use of a property\\
Spectrum &
$\sigma_{\mathrm{ess}}(A)$ is defined here using the Calkin class;
it is distinguished from $\sigma_{\mathcal B(X)}(A)$. &
Developed in these notes from the quotient cited in Karlovich 2025, p.~2 \cite{Karlovich2025}. &
Related construction\\
Characters and Gelfand &
The fibers $M_0,M_\infty$ of $\Omega(\SO(\mathbb R_+))$ encode
limiting behavior of coefficients. &
Karlovich 2017, §2, pp.~465--466 \cite{Karlovich2017}. &
Use of a property\\
Projection in a $C^*$-algebra &
On $L^2(\Gamma)$, for $\Gamma=\mathbb R$ or $\mathbb T$,
$P_\Gamma^\pm$ are orthogonal.
The same formula on the open star does not guarantee idempotence. &
BFK 2006, §2, p.~22 \cite{BFK2006}; Karlovich 2017, Theorem~5.3, p.~476 \cite{Karlovich2017}. &
Example with assumptions\\
Unitaries and shifts &
Rotations $U_z$ are unitary; the Haseman shift $V_\alpha$
requires different boundedness assumptions. &
BFK 2006, §1, p.~20 \cite{BFK2006}; Karlovich 2017, §2, p.~467 \cite{Karlovich2017}. &
Comparison\\
Generated algebra &
Multipliers, $S_\Gamma$, and conjugation generate a real operator
algebra; typing conjugation prevents treating it as complex linear. &
Karlovich 2025, (1.4), p.~3 \cite{Karlovich2025}. &
Direct example\\
Homomorphisms &
Matrix-valued symbol maps help test the Fredholm property in a concrete
algebra; they are not automatically $*$-homomorphisms. &
Karlovich 2025, Theorems~6.2 and 7.2, pp.~25--26 \cite{Karlovich2025}. &
Related construction\\
Fredholm property &
The criterion for $T$ follows its Mellin representation modulo
compact operators and the associated symbol assumptions. &
Karlovich 2017, Theorem~7.1, pp.~487--489 \cite{Karlovich2017}. &
Further reading\\
\bottomrule
\end{longtable}
\normalsize


\renewcommand{\bibname}{Working References}
\begin{thebibliography}{99}
\addcontentsline{toc}{chapter}{Working References}
\bibitem{CastillejosParte1} J. Castillejos, \emph{$C^*$-álgebras: Parte 1. Teoría espectral en álgebras de Banach} [Part 1: Spectral Theory in Banach Algebras], course notes provided by the user.
\bibitem{CastillejosCap2} J. Castillejos, \emph{$C^*$-álgebras: Sección 2} [Section 2], course notes provided by the user.
\bibitem{Karlovich2017} Yu. I. Karlovich, \emph{The Haseman Boundary Value Problem with Slowly Oscillating Coefficients and Shifts}, Operator Theory: Advances and Applications 259 (2017), 463--500.
\bibitem{Karlovich2025} Yu. I. Karlovich, \emph{Algebras of singular integral operators with complex conjugation on curves with cusps}, Bol. Soc. Mat. Mex. 31 (2025), 102.
\bibitem{BFK2006} M. A. Bastos, C. A. Fernandes, Yu. I. Karlovich, \emph{$C^*$-Algebras of Integral Operators with Piecewise Slowly Oscillating Coefficients and Shifts Acting Freely}, Integral Equations Operator Theory 55 (2006), 19--67.
\end{thebibliography}

\end{document}
